\section{Arithmetic quadrics and geometric detection}\label{d3:sec:geometry}
The supported classes on an arithmetic quadric are detected by a ruling
line bundle. The geometric argument has two steps: the ruling survives
every finite localization permitted by the arithmetic local ring, and
its regulator pairs nontrivially with a number-field coefficient.

\subsection{The arithmetic quadric and its coefficient class}

For an abelian group $M$, write $M\Qrat=M\otimes_{\Z}\Q$. All $K$-groups in this article are Quillen algebraic $K$-groups. A coherent sheaf on a regular noetherian scheme is viewed in $K$-theory through a finite locally free resolution, or equivalently through the comparison $K\simeq G$.

The injectivity part of the rational Gersten conjecture asserts that, for every regular local ring $R$ with fraction field $F$ and every $n\geq 0$,
\[
K_n(R)\Qrat\longrightarrow K_n(F)\Qrat
\]
is injective. Its failure is sufficient to disprove the augmented rational Gersten conjecture.

\begin{theorem}\label{d3:thm:main}
Let $p$ be an odd prime and set
\[
A=\left(\Z_{(p)}[x,y,z]/(p+xy+z^2)\right)_{(x,y,z)}.
\]
There exists an integral class $c\in K_3(A)$ of infinite order whose image in $K_3(\Frac A)$ is zero. In particular,
\[
\ker\bigl(K_3(A)\Qrat\longrightarrow K_3(\Frac A)\Qrat\bigr)\neq 0.
\]
\end{theorem}

Only $p=3$ is needed for a counterexample. The proof is given for odd $p$ because its algebraic geometry is identical for all such primes.

In the case $p=3$, Section~\ref{d3:sec:explicit} gives the following fully specified class:
\begin{equation}\label{d3:eq:explicitfront}
\boxed{\quad
c_{\mathrm{exp}}=2\left\{\frac{1+z}{2},\,1+x,\,
1+\frac{(1+x)y}{4}\right\}\in K_3(A).
\quad}
\end{equation}
Braces mean the product of the three indicated unit classes in Quillen $K_1$. All entries are units in $A$. The factor two makes generic vanishing an integral assertion.


Put
\[
V=\Z_{(p)},\qquad W=V[u]/(u^2+p),\qquad E=\Q(u),\qquad \sigma(u)=-u.
\]
Thus $W$ is a discrete valuation ring (DVR) with uniformizer $u$ and residue field $\F_p$, and $E$ is an imaginary quadratic number field. Write
\[
A_0=V[x,y,z]/(p+xy+z^2),\quad \m=(x,y,z),\quad S=A_0\setminus\m.
\]
Then $A=S^{-1}A_0$. Let
\begin{align*}
B_0&=A_0\otimes_VW=W[x,y,z]/(xy+z^2-u^2),\\
B&=A\otimes_VW=S^{-1}B_0,\\
Q&=\Spec E[x,y,z]/(xy+z^2-u^2),\\
T&=\Spec B[1/u].
\end{align*}
For $s\in S$, let $Q_s$ denote the corresponding principal open of $Q$.

We use the following standard inputs.

\begin{enumerate}
\item Quillen localization and d\'evissage, the projection formula for a regular closed immersion, and flat base change for the associated pushforward. These may be formulated in $G$-theory and transported to $K$-theory on the regular schemes involved; see \cite[Sections~5--7]{Quillen}.
\item $K_2(\F_p)=0$ \cite{QuillenFinite}\cite[IV, Corollary~1.13]{Weibel}. Consequently $K_3(W)\to K_3(E)$ is surjective, integrally.
\item Borel's theorem gives a class in $K_3(E)$ with nonzero regulator at a complex embedding. The comparison of the Borel and Beilinson regulators identifies this, up to a nonzero scalar, with a class in
\[
\HD^1(\Spec\C,\R(2))=\C/\R(2),\qquad \R(j)=(2\pi i)^j\R.
\]
For an imaginary quadratic field, conjugating the embedding negates this regulator. See \cite[Proposition~12.2]{Borel} and \cite[Section~6]{BunkeTamme}.
\item The analytic version of the Beilinson regulator is natural and multiplicative and has the usual first Chern-class normalization. A precise reference is \cite[Theorem~2.31]{BunkeTamme}. We use its weight-three component
\[
r_3:K_3(U)\longrightarrow\HD^3(U,\R(3))
\]
on smooth complex fibers, and its weight-two component on number fields.
\item Algebraic $K$-theory of rings commutes with filtered colimits \cite[IV, Section~6.4]{Weibel}. In particular,
\begin{equation}\label{d3:eq:continuity}
K_3(T)=\colim_{s\in S} K_3(Q_s).
\end{equation}
\end{enumerate}

\begin{remark}\label{d3:rem:spreading}
The site in \cite{BunkeTamme} consists of regular separated schemes of finite type over $\Z$. Although $E$ and $Q_s$ are not finite type over $\Z$, each finite collection of coefficients, $K$-classes, line bundles, and morphisms used below spreads to a smooth model over $\OO_E[1/N]$ for a suitable positive integer $N$. These models are regular and of finite type over $\Z$, and all further integer localizations have the same complex fiber. Regulator naturality and filtered-colimit compatibility therefore give precisely the maps used here. Projection from the conjugation-invariant collection of complex components to a chosen embedding preserves products. No regulator on the infinite scheme $T$ is required.
\end{remark}

\subsection{Regularity and the supported class}

\begin{lemma}\label{d3:lem:regularA}
The ring $A$ is a three-dimensional regular local domain with residue field $\F_p$. Its maximal ideal is generated by $x,y,z$, its fraction field is $\Q(x,z)$, and $p\in\m^2$.
\end{lemma}
\begin{proof}
The ambient ring $V[x,y,z]_{(p,x,y,z)}$ is regular local of dimension four. The element $p+xy+z^2$ belongs to its maximal ideal but not to its square: its initial linear term is $p$. Quotienting by this element gives a regular local ring of dimension three. The relation $p=-xy-z^2$ gives both the asserted generators of the maximal ideal and $p\in\m^2$. Solving for $y$ after inverting $x$ gives the fraction field $\Q(x,z)$. Alternatively, the defining polynomial is primitive and irreducible as a polynomial of degree one in $y$, so the quotient is a domain.
\end{proof}

The regular principal divisor $x=0$ has coordinate ring
\[
R=A/(x)=\left(V[y,z]/(z^2+p)\right)_{(y,z)}
\cong W[y]_{(u,y)},
\]
where the map $W\to R$ sends $u$ to $z$. Let
\[
i:\Spec R\hookrightarrow\Spec A
\]
be the closed immersion.

\begin{lemma}\label{d3:lem:beta}
There is an integral class $\beta\in K_3(W)$ such that
\[
\sigma\beta=-\beta\quad\text{integrally},\qquad r_2(\beta|_E)\neq 0
\]
at a fixed complex embedding $E\hookrightarrow\C$.
\end{lemma}
\begin{proof}
Choose a class of $K_3(E)$ with nonzero Borel--Beilinson regulator. Localization gives
\[
K_3(W)\longrightarrow K_3(E)\longrightarrow K_2(\F_p)=0,
\]
so choose an integral lift $b_0\in K_3(W)$. Set $\beta=b_0-\sigma b_0$. Since $\sigma^2=1$, we have $\sigma\beta=-\beta$ without rationalization. The nontrivial embedding is complex conjugation, and on $\C/\R(2)$ this negates the nonzero regulator. Thus $r_2(\beta)=2r_2(b_0)\neq 0$.
\end{proof}

Define
\begin{equation}\label{d3:eq:c}
c=i_*(\beta|_R)\in K_3(A).
\end{equation}
The pushforward exists because $A$ and $R$ are regular, and the closed immersion is a regular Cartier immersion.

\begin{lemma}\label{d3:lem:genericzero}
The image of $c$ in $K_3(\Frac A)$ is zero integrally.
\end{lemma}
\begin{proof}
By construction, $c$ comes from a class with support on $V(x)$. The localization sequence therefore sends it to zero in $K_3(A[1/x])$, and hence in the fraction field.
\end{proof}

The remaining question is whether $c$ has infinite order. Merely knowing that $V(x)$ is principal does not answer this: a trivial normal bundle implies the vanishing of the self-intersection operator $i^*i_*$, not the vanishing of $i_*$ itself.

\subsection{The arithmetic node and its divisor class}\label{d3:sec:node}

\begin{lemma}\label{d3:lem:normalB}
The ring $B$ is a normal local domain of dimension three, with maximal ideal
\[
\n=(u,x,y,z).
\]
It is regular away from its closed point, and
\[
\operatorname{gr}_{\n}B
\cong\F_p[U,X,Y,Z]/(XY+Z^2-U^2).
\]
\end{lemma}
\begin{proof}
The map $A\to B$ is finite free of rank two. Its closed fiber is $\F_p[u]/(u^2)$, with a unique maximal ideal, so $B$ is local. In particular, $B=(B_0)_{\n}$, not merely an unspecified localization of $B_0$. The polynomial $u^2+p$ is irreducible over the purely transcendental extension $\Q(x,z)$, hence $B$ is a domain.

At a point where $u$ is invertible, the ring is on a smooth affine quadric over $E$. At a point of the special fiber other than the origin, one of $x,y,z$ is invertible, and the relative gradient $(y,x,2z)$ has a unit component. Thus the punctured spectrum is regular. Being a hypersurface, $B$ satisfies $S_2$; regularity in codimension one gives normality. The displayed associated graded ring follows from the nonzero quadratic initial form of the equation in the regular ambient local ring $W[x,y,z]_{(u,x,y,z)}$.
\end{proof}

Consider the height-one primes
\[
P_+=(x,z-u),\qquad P_-=(x,z+u).
\]
Both have multiplicity one in $\divisor_B(x)$, so
\begin{equation}\label{d3:eq:divx}
\divisor_B(x)=P_++P_-.
\end{equation}

\begin{lemma}\label{d3:lem:smallresolution}
Let $W'$ be a DVR with uniformizer $\varpi$ and residue field $k$, and let $0\neq d\in\varpi W'$. In the normal local domain
\[
C=\left(W'[x,y,w]/(xy+w(w+d))\right)_{(\varpi,x,y,w)},
\]
the height-one prime $P=(x,w)$ has infinite order in $\Cl(C)$. There is no restriction on the residue characteristic or on the positive valuation of $d$.
\end{lemma}
\begin{proof}
The defining polynomial is irreducible, and $C$ is a hypersurface. Its generic fiber is smooth since $d\neq0$. On the special fiber, a prime containing $x$ and $y$ also contains $w$; hence every nonclosed point of that fiber has $x$ or $y$ invertible. Thus $C$ is regular off its closed point, and $S_2$ and $R_1$ give normality.

Blow up the ideal $I=(x,w)$:
\[
\rho:Y=\operatorname{Bl}_{I}(\Spec C)\longrightarrow\Spec C.
\]
Its two standard charts are localizations of the following polynomial rings:
\begin{align*}
W'[x,t],&\qquad w=xt,\quad y=-t(xt+d),\\
W'[q,y],&\qquad w=-qy-d,\quad x=-q(qy+d).
\end{align*}
These are the actual Rees charts $C[w/x]$ and $C[x/w]$ inside the fraction field. For example, introducing $t=w/x$ gives $x(y+t(xt+d))=0$, and $x$-saturation gives exactly the first polynomial ring. The second is obtained by $w$-saturation. There are no extra equations because each indicated pair is a transcendence basis over $\Frac(W')$.

Consequently $Y$ is regular. The closed fiber has equations $(\varpi,x)$ in the first chart and $(\varpi,y)$ in the second. Its two affine charts are $k[t]$ and $k[q]$, with $q=t^{-1}$ on their overlap. Hence it is a reduced curve
\[
F\cong\PP^1_k.
\]

The ideal $I$ is invertible off the closed point. Outside $V(I)$ this is immediate. At a nonclosed prime containing $I$, either $y$ is a unit, giving $I=(w)$, or $\varpi$ is a unit; then $d$ and $w+d$ are units and $I=(x)$. The blowup is therefore an isomorphism off the closed point. Its only exceptional locus is the codimension-two curve $F$, so there are no exceptional prime divisors.

The standard tautological identity for a blowup is $I\OO_Y=\OO_Y(1)$, whose restriction to $F$ is $\OO_{\PP^1_k}(1)$; see \cite[Lemma~31.33.4, Tag~02OS]{StacksBlowup}. Since $I$ has multiplicity one at the generic point of $P$ and there are no exceptional divisors,
\[
I\OO_Y=\OO_Y(-\widetilde P),\qquad
\OO_Y(\widetilde P)|_F=\OO_{\PP^1_k}(-1).
\]

If $\divisor_C(f)=NP$ for a nonzero integer $N$, then the same rational function on $Y$ has divisor exactly $N\widetilde P$. Restricting its trivial divisor line bundle to $F$ gives
\[
\OO_{\PP^1_k}(-N)\cong\OO_{\PP^1_k},
\]
which forces $N=0$, a contradiction.
\end{proof}

\begin{corollary}\label{d3:lem:clnonzero}
The class $[P_+]$ is not torsion in $\Cl(B)$.
\end{corollary}
\begin{proof}
Set $w=z-u$ and $d=2u$. Then $B$ is the ring in Lemma~\ref{d3:lem:smallresolution}, with $W'=W$ and $\varpi=u$, and $P_+=(x,w)$.
\end{proof}

\begin{remark}
Blowing up the maximal ideal instead gives the regular resolution with exceptional quadric $\PP^1_{\F_p}\times\PP^1_{\F_p}$. Its normal bundle is $\OO(-1,-1)$, while the strict transform of $P_+$ restricts to one ruling, $\OO(1,0)$. If $NP_+$ were principal, its divisor upstairs would have the form $N\widetilde P_++aF$, whose restriction would be $\OO(N-a,-a)$. Triviality again forces $N=0$. The small resolution above avoids this larger exceptional divisor and also works in residue characteristic two.
\end{remark}

\begin{lemma}\label{d3:lem:invertu}
Restriction induces an isomorphism
\[
\Cl(B)\xrightarrow{\sim}\Cl(B[1/u]).
\]
In particular, the image of $[P_+]$ is not torsion in $\Pic(T)$.
\end{lemma}
\begin{proof}
The quotient
\[
B/(u)=\F_p[x,y,z]_{(x,y,z)}/(xy+z^2)
\]
is a domain. The only height-one prime removed by inverting $u$ is therefore the principal prime $(u)$. The localization sequence for divisor class groups is surjective with kernel generated by the removed prime classes, so its kernel is zero. Since $B[1/u]$ is regular, its class group equals its Picard group.
\end{proof}

\begin{remark}
One can also compute $\Cl(B)=\Z$. Indeed, $B[1/x]$ is a localization of the unique factorization domain $W[x,x^{-1},z]$, so the class group is generated by $P_+$ and $P_-$. Equation~\eqref{d3:eq:divx} gives their sum as zero, and Lemma~\ref{d3:lem:clnonzero} excludes any further relation. The argument needs only non-torsion, not this complete computation.
\end{remark}

\subsection{Persistence of Picard and Betti classes}\label{d3:sec:persistence}

Let $D_+$ and $D_-$ denote the divisors of $Q$ with equations
\[
D_+=(x,z-u),\qquad D_-=(x,z+u).
\]
They are disjoint affine lines over $E$.

\begin{lemma}\label{d3:lem:quadricpic}
We have $\Pic(Q)=\Z[D_+]$, and the same statement holds after extension to $\overline E$ or $\C$. The absolute Galois group of $E$ acts trivially on the geometric Picard group. At a complex embedding,
\[
H^2(Q(\C),\Z)=\Z,
\]
with $c_1(\OO_Q(D_+))$ a primitive generator after the usual Tate normalization.
\end{lemma}
\begin{proof}
The projective closure
\[
\overline Q=\{xy+z^2=u^2t^2\}\subset\PP^3_E
\]
is split: it has the rank-one matrix presentation
\[
\det\begin{pmatrix}x&ut-z\\ut+z&y\end{pmatrix}=0.
\]
Thus $\overline Q\cong\PP^1\times\PP^1$. The divisor at infinity is a smooth divisor of type $(1,1)$; after an automorphism of one factor it is the diagonal. Hence
\[
Q\cong (\PP^1\times\PP^1)\setminus\Delta,
\qquad
\Pic(Q)=\Z^2/\Z(1,1)=\Z.
\]
The closure of $D_+$ is a ruling, and gives a generator. This calculation is unchanged over algebraically closed extensions. Because the generator is defined over $E$, the absolute Galois action is trivial.

Projection to one factor makes $Q(\C)\to\PP^1(\C)$ a locally trivial affine-line bundle, so it is a homotopy equivalence. Equivalently, the topological localization sequence for the diagonal gives the same quotient $\Z^2/\Z(1,1)$ in degree two. The first Chern class of a ruling maps to its primitive generator.
\end{proof}

\begin{lemma}\label{d3:lem:Picinject}
The restriction map $\Pic(Q)\to\Pic(T)$ is injective.
\end{lemma}
\begin{proof}
The generator $[D_+]$ maps to the image of $[P_+]$ in $\Cl(B[1/u])$. The latter has infinite order by Lemmas~\ref{d3:lem:clnonzero} and \ref{d3:lem:invertu}.
\end{proof}

\begin{lemma}\label{d3:lem:components}
For each $s\in S$, every geometrically irreducible curve in $Q_{\C}\setminus(Q_s)_{\C}$ has zero class in $\Pic(Q_{\C})$.
\end{lemma}
\begin{proof}
Let $D$ be an irreducible divisor over $E$ in $Q\setminus Q_s$. It is absent from $T$, so $\OO_Q(D)|_T$ is trivial. Lemma~\ref{d3:lem:Picinject} implies $[D]=0$ in $\Pic(Q)$.

In characteristic zero the reduced divisor $D$ becomes a reduced union of finitely many geometric irreducible components. These components are permuted transitively by the absolute Galois group of $E$. Its action on $\Pic(Q_{\overline E})=\Z$ is trivial, so all components have the same integer class, say $a$. If there are $r>0$ components, their sum has class $ra$. This sum is the pullback of $[D]=0$, so $ra=0$ in the torsion-free group $\Z$, and therefore $a=0$. Extending further to $\C$ preserves the calculation.
\end{proof}

\begin{remark}
A geometric boundary component $\Gamma$ is defined over some finite
extension $E'/E$. Put $L=\OO_Q(D_+)$ and write
$[\OO(\Gamma)]=a[L_{E'}]$ in $\Pic(Q_{E'})=\Z[L_{E'}]$.
This line bundle is trivial on $T_{E'}$. Its finite-flat norm
(equivalently, determinants after restriction of scalars) therefore
makes $L_T^{a[E':E]}$ trivial. The infinite order of $L_T$ forces
$a=0$, giving a second proof that each geometric boundary component has zero Picard class.
\end{remark}

\begin{remark}
It is not enough that $\divisor(s)$, as a whole, is principal. For instance $\divisor(x)=D_++D_-$ is principal, while its two components have opposite nonzero classes. Removing $D_+\cup D_-$ kills $H^2(Q,\R)$. The element $x$ is not in the permitted multiplicative set $S$.
\end{remark}

\begin{lemma}\label{d3:lem:Betti}
For every $s\in S$, restriction is injective:
\[
H^2(Q_{\C},\R)\longrightarrow H^2((Q_s)_{\C},\R).
\]
\end{lemma}
\begin{proof}
Let $Z=Q_{\C}\setminus(Q_s)_{\C}$ with its reduced support. The kernel is the image of the support map
\[
H_Z^2(Q_{\C},\R)\longrightarrow H^2(Q_{\C},\R).
\]
Remove from the ambient surface the finite set of singular points and intersections of the reduced curve components of $Z$. Support cohomology in finitely many points of a real four-manifold vanishes below degree four, so this changes neither the degree-two support group nor the ambient $H^2$. The remaining components are disjoint smooth connected complex curves.
The tubular-neighborhood calculation in the proof of
Proposition~\ref{d2:prop:geometry} identifies degree-two support
cohomology with one copy of $\R$ for each component.
There the canonical transverse section of the divisor line bundle
identifies the image of its support generator with the global first
Chern class of the component. The isomorphism on ambient $H^2$
transports this identification across the deleted finite set.
Every such class is zero by Lemma~\ref{d3:lem:components}.
Thus the support map is zero.
\end{proof}

\subsection{Analytic Deligne detection}\label{d3:sec:Delignedetection}

Throughout this section, $\HD$ denotes analytic Deligne cohomology.

\begin{lemma}\label{d3:lem:Deligne}
For a smooth complex surface $U$, there is a natural isomorphism
\begin{equation}\label{d3:eq:Delignequotient}
\HD^3(U,\R(3))
\cong H^2(U,\C)/H^2(U,\R(3)).
\end{equation}
\end{lemma}
\begin{proof}
The analytic Deligne complex is
\[
[\R(3)\longrightarrow\OO_U\longrightarrow\Omega_U^1\longrightarrow\Omega_U^2].
\]
The holomorphic de Rham complex resolves the constant sheaf $\C$ by the holomorphic Poincar\'e lemma, and there are no holomorphic forms in degree three. The Deligne complex is consequently quasi-isomorphic to the constant sheaf $\C/\R(3)$ placed in degree one. Equivalently, it is the shifted cone of
\[
R\Gamma(U,\R(3))\longrightarrow R\Gamma(U,\C).
\]
Since real cohomology injects into its complexification in every degree, the long exact sequence of this cone gives \eqref{d3:eq:Delignequotient}. All constructions are natural for restriction to open subsets. Properness is not used.
\end{proof}

\begin{remark}
On a smooth normal-crossings compactification, the logarithmic de Rham complex also has $F^3=0$. Thus Deligne--Beilinson cohomology agrees with the analytic group in \eqref{d3:eq:Delignequotient}; compare \cite[Remark~2.13]{EsnaultViehweg}. The direct analytic formulation above already suffices for the proof.
\end{remark}

\begin{lemma}\label{d3:lem:Deligneinject}
For every $s\in S$, restriction induces an injection
\[
\HD^3(Q_{\C},\R(3))\longrightarrow\HD^3((Q_s)_{\C},\R(3)).
\]
\end{lemma}
\begin{proof}
Combine Lemmas~\ref{d3:lem:Betti} and \ref{d3:lem:Deligne}. Explicitly, an injection of real vector spaces $V\hookrightarrow W$ induces an injection
\[
\frac{V\otimes_{\R}\C}{(2\pi i)^3V}
\longrightarrow
\frac{W\otimes_{\R}\C}{(2\pi i)^3W},
\]
because the intersection of the image of $V\otimes\C$ with $(2\pi i)^3W$ is exactly $(2\pi i)^3V$.
\end{proof}

Set
\[
\eta=[\OO_{D_+}]\in K_0(Q),\qquad
\gamma=(\beta|_E)\eta\in K_3(Q).
\]
The structure sheaf class is a perfect-complex class, with
\[
\eta=1-[\OO_Q(-D_+)],\quad \operatorname{rank}(\eta)=0,\quad
\operatorname{ch}_1(\eta)=c_1(\OO_Q(D_+)).
\]

\begin{lemma}\label{d3:lem:product}
The class $\gamma$ has nonzero weight-three analytic Deligne regulator.
\end{lemma}
\begin{proof}
Naturality and multiplicativity give
\begin{equation}\label{d3:eq:product}
r_3(\gamma)=r_2(\beta|_E)\cup c_1(\OO_Q(D_+)).
\end{equation}
There are no additional terms of total weight three: a degree-three $K$-class pulled back from a point has the relevant regulator only in $H^1_{\mathcal D}(\mathrm{pt},\R(2))$, and the rank of $\eta$ is zero.

Let $h$ be a real primitive generator of $H^2(Q_{\C},\R)$, normalized so that $c_1(\OO_Q(D_+))$ has Betti image $2\pi i\,h$. A nonzero element of $\C/\R(2)$ is represented by $it$ for a nonzero real $t$. Under \eqref{d3:eq:Delignequotient}, the product in \eqref{d3:eq:product} is represented, up to the harmless convention for the sign of $h$, by
\[
(it)(2\pi i)h=-2\pi t\,h.
\]
This is a nonzero real multiple of $h$, whereas $H^2(Q_{\C},\R(3))$ is the imaginary real line generated by $h$. The quotient class is therefore nonzero.

One may equivalently compute this cup product on $\PP^1$: the projection of the affine-line torsor $Q\to\PP^1$ induces the required degree-two Betti isomorphism, and cup product with the hyperplane class is the usual projective-bundle map.
\end{proof}

\begin{proposition}\label{d3:prop:infinite}
The restriction $\gamma|_T\in K_3(T)$ has infinite order.
\end{proposition}
\begin{proof}
By regulator naturality and Lemma~\ref{d3:lem:Deligneinject},
\[
r_3(\gamma|_{Q_s})\neq 0
\qquad\text{for every }s\in S.
\]
If a positive integer multiple $N\gamma$ vanished in $K_3(T)$, continuity \eqref{d3:eq:continuity} would make it vanish already in $K_3(Q_s)$ for some $s$. Applying the regulator there would yield $N r_3(\gamma|_{Q_s})=0$ in a real vector space, which is impossible.
\end{proof}

\subsection{Flat base change and nonvanishing}

\begin{lemma}\label{d3:lem:basechange}
The image of the supported class \eqref{d3:eq:c} in $K_3(T)$ equals $2\gamma|_T$.
\end{lemma}
\begin{proof}
The map $A\to B[1/u]$ is flat, being a finite free base change followed by localization. Its pullback of the divisor $x=0$ is the disjoint union
\[
(D_+\cap T)\amalg(D_-\cap T).
\]
The original coefficient map $W\to R$ sends $u$ to $z$. Therefore its two restrictions after the splitting are the coefficient maps $u\mapsto u$ and $u\mapsto -u$. Flat base change and the projection formula give
\begin{align*}
c|_T
&=(\beta|_E)[\OO_{D_+\cap T}]+(\sigma\beta|_E)[\OO_{D_-\cap T}]\\
&=(\beta|_E)\bigl([\OO_{D_+\cap T}]-[\OO_{D_-\cap T}]\bigr).
\end{align*}
The two divisors are disjoint since $2u$ is a unit on $T$. Moreover their union is the principal divisor $V(x)$. The resolution by multiplication by $x$ gives
\[
[\OO_{D_+\cap T}]+[\OO_{D_-\cap T}]
=[\OO_{T/(x)}]=0\quad\text{in }K_0(T).
\]
Consequently $c|_T=2(\beta|_E)[\OO_{D_+\cap T}]=2\gamma|_T$.
\end{proof}

\begin{proof}[Proof of Theorem~\ref{d3:thm:main}]
Lemma~\ref{d3:lem:regularA} gives the required regular local ring. By Proposition~\ref{d3:prop:infinite} and Lemma~\ref{d3:lem:basechange}, the pullback of $c$ to $K_3(T)$ has infinite order; hence $c$ itself has infinite order. Lemma~\ref{d3:lem:genericzero} shows that its generic image is zero integrally. Tensoring with $\Q$ therefore gives the asserted nonzero rational kernel.
\end{proof}

The nonzero-product assertion in this proof is not deduced from a nonzero Picard class alone. The essential extra step is that its regulator remains nonzero \emph{after every permitted finite localization}. This is what makes the filtered-colimit argument valid.

\subsection{The unsplit generic fiber}

The following calculation identifies the coefficient of the supported class on the unsplit generic fiber. Put
\[
X=\Spec\Q[x,y,z]/(xy+z^2+p).
\]
Its divisor $Z=V(x)$ is $\A^1_E$, while
\[
X\setminus Z\cong\mathbb G_{m,\Q}\times\A^1_{\Q}.
\]
For every $n\geq 0$, localization, homotopy invariance, and the fundamental theorem give
\[
K_n(E)\longrightarrow K_n(X)\longrightarrow
K_n(\Q)\oplus K_{n-1}(\Q)
\xrightarrow{\partial}K_{n-1}(E).
\]
The boundary is zero on $K_n(\Q)$ and is restriction $\operatorname{res}_{E/\Q}$ on the second summand: that summand is represented by $\{x\}\cdot b$, and $x$ cuts $Z$ with multiplicity one. The preceding boundary has image $\operatorname{res}K_n(\Q)$ for the same reason.

After rationalization, restriction is injective because transfer composed with restriction is multiplication by two. Constants split the resulting short exact sequence, giving
\begin{equation}\label{d3:eq:genericdecomposition}
K_n(X)\Qrat\cong K_n(\Q)\Qrat\oplus K_n(E)\Qrat^{-}.
\end{equation}
Here the second summand is identified using $\operatorname{res}\operatorname{Tr}=1+\sigma$ and enters through $Z$-Gysin. In degree three, Borel's calculation \cite[Proposition~12.2]{Borel} yields
\[
K_3(\Q)\Qrat=0,\qquad K_3(E)\Qrat^{-}\cong\Q.
\]
Thus the generic-fiber Gysin map identifies $K_3(E)\Qrat$ with $K_3(X)\Qrat$. The supported class is already a nonzero one-dimensional summand before localization at $\m$. The main proof shows that this particular summand cannot disappear at that localization.

\section{The explicit three-adic class and its complex period}\label{d3:sec:explicit}

This section gives a geometric proof of the explicit class at $p=3$. It uses the same geometric persistence lemmas, but obtains nonvanishing by an explicit period rather than by selecting a Borel class. Put
\[
A=\left(\Z_{(3)}[x,y,z]/(3+xy+z^2)\right)_{(x,y,z)},
\]
and define
\[
a=\frac{1+z}{2},\qquad s_0=a(1-a)=1+\frac{xy}{4},\qquad
\lambda=1+x,\qquad \mu=1+\frac{(1+x)y}{4}.
\]
Every one of $a,1-a,s_0,\lambda,\mu$ is a unit of $A$, and
\begin{equation}\label{d3:eq:unitsidentity}
1+x\mu=\lambda s_0.
\end{equation}

\subsection{Dennis--Stein notation and integral generic vanishing}
We use the modern Dennis--Stein convention: $\langle r,t\rangle$ is defined if $1-rt$ is a unit, and, if $r$ is a unit,
\[
\langle r,t\rangle=\{r,1-rt\}.
\]
If $t$ is a unit, antisymmetry gives $\langle r,t\rangle=\{1-rt,t\}$. The relation used below is
\begin{equation}\label{d3:eq:DS2}
\langle r,t\rangle+\langle r,v\rangle
=\langle r,t+v-rtv\rangle,
\qquad \langle r,1\rangle=0.
\end{equation}
The modern convention and relations are those of \cite[III, Definition~5.11 and (D1)--(D3)]{Weibel}; the historical source is \cite{DennisStein}. We do not need a presentation theorem for $K_2$.

Let
\[
\delta=\langle -x,y/4\rangle\in K_2(A),\qquad
\alpha=[a]\delta\in K_3(A).
\]
On $A[1/x]$, the symbol is $\{-x,s_0\}$. Since $a$ and $1-a$ are units, the Steinberg relation and graded commutativity give
\[
2\{a,s_0\}=2\{a,a(1-a)\}
=2[a]^2+2\{a,1-a\}=0.
\]
Reordering products therefore gives
\begin{equation}\label{d3:eq:explicitgeneric}
2\alpha|_{A[1/x]}=2\{a,-x,s_0\}=0
\qquad\text{integrally}.
\end{equation}

The Dennis--Stein class can also be rewritten using units alone. By \eqref{d3:eq:DS2},
\[
\delta=\langle -x,\mu\rangle=\{1+x\mu,\mu\}
=\{\lambda s_0,\mu\}.
\]
Multiplying by $2[a]$ kills the $\{a,s_0,\mu\}$ term and yields
\begin{equation}\label{d3:eq:explicitDS}
2\alpha=2\{a,\lambda,\mu\}=c_{\mathrm{exp}}.
\end{equation}
Thus the explicitly displayed class \eqref{d3:eq:explicitfront} has zero generic image integrally.

\subsection{A finite open on which to calculate the period}
Keep $E=\Q(u)$ with $u^2=-3$, and fix $u=i\sqrt3$. Let
\[
Q=\Spec E[x,y,z]/(3+xy+z^2),\qquad
U=Q[(1-z^2)^{-1}].
\]
The classes $\delta$ and $\alpha$ are defined on $U$, since $a,1-a,s_0$ are units there. Both lines $D_+$ and $D_-$ lie entirely inside $U$, because on them $1-z^2=4$. Equation~\eqref{d3:eq:explicitgeneric} holds on $U[1/x]$ as well.

Localization therefore provides a class supported on $D=V(x)$ that maps to $2\alpha\in K_3(U)$. Since the entire divisor $D$ lies in $U$, excision identifies the supported groups for $Q$ and $U$. Equivalently, apply the localization sequence with $K_3(D)$ and the naturality of its pushforward. Thus there is a class
\begin{equation}\label{d3:eq:explicitlift}
\widetilde c\in K_3(Q),\qquad \widetilde c|_U=2\alpha.
\end{equation}
No coefficient of this supported lift has to be chosen or computed explicitly.

\subsection{The curvature and the Deligne product}
The weight-two analytic regulator of $\delta$ has the holomorphic curvature
\begin{equation}\label{d3:eq:omega}
\omega=\frac{dx\wedge d(y/4)}{1+xy/4}.
\end{equation}
Indeed, where $x\neq0$, multiplicativity for the two unit classes gives
\[
d\log(-x)\wedge d\log s_0
=\frac{dx\wedge d(y/4)}{s_0}.
\]
The Hodge endpoint is a canonical holomorphic two-form: in this degree $F^2A^1=0$, so there is no preceding filtered form that could change it by a boundary. Both sides therefore extend as actual holomorphic two-forms on $U$, and agreement on the dense open proves \eqref{d3:eq:omega}. The logarithmic normalization follows from the ordinary unit regulator and \cite[Theorem~2.31 and Lemma~4.5]{BunkeTamme}. An independent algebraic differential calculation for Dennis--Stein characters over smooth commutative $\Q$-algebras, with its own symbol conventions, appears in \cite[Section~3.2, Theorem~3.4 and Corollary~3.6]{Ginot}; the preceding dense-open calculation fixes the regulator sign and gives the needed form.

Suppose that $a$ has a single-valued holomorphic logarithm $L$ on an analytic open $N\subset U(\C)$. Then the regulator of $\alpha=[a]\delta$ on $N$ is represented in
\[
\HD^3(N,\R(3))=H^2(N,\C/\R(3))
\]
by the closed holomorphic two-form
\begin{equation}\label{d3:eq:Lomega}
-L\omega.
\end{equation}
We use the cone-coordinate convention of \cite[Proposition~2.8]{BunkeTamme}: an interval form $\Theta$ maps to
\[
\bigl(\Theta|_0\oplus\Theta|_1,-\textstyle\int_{[0,1]}\Theta\bigr).
\]
Let $v$ denote the auxiliary interval coordinate and put $g=i\operatorname{Im}L+v\operatorname{Re}L$.
By \cite[Lemma~4.5]{BunkeTamme}, the unit class $[a]$ is represented by $dg$.
We may replace this by $d(vL)$ in the same absolute interval complex:
\[
dg-d(vL)=d\bigl((1-v)i\operatorname{Im}L\bigr).
\]
The primitive is allowed in $IDR(1)$: its endpoint at zero is
$\R(1)$-valued, and its endpoint at one is zero, as required by
$F^1A^0=0$. Thus this replacement does not substitute a different
relative class for the original absolute regulator.
Choose a closed interval representative $\Theta_\delta$ whose Hodge endpoint is $\omega$.
The product $d(vL)\wedge\Theta_\delta$ has zero endpoint at zero and
endpoint $dL\wedge\omega=0$ at one, since $N$ is a complex surface.
Multiplicativity and Stokes' formula therefore give
\[
-\int_{[0,1]}d(vL)\wedge\Theta_\delta
=-L\omega+d\int_{[0,1]}vL\Theta_\delta.
\]
Thus the cone coordinate is \eqref{d3:eq:Lomega}, for every chosen
closed representative of the original absolute class. The form
$L\omega$ is closed because $N$ has complex dimension two.

\subsection{A strictly positive real period}
Inside $U(\C)$ consider the smooth sphere
\[
\Sigma:\quad z=i\sqrt3\,\tau,\quad
x=\sqrt3\sqrt{1-\tau^2}\,e^{i\vartheta},\quad
y=-\overline x,\qquad -1\leq\tau\leq1.
\]
The polar coordinates degenerate only at the poles; intrinsically this is the sphere $|x|^2+(\operatorname{Im}z)^2=3$. Along it,
\[
a=\frac{1+i\sqrt3\,\tau}{2},\qquad
s_0=\frac{1+3\tau^2}{4}>0.
\]
There is an analytic neighborhood $N$ of $\Sigma$ on which $\operatorname{Re}a>0$, so the principal $L=\log a$ is single-valued. On $\Sigma$,
\[
L=\frac12\log\frac{1+3\tau^2}{4}+i\arctan(\sqrt3\,\tau),
\qquad
\omega=-\frac{6i\tau}{1+3\tau^2}\,d\tau\wedge d\vartheta.
\]
Consequently,
\begin{equation}\label{d3:eq:positiveperiod}
\int_{\Sigma}\operatorname{Re}(L\omega)
=2\pi\int_{-1}^{1}
\frac{6\tau\arctan(\sqrt3\,\tau)}{1+3\tau^2}\,d\tau>0.
\end{equation}
The integrand is positive except at $\tau=0$. With the specified Bunke--Tamme cone convention, the real period of $r_3(2\alpha)$ is the negative double of this positive integral. Since $\R(3)$ is the imaginary axis, this detects a nonzero class in $H^2(N,\C/\R(3))$. Thus $r_3(2\alpha)\neq0$ on $U$.

For comparison, integration by parts expresses \eqref{d3:eq:positiveperiod} as
\[
4\pi D(e^{i\pi/3}),
\]
with the displayed orientation and logarithmic convention. Strict positivity alone is sufficient; neither the special value nor Borel's rank theorem is needed in this argument.

\subsection{A normalized supported coefficient}\label{d3:sec:exactcoefficient}
The exact coefficient can also be specified rationally. Put $\zeta=(1+u)/2$, and define $\beta_{\mathrm{BT}}=\tfrac12\operatorname{bl}(2[\zeta])\in K_3(E)_{\Q}$ using the Bloch symbol map of \cite[Theorem~4.9 and (4.8)]{BunkeTamme}, first on the full integer ring $\OO_E=\Z[\zeta]$ and then by restriction to $E$. Both $\zeta$ and $1-\zeta=\zeta^{-1}$ are units. The doubled wedge boundary is zero even with the antisymmetric-tensor convention, so this definition avoids any integral two-torsion convention for the Bloch group. At the chosen embedding its cone-coordinate regulator is $-iD(\zeta)$. Borel injectivity \cite[Section~6]{BunkeTamme} makes this rational class independent of the integral lift.

A supported lift as in \eqref{d3:eq:explicitlift} has the form $\widetilde c=\gamma\eta_+$, with $\gamma\in K_3(E)$: use homotopy invariance on the two affine lines and $\eta_-=-\eta_+$. The two generators of $I=(x,z-u)$ give a map $f:Q\to\mathbf P^1_E$ with $f^*\OO(1)=I$. On $\Sigma$ its coordinate is
\[
\frac{z-u}{x}=-i\sqrt{\frac{1-\tau}{1+\tau}}e^{-i\vartheta}.
\]
Both the radial derivative and the angular derivative are negative, so this is a map of degree $+1$ for the displayed orientation. Hence
\[
\int_\Sigma c_1(I)=2\pi i,\qquad
\int_\Sigma\operatorname{ch}_1(\eta_+)=-2\pi i.
\]
Multiplicativity gives period $(-iD(\zeta))(-2\pi i)=-2\pi D(\zeta)$ for $\beta_{\mathrm{BT}}\eta_+$, whereas \eqref{d3:eq:positiveperiod} gives period $-8\pi D(\zeta)$ for $\widetilde c$. Borel injectivity on $K_3(E)_{\Q}$ therefore gives
\begin{equation}\label{d3:eq:exactcoefficient}
\widetilde c=4\beta_{\mathrm{BT}}\eta_+
\quad\text{in }K_3(Q)_{\Q}.
\end{equation}
This is a regulator-normalized identity for every supported lift. It does not identify $\beta_{\mathrm{BT}}$ with the Besser--de Jeu class used in Section~\ref{d3:sec:explicitcomplete}; that latter comparison will only use a nonzero rational scalar.

\subsection{Persistence and the explicit conclusion}
By \eqref{d3:eq:explicitlift}, the regulator of $\widetilde c$ is nonzero on $Q$. Lemma~\ref{d3:lem:Deligneinject} keeps it nonzero on every $Q_s$ for $s\in S$. The element $1-z^2$ belongs to $S$, so $T$ lies in $U$ and the restrictions of $\widetilde c$ and $2\alpha$ agree on $T$. By \eqref{d3:eq:explicitDS}, this is also the image of $c_{\mathrm{exp}}\in K_3(A)$.

If a positive multiple of $c_{\mathrm{exp}}$ vanished in $K_3(A)$, the same multiple of $\widetilde c$ would vanish in $K_3(T)$, and hence on some finite $Q_s$ by continuity. Its nonzero real regulator rules this out. We have therefore obtained the explicit version of the counterexample:
\begin{equation}\label{d3:eq:explicitconclusion}
\boxed{\quad
0\neq 2\left\{\frac{1+z}{2},\,1+x,\,
1+\frac{(1+x)y}{4}\right\}\in
\ker\bigl(K_3(A)\longrightarrow K_3(\Frac A)\bigr),
\quad}
\end{equation}
where nonzero is strengthened to infinite order by the period argument. This conclusion depends on the geometric persistence proof in Sections~\ref{d3:sec:node}--\ref{d3:sec:Delignedetection}, in addition to the explicit calculation in this section.

\subsection{An independent algebraic detection modulo three}\label{d3:sec:modthree}

There is a short independent test that proves $c_{\mathrm{exp}}$ is nonzero integrally, and even that it is not divisible by three. This test does not prove infinite order, so it does not replace the regulator argument for the rational conjecture.

Let
\[
R=A[T]_{\m A[T]},\qquad V'=\Z_{(3)}[T]_{(3)}.
\]
Then $R$ is a local $V'$-algebra with residue field $\F_3(T)$. Its module of relative differentials is presented by
\[
\Omega^1_{R/V'}=
\frac{R\,dx\oplus R\,dy\oplus R\,dz}
 {R(y\,dx+x\,dy+2z\,dz)},
\]
so
\[
\Omega^3_{R/V'}=
\bigl(R/(x,y,2z)\bigr)\,dx\wedge dy\wedge dz
\cong\F_3(T)\,dx\wedge dy\wedge dz.
\]

For a local ring with infinite residue field, the usual Milnor symbol product
\[
\iota:K^M_3(R)\longrightarrow K_3(R)
\]
has a comparison homomorphism $\phi:K_3(R)\to K^M_3(R)$ satisfying $\phi\iota=2$. This is the degree-three case of the Nesterenko--Suslin comparison recorded for general local rings in \cite[Proposition~4.2.8(5),(6)]{KerzThesis}, after identifying ordinary and improved Milnor $K$-theory for the infinite residue field of $R$.

Write $c_M=2\{a,\lambda,\mu\}\in K^M_3(R)$. The logarithmic differential is well defined on Milnor symbols, since it kills every Steinberg relation. A direct calculation gives
\[
d\log(c_M)
=\frac{1}{2(1+z)\mu}\,dz\wedge dx\wedge dy.
\]
Consequently the composite additive homomorphism
\[
K_3(A)\longrightarrow K_3(R)
\xrightarrow{\phi}K^M_3(R)
\xrightarrow{d\log}\Omega^3_{R/V'}
\]
sends $c_{\mathrm{exp}}$ to
\[
\frac{1}{(1+z)\mu}\,dx\wedge dy\wedge dz
=dx\wedge dy\wedge dz\neq0.
\]
The equality is in $\Omega^3_{R/V'}$, where $x,y,z$ vanish in the coefficient ring and $\mu=1$. Since this target is killed by three, $c_{\mathrm{exp}}\notin3K_3(A)$.

In particular, the coefficient exact sequence injects the nonzero image of this class in $K_3(A)/3$ into $K_3(A;\Z/3)$. Its generic image is zero by the integral calculation in \eqref{d3:eq:explicitgeneric}. Thus this explicit example also fails Gersten injectivity with $\Z/3$-coefficients, independently of the analytic regulator.

\begin{remark}
The same conclusion is already visible in
\[
C=A/\m^2\cong
\F_3[\epsilon_x,\epsilon_y,\epsilon_z]/(\epsilon_x,\epsilon_y,\epsilon_z)^2.
\]
To check this without assuming that the finite residue field is large enough for an ordinary Milnor comparison, pass to
\[
D=C[T]_{\m_C C[T]}
\cong k[\epsilon_x,\epsilon_y,\epsilon_z]/(\epsilon_x,\epsilon_y,\epsilon_z)^2,
\qquad k=\F_3(T).
\]
Let $M=k e_x\oplus k e_y\oplus k e_z$ be a $D$-module through the residue map. The $k$-derivation $d:D\to M$ given by $d\epsilon_i=e_i$ induces
\[
K_3^M(D)\longrightarrow\bigwedge_k^3 M,\qquad
\{g_1,g_2,g_3\}\longmapsto
\frac{dg_1}{g_1}\wedge\frac{dg_2}{g_2}\wedge\frac{dg_3}{g_3}.
\]
This kills the Steinberg relations because $dg/g$ and $d(1-g)/(1-g)$ are proportional. The three entries of $c_{\mathrm{exp}}$ map to $2(1+\epsilon_z)$, $1+\epsilon_x$, and $1+\epsilon_y$. After the comparison $K_3(D)\to K_3^M(D)$, the image of $c_{\mathrm{exp}}$ is therefore
\[
4e_z\wedge e_x\wedge e_y=e_x\wedge e_y\wedge e_z\ne0.
\]
In fact its image in $K_3(C)$ has exact order three: it is nonzero by this test, and $(1+\epsilon_x)^3=1$ kills its triple product by three. Both the henselization and the completion of $A$ have the same quotient $C$. Hence the explicit class also remains nonzero and not divisible by three in their $K_3$-groups, with zero generic image. This observation alone does not prove infinite order in either ring.
\end{remark}

\section{Quadratic and Eisenstein families}\label{d3:sec:families}

\subsection{Residue characteristic two}\label{d3:sec:two}

The small resolution in Lemma~\ref{d3:lem:smallresolution} allows the same argument at the prime two. This variant is not needed for the explicit counterexample in \eqref{d3:eq:explicitfront}.

\begin{proposition}\label{d3:prop:two}
For the regular local ring
\[
A_2=\left(\Z_{(2)}[x,y,z]/(1+xy+z^2)\right)_{(2,x,y,z-1)},
\]
the map $K_3(A_2)\Qrat\to K_3(\Frac A_2)\Qrat$ has nonzero kernel.
\end{proposition}
\begin{proof}
Put $v=z-1$. The defining equation is $2+2v+v^2+xy$, whose linear term in the regular ambient local ring is $2$. Thus $A_2$ is regular of dimension three, with maximal ideal generated by $x,y,v$ and residue field $\F_2$.

Take $W_2=\Z_{(2)}[i]$, $E_2=\Q(i)$, and $\varpi=1-i$. The equation $\varpi^2-2\varpi+2=0$ is Eisenstein, so $W_2$ is a DVR. The divisor $x=0$ in $\Spec A_2$ is regular, with coordinate ring $W_2[y]_{(\varpi,y)}$. Choose an integral anti-invariant Borel class $\beta\in K_3(W_2)$ exactly as in Lemma~\ref{d3:lem:beta}, and push it forward from this divisor.

After the finite free quadratic base change, put $w=z-i$. The local ring becomes
\[
\left(W_2[x,y,w]/(xy+w(w+2i))\right)_{(\varpi,x,y,w)}.
\]
Here $d=2i=-\varpi^2$ is nonzero and has positive valuation. Lemma~\ref{d3:lem:smallresolution} proves that $(x,w)$ has infinite order in the class group. The quotient by $\varpi$ is the domain
\[
\F_2[x,y,w]_{(x,y,w)}/(xy+w^2),
\]
so inverting $\varpi$ preserves this class group.

The generic fiber over $E_2$ is the split smooth quadric $xy+z^2+1=0$. Its Picard group, geometric Galois action, and Betti cohomology are exactly those used in Section~\ref{d3:sec:persistence}. The argument of Lemma~\ref{d3:lem:components} therefore proves persistence under every denominator from the displayed localization defining $A_2$. Flat base change along $A_2\to A_2\otimes_{\Z_{(2)}}W_2$, followed by inversion of $\varpi$, gives twice the product of $\beta$ with the ruling class, and the analytic regulator in twist three proves that this product has infinite order. Its generic vanishing follows from its original support on $x=0$, as before.
\end{proof}

\subsection{Every higher odd $K$-degree}

The regulator argument also has the following extension. It is stated separately to keep the degree-three counterexample independent of extra indexing conventions.

\begin{proposition}\label{d3:prop:odddegrees}
Let $p$ be an odd prime and let $n=2m-1\geq3$. Set $\varepsilon=(-1)^{m+1}$ and
\[
A_{p,m}=\left(\Z_{(p)}[x,y,z]/(xy+z^2-\varepsilon p)\right)_{(x,y,z)}.
\]
There is an integral element of infinite order in
\[
\ker\bigl(K_n(A_{p,m})\longrightarrow K_n(\Frac A_{p,m})\bigr).
\]
Thus the family with $xy+z^2+p$ works in degrees $3,7,11,\ldots$, and the family with $xy+z^2-p$ works in degrees $5,9,13,\ldots$.
\end{proposition}
\begin{proof}
Take $W=\Z_{(p)}[u]/(u^2-\varepsilon p)$, $E=\Q(u)$, and $\sigma(u)=-u$. The regularity and small-resolution arguments above are unchanged: after base change the equation is $xy+z^2-u^2$, with $d=2u$ in Lemma~\ref{d3:lem:smallresolution}.

If $m$ is even, $E$ is imaginary quadratic. Borel's regulator identifies $K_{2m-1}(E)\otimes\R$ with a one-dimensional real space, and the involution acts by $-1$: its target $\C/\R(m)$ is the imaginary quotient. If $m$ is odd, $E$ is real quadratic. The regulator target is a pair of real lines indexed by the two real embeddings, and $\sigma$ exchanges the two coordinates. Its anti-invariant part is again one-dimensional. These regulator ranks and conjugation conventions are those of \cite[Proposition~12.2]{Borel} and \cite[Section~6]{BunkeTamme}.

In either case choose $b\in K_n(E)$ with nonzero anti-invariant regulator. Since $n-1>0$ is even, $K_{n-1}(\F_p)=0$ and localization lifts $b$ to $K_n(W)$. Antisymmetrizing gives an integral $\beta$ with $\sigma\beta=-\beta$ and
\[
r_m(\beta)\neq0\quad\text{in}\quad\HD^1(\Spec\C,\R(m)).
\]
Push $\beta$ forward from the regular divisor $x=0$ of $\Spec A_{p,m}$. The resulting class is zero in the fraction field. On the split quadric localization it becomes $2\beta\eta$.

The relevant regulator component now has target
\[
\HD^{2(m+1)-n}(Q_s,\R(m+1))
=\HD^3(Q_s,\R(m+1)).
\]
Since $m+1\geq3>\dim_{\C}Q_s$, the same holomorphic Poincar\'e-lemma calculation gives
\[
\HD^3(Q_s,\R(m+1))
=H^2(Q_s,\C)/H^2(Q_s,\R(m+1)).
\]
The proof of Lemma~\ref{d3:lem:Betti}, applied to each geometric boundary curve, therefore gives the same injection. Multiplication by $c_1(D_+)=2\pi i\,h$ identifies the nonzero coefficient in $\C/\R(m)$ with a nonzero coefficient in $\C/\R(m+1)$. No other source regulator component contributes: for a complex point and a positive twist, analytic Deligne cohomology is concentrated in degree one. The target regulator is nonzero at every finite stage, and $K$-theoretic continuity proves infinite order as before.
\end{proof}

\begin{remark}
\label{d3:rem:fixedgaussianallodd}
The same argument gives a \emph{single} arithmetic local
ring with a nonzero rational generic kernel in every odd degree
\(2n-1\ge3\). Fix an odd prime \(p\), a prime \(\mathfrak v\)
of \(\Z_{(p)}[i]\) above \(p\), and put
\[
 F=\Q(i),\quad V=(\Z_{(p)}[i])_{\mathfrak v},\quad
 R=\left(V[x,y,z]/(p+xy+z^2)\right)_{(x,y,z)}.
\]
Here \(V\) is an unramified DVR with uniformizer \(p\), so \(R\)
is regular of dimension three. Set \(E=F(u)\), \(u^2=-p\),
\(W=V[u]\), and let \(\sigma\) fix \(F\) and send \(u\) to \(-u\).
Eisenstein's criterion makes \(W\) a DVR and gives \([E:F]=2\).
Both \(F\) and \(E\) are totally imaginary, with respectively
one and two complex places. Borel's theorem \cite[Proposition~12.2]{Borel} gives
\(\dim_{\Q}K_{2n-1}(F)_{\Q}=1\) and
\(\dim_{\Q}K_{2n-1}(E)_{\Q}=2\), for every \(n\ge2\).
Restriction and transfer satisfy
\(\mathrm{res}\,\mathrm{tr}=1+\sigma\) and
\(\mathrm{tr}\,\mathrm{res}=2\).
Thus the \(\sigma\)-fixed subspace is exactly the restriction
of \(K_{2n-1}(F)_{\Q}\), and the anti-invariant subspace has
dimension one.

Choose a nonzero class there, clear denominators, lift it to
\(K_{2n-1}(W)\), and antisymmetrize. The lift exists because
\(K_{2n-2}(k_V)=0\)
\cite[Chapter~IV, Corollary~1.13]{Weibel}. This produces an actual integral
\(\beta_n\) with \(\sigma\beta_n=-\beta_n\) and a nonzero
regulator at some complex embedding. Its pushforward from
\(R/(x)=W[y]_{(u,y)}\) is generically zero and becomes
\(2\beta_n(1-[\OO(-D_+)])\) on the split generic quadric.
The small resolution and the persistence proved in
Lemma~\ref{d3:lem:Betti} are unchanged. In twist \(n+1>2\), the detecting target is
\[
 \HD^3(Q_s,\R(n+1))
   =H^2(Q_s,\C)/H^2(Q_s,\R(n+1)).
\]
Cup with the ruling multiplies the point coefficient
\(\C/\R(n)\) by \(2\pi i\), giving an isomorphism onto the
ruling line modulo \(\R(n+1)\). The product is nonzero at
every finite stage, hence has infinite order after localization.
Here \(\sigma\) is \emph{not} complex conjugation: it fixes \(i\).
The invariant-rank argument, and multiplication by \(2\pi i\),
work for both parities of \(n\) \cite{Borel,BunkeTamme,Quillen}.
\end{remark}


\subsection{Eisenstein polynomials and arbitrary kernel rank}\label{d3:sec:eisenstein}

The quadratic construction is one instance of a root-lattice calculation. This section records the stronger statement and its geometric proof.

\begin{theorem}\label{d3:thm:eisenstein}
Let $p$ be any prime, let $f\in\Z_{(p)}[z]$ be a monic Eisenstein polynomial of degree $N\geq2$, and put
\[
A_f=\left(\Z_{(p)}[x,y,z]/(xy+f(z))\right)_{(x,y,z)},\qquad
W=\Z_{(p)}[\theta]/(f(\theta)),\qquad K=\Q(\theta).
\]
The ring $A_f$ is regular local of dimension three and has fraction field $\Q(x,z)$. Let $(r_1,r_2)$ be the signature of $K$. For $n\geq3$ odd, consider
\[
g_n:K_n(W)_{\Q}\longrightarrow K_n(A_f)_{\Q},\qquad
\beta\longmapsto i_*(\beta|_{A_f/(x)}),
\]
where $i$ is the regular principal divisor $x=0$. Its image is contained in the generic kernel. Moreover,
\[
\begin{array}{c|c|c}
n&\ker g_n&\dim_{\Q}\operatorname{im}g_n\\ \hline
3\pmod4&0&r_2\\
1\pmod4&\operatorname{res}_{K/\Q}K_n(\Q)_{\Q}&r_1+r_2-1.
\end{array}
\]
In the second row $n\geq5$, and $K_n(W)_{\Q}$ is identified with $K_n(K)_{\Q}$ by localization.
\end{theorem}

\subsubsection{The divisor lattice after splitting}

Let $O$ be a DVR with uniformizer $\pi$, residue field $k$, and fraction field $L$. Suppose that $a_1,\ldots,a_N$ are distinct elements of $\pi O$. Define
\[
B=\left(O[x,y,z]/\left(xy+\prod_{i=1}^N(z-a_i)\right)\right)_{(\pi,x,y,z)},
\qquad P_i=(x,z-a_i).
\]
The same normality argument as in Lemma~\ref{d3:lem:smallresolution} shows that $B$ is normal and regular off its closed point. Since $B[1/x]$ is a localization of $O[x,x^{-1},z]$, the primes $P_i$ generate $\Cl(B)$, with the relation
\[
\divisor(x)=P_1+\cdots+P_N.
\]
To prove that there are no further relations, successively blow up the strict transforms of $P_1,\ldots,P_{N-1}$. The resulting regular scheme $Y$ has $N$ charts, each a localization of $O[u_i,v_i]$, on which
\begin{align*}
z&=a_i+u_i v_i,\\
x&=u_i\prod_{k<i}(z-a_k),&
y&=-v_i\prod_{k>i}(z-a_k).
\end{align*}
For completeness, before the $i$th blowup the unresolved chart is
\[
O[s_{i-1},y,z]/\left(s_{i-1}y+\prod_{k\geq i}(z-a_k)\right),
\quad s_0=x.
\]
Blowing up $(s_{i-1},z-a_i)$ gives the displayed $i$th polynomial chart and the next unresolved chart, where $s_i=s_{i-1}/(z-a_i)$. These are the actual Rees charts inside the common function field. A future strict transform $P_j$, $j>i$, is absent from the resolved chart, since its generic point has $z=a_j$ and cannot satisfy $u_i=0$, $z=a_i$. Thus the successive centers glue globally and do not modify previous resolved charts.

The adjacent transition is
\[
u_{i+1}=v_i^{-1},\qquad
v_{i+1}=v_i(a_i-a_{i+1}+u_i v_i).
\]
The exceptional fiber is a reduced chain $C_1,\ldots,C_{N-1}$ of projective lines. Its ideal is $(\pi,u_1)$ on the first chart, $(\pi,v_N)$ on the last, and $(\pi,u_i v_i)$ on an interior chart. In particular it has dimension one, and the resolution has no exceptional divisors. The two charts on $C_j$ have coordinates $v_j,u_{j+1}$ with $u_{j+1}=v_j^{-1}$.

Write $\widetilde P_\ell$ for the strict transform of $P_\ell$. On the $i$th chart its defining equation is
\[
d_{\ell i}=\begin{cases}
z-a_\ell,&\ell<i,\\
u_i,&\ell=i,\\
1,&\ell>i.
\end{cases}
\]
These equations also show that no vertical divisor has been added: when $\ell<i$, the quotient by $u_i v_i+a_i-a_\ell$ is a domain. Taking transition functions of $\OO_Y(\widetilde P_\ell)$ gives
\[
\deg\bigl(\OO_Y(\widetilde P_\ell)|_{C_j}\bigr)
=\delta_{\ell,j+1}-\delta_{\ell,j}.
\]
Therefore a principal relation $\sum n_iP_i$ forces $n_{j+1}=n_j$ for every $j$. Conversely equal coefficients give a power of $x$. We have proved the integral presentation
\begin{equation}\label{d3:eq:rootlattice}
\Cl(B)=\Z^N/\Z(1,\ldots,1).
\end{equation}

Since $B/\pi B$ is the domain $k[x,y,z]_{(x,y,z)}/(xy+z^N)$, inverting $\pi$ removes just the principal prime $(\pi)$. Thus \eqref{d3:eq:rootlattice} is also $\Pic(B[1/\pi])$. On the full generic surface
\[
Q_L=\Spec L[x,y,z]/\left(xy+\prod_i(z-a_i)\right),
\]
divisor localization at $x$ gives the same presentation: the units of $L[x,x^{-1},z]$ are $L^\times x^{\Z}$. Consequently
\begin{equation}\label{d3:eq:generalpiciso}
\Pic(Q_L)\xrightarrow{\sim}\Pic(B[1/\pi]).
\end{equation}

\subsubsection{Betti classes and finite localizations}

Assume that $L$ embeds in $\C$. The complex surface $Q_L(\C)$ is the blowup of $\A^2_{\C}$ at the $N$ points $(0,a_i)$ with the strict transform of the line $x=0$ removed. Indeed, blow up the ideal $(x,\prod(z-a_i))$; its $x$-chart is precisely $Q_L$, and the complement is that strict transform.

The blowup has $H^2=\Z^N$ on the exceptional classes $E_i$, and no other positive-degree cohomology. The removed line is isomorphic to $\A^1$, with class $-\sum E_i$, since its sum with the exceptional divisors is the principal divisor of $x$. The Gysin sequence therefore gives
\[
H^2(Q_L(\C),\Z)=\Z^N/\Z(1,\ldots,1),
\]
and first Chern class identifies this with $\Pic(Q_L)$. All the root divisors are defined over $L$, so the absolute Galois action on the geometric Picard lattice is trivial.

For any denominator allowed in the localization $B[1/\pi]$, every removed $L$-prime divisor has zero Picard class by \eqref{d3:eq:generalpiciso}. Its geometric components are a single Galois orbit; they have equal classes in a torsion-free lattice, and their sum is zero. Each component therefore has zero class. The degree-two support argument of Section~\ref{d3:sec:persistence} proves
\[
H^2(Q_L(\C),\R)\hookrightarrow H^2((Q_L)_s(\C),\R)
\]
for every such finite denominator $s$. The analytic Deligne injection in every twist greater than two follows exactly as in Section~\ref{d3:sec:Delignedetection}.

\subsubsection{The arithmetic rank calculation}

We now prove Theorem~\ref{d3:thm:eisenstein}. Eisenstein's condition makes $W$ a DVR with uniformizer $\theta$ and residue field $\F_p$. It also makes the defining equation of $A_f$ have nonzero linear term $p$ in its regular four-dimensional ambient local ring. Thus $A_f$ is regular local of dimension three, and
\[
A_f/(x)=W[y]_{(\theta,y)}.
\]
Solving for $y$ after inverting $x$ gives $\Frac A_f=\Q(x,z)$. Generic vanishing of $g_n$ follows from localization.

Choose a splitting field $L$ of $f$ and a prime above $p$, and let $O$ be the corresponding coefficient DVR. Every root $a_i$ has positive valuation. The map $A_f\to B$ used above is flat: $O$ is torsion-free over $\Z_{(p)}$, followed by localization. There is no need for $O$ itself to be a finite $\Z_{(p)}$-module. After inverting $\pi$, flat base change expresses the Gysin image as
\begin{equation}\label{d3:eq:generalbasechange}
g_n(\beta)|_{B[1/\pi]}=
\sum_{i=1}^N \sigma_i(\beta)\,\eta_i,
\qquad \eta_i=[\OO_{(x,z-a_i)}],
\end{equation}
where $\sigma_i:K\hookrightarrow L$ are the embeddings indexed by the roots.

Write $n=2m-1$. At a fixed complex embedding of $L$, the regulator in weight $m+1$ sends \eqref{d3:eq:generalbasechange} to the vector of Borel coefficients
\[
\bigl(r_m(\sigma_1\beta),\ldots,r_m(\sigma_N\beta)\bigr)
\quad\text{modulo the diagonal line},
\]
with the common nonzero Tate normalization coming from first Chern class. Its target is
\[
\HD^3((Q_L)_s,\R(m+1))
=H^2((Q_L)_s,\C)/H^2((Q_L)_s,\R(m+1)),
\]
because $m+1\geq3$. The preceding finite-localization injection ensures that a nonzero vector modulo the diagonal remains nonzero at every finite stage, hence in $K_n(B[1/\pi])_{\Q}$.

If $m$ is even, the two coefficients at conjugate embeddings have opposite signs, and a real embedding contributes zero. Borel's theorem \cite[Proposition~12.2]{Borel} identifies the real scalar extension of $K_n(K)_{\Q}$ with this $r_2$-dimensional coefficient space. A vector in it is diagonal only when it is zero. Thus $g_n$ is injective for $n\equiv3\pmod4$.

If $m$ is odd, conjugate coefficients agree, and the coefficient space has dimension $r_1+r_2$. Its diagonal line is exactly the real span of restriction from $K_n(\Q)_{\Q}$, which is one-dimensional. Since this line is generated by a rational class and Borel's real regulator is injective \cite[Section~6]{BunkeTamme}, its intersection with $K_n(K)_{\Q}$ is exactly that rational restriction line. Finally, that line really is killed by $g_n$: its classes lift from $K_n(\Z_{(p)})_{\Q}$, and the projection formula multiplies them by
\[
[\OO_{V(x)}]=[\OO]-[\OO]=0\quad\text{in }K_0(A_f).
\]
This proves both the exact kernel and the stated image dimensions.

\begin{corollary}
For every prime $p$ and integer $r\geq1$, the three-dimensional regular local ring
\[
\left(\Z_{(p)}[x,y,z]/(xy+z^{2r}+p)\right)_{(x,y,z)}
\]
has a rational degree-three generic kernel of dimension at least $r$. Thus these kernel dimensions are unbounded in fixed local dimension three. Moreover the single ring with $xy+z^3+p$ has a nonzero rational generic kernel in every odd degree $n\geq3$: its number field has signature $(1,1)$.
\end{corollary}

\section{Explicit triples of units at every residue prime}\label{d3:sec:explicitallprimes}

The preceding geometry also detects explicit symbols. The following formulas use the Dennis--Stein convention of Section~\ref{d3:sec:explicit}; unlike \eqref{d3:eq:explicitfront}, they require no prefactor.

\begin{theorem}\label{d3:thm:explicitallprimes}
For an odd prime $p$, put
\[
 A_p=\left(\Z_{(p)}[x,y,z]/(xy+\Phi_{2p}(z-1))\right)_{(x,y,z)}.
\]
Then
\begin{equation}\label{d3:eq:explicitoddprime}
 c_p=\left\{2-z,\ 1-x,\
 1+(1-x)yz\bigl((z-1)^p-1\bigr)\right\}\in K_3(A_p)
\end{equation}
has infinite order and maps to zero in $K_3(A_p[1/x])$ integrally.
For $p=2$, put
\[
 A_2=\left(\Z_{(2)}[x,y,t]/(xy+t^2+1)\right)_{(2,x,y,t-1)},
 \qquad H(t)=t^{10}-t^8+t^6-t^4+t^2-1.
\]
The same assertions hold for
\begin{equation}\label{d3:eq:explicittwoprime}
 c_2=\left\{1+t+t^2,\ 1-x,\ 1+(1-x)yH(t)\right\}\in K_3(A_2).
\end{equation}
All entries in these formulas are units. Each $A_p$ is regular local of dimension three, with residue field $\F_p$, so the displayed classes give nonzero rational generic kernels.
\end{theorem}

\subsection{The integral identities}

We first record the common calculation. Suppose $1-xb=u^n$, and put
\[
 \delta=\langle x,b\rangle,\qquad \lambda=1-x,\qquad
 \mu=1+\lambda b.
\]
Assume $u,g,\lambda,\mu$ are units and $\{g,u^n\}=0$ in $K_2$.
Relation \eqref{d3:eq:DS2}, together with $\langle x,1\rangle=0$, gives
\[
 \delta=\langle x,b+1-xb\rangle=\langle x,\mu\rangle
       =\{\lambda u^n,\mu\},
 \qquad 1-x\mu=\lambda u^n.
\]
Consequently
\begin{equation}\label{d3:eq:allprimesDS}
 [g]\delta=\{g,\lambda,\mu\},\qquad
 [g]\delta|_{D(x)}=\{g,x,u^n\}=-[x]\{g,u^n\}=0.
\end{equation}
These are integral identities in Quillen $K$-theory.

For odd $p$, write $a=z-1$ and set
\[
 H(a)=(a+1)(a^p-1),\qquad b=yH(a),\qquad g=1-a,\qquad n=2p.
\]
The polynomial identity
\[
 a^{2p}-1=H(a)\Phi_{2p}(a)=-xyH(a)
\]
gives $1-xb=a^{2p}$. At the closed point,
$a=-1$, $g=2$, and $\lambda=\mu=1$ in $\F_p$; these elements are units.
Steinberg gives $\{1-a,a^{2p}\}=2p\{1-a,a\}=0$.
Thus \eqref{d3:eq:allprimesDS} proves \eqref{d3:eq:explicitoddprime} and its integral vanishing after inverting $x$.
Moreover,
\[
 \Phi_{2p}(z-1)
 =\sum_{k=0}^{p-1}(1-z)^k
 =\sum_{j=1}^{p}(-1)^{j+1}\binom pj z^{j-1}
\]
is monic Eisenstein of degree $p-1$, with constant term $p$. The regularity assertion follows as in Theorem~\ref{d3:thm:eisenstein}.

At two, take $u=t$, $n=12$, $b=yH(t)$ and $g=1+t+t^2$.
Here
\[
 (t^2+1)H(t)=t^{12}-1,\qquad 1-xb=t^{12}.
\]
All required units have nonzero residue at $t=1$. To prove
$\{g,t^{12}\}=0$ integrally, use the finite \'etale local extension
$A_2'=A_2[\zeta]$, where $\zeta^2+\zeta+1=0$.
The factor $1-\zeta t$ is a unit, since
\[
 (1-\zeta t)(1-\zeta^2t)=g,
 \qquad (\zeta t)^{12}=t^{12}.
\]
Steinberg and the transfer projection formula therefore give
\[
 0=\operatorname{Tr}_{A_2'/A_2}
       \{1-\zeta t,t^{12}\}
   =\{\operatorname{Norm}_{A_2'/A_2}(1-\zeta t),t^{12}\}
   =\{g,t^{12}\}.
\]
There is no factor of two in this formula. Indeed, multiplication by
$1-\zeta t$ in the basis $1,\zeta$ has matrix
$\left(\begin{smallmatrix}1&t\\-t&1+t\end{smallmatrix}\right)$,
elementary-equivalent to $\operatorname{diag}(1,g)$; its $K_1$ transfer is exactly $[g]$.
See \cite[Chapter~III, Definition~5.11 and relations (D1)--(D3)]{Weibel}.
Equation \eqref{d3:eq:allprimesDS} now proves \eqref{d3:eq:explicittwoprime}.
Finally, with $v=t-1$, the equation becomes
$xy+v^2+2v+2=0$, which proves regularity at the specified maximal ideal.

\subsection{A supported lift on a finite surface}

For odd $p$ let $E=\Q(\zeta_{2p})$, and at two let $E=\Q(i)$.
Use the following smooth split surfaces and finite opens:
\[
\begin{array}{c|c|c}
 & Q & U\\ \hline
p\text{ odd}
 & \Spec E[x,y,a]/(xy+\Phi_{2p}(a))
 & D(a(1-a))\\
p=2
 & \Spec E[x,y,t]/(xy+t^2+1)
 & D(t(1+t+t^2)).
\end{array}
\]
The Dennis--Stein product $\alpha_p=[g]\delta$ is defined on $U$
and vanishes on $U[1/x]$ by the preceding calculation.
At two the same norm argument uses $E(\zeta)/E$.
The entire divisor $D_x=V(x)\subset Q$ lies in $U$:
its coordinates are primitive $2p$-th roots in the odd case, and
$t=\pm i$, $g(\pm i)=\pm i$, at two.
Localization and d\'evissage supply $\xi_p\in K_3(D_x)$ such that its
pushforward to $U$ is $\alpha_p$. Pushing the same class into $Q$ gives
\begin{equation}\label{d3:eq:allprimeslift}
 \widetilde c_p=i_*\xi_p\in K_3(Q),
 \qquad \widetilde c_p|_U=\alpha_p.
\end{equation}
Only this supported lift will be needed; its coefficients need not be specified.

We calculate the regulator of $\alpha_p$ using its Dennis--Stein expression.
The entries $\lambda,\mu$ need only be units in the arithmetic local ring;
they are not needed on the analytic cycle below.
On $D(x)\subset U$ one has $\delta=\{x,u^n\}$, and its holomorphic
weight-two curvature extends to $U$ as
\[
 \omega=n\,\frac{dx}{x}\wedge\frac{du}{u}
        =-\frac{dx\wedge db}{1-xb}.
\]
Whenever $g=e^L$ on an analytic neighborhood, the interval product
calculation \eqref{d3:eq:Lomega} gives
$r_3(\alpha_p)=[-L\omega]$ there.
This uses the multiplicativity and unit normalization of
\cite[Proposition~2.8, Theorem~2.31 and Lemma~4.5]{BunkeTamme}.

\subsection{Two nonzero matching-sphere periods}

We use a single endpoint formula. Let $\gamma:[-1,1]\to\C$ be a smooth
embedded arc joining two distinct simple roots of the relevant polynomial
$P(u)$, avoiding the other roots. On $xy+P(u)=0$, take
\[
 u=\gamma(\tau),\qquad
 x=\sqrt{|P(\gamma(\tau))|}\,e^{i\vartheta},\qquad
 y=-P(\gamma(\tau))/x.
\]
Adding $x=y=0$ at both ends gives a smooth sphere $\Sigma_\gamma$,
oriented by $d\tau\wedge d\vartheta$. Smoothness at an end follows by
using $|x|^2$ as a radial coordinate and writing $y$ as a smooth
nonzero phase times $\overline x$.

Suppose $L$ and $F$ are holomorphic near the arc, with
$e^L=g$ and $dF=-L\,du/u$. On the sphere away from its ends,
\[
 L\omega=n\,d\left(F\,\frac{dx}{x}\right).
\]
Stokes' theorem on a truncated cylinder, followed by shrinking its
two missing caps, gives
\begin{equation}\label{d3:eq:cyclotomicendpoint}
 \int_{\Sigma_\gamma}r_3(\alpha_p)
 =-2\pi i n\bigl(F(\gamma(1))-F(\gamma(-1))\bigr)
 \pmod{\R(3)}.
\end{equation}
The caps contribute zero because $L\omega$ extends smoothly.
This calculation uses $dx/x$ on the cylinder and never requires a
logarithm of $x$. Since $\R(3)=i\R$, a nonzero real period detects
the regulator.

For odd $p$, put $\theta=\pi/p$, $r=e^{i\theta}$, and choose
\[
 \gamma(\tau)=\cos\theta+i\sin\theta\,\tau .
\]
This chord joins $\overline r$ to $r$, lies in $0<\Re a<1$, and its
interior lies strictly inside the unit circle. It therefore avoids
$0,1$ and all other roots. On one neighborhood of the whole arc,
take the principal branches
\[
 L=\operatorname{Log}(1-a),\qquad F=\operatorname{Li}_2(a).
\]
Here $\operatorname{Li}_2$ is the classical dilogarithm, initially
given by $\sum_{r\geq1}a^r/r^2$ on the unit disc. The chosen
branches satisfy the required differential equation.
Writing $D$ for the Bloch--Wigner dilogarithm,
$F(r)-F(\overline r)=2iD(r)$, so \eqref{d3:eq:cyclotomicendpoint} yields
\begin{equation}\label{d3:eq:oddcyclotomicperiod}
 \int_{\Sigma_\gamma}\Re r_3(\alpha_p)
 =8\pi p\,D(e^{i\pi/p})>0.
\end{equation}
For example, positivity follows from
$D(e^{i\theta})=-\int_0^\theta\log(2\sin(v/2))\,dv$
for $0<\theta\leq\pi/3$.

At two, the straight segment from $-i$ to $i$ would meet $t=0$.
Instead use
\[
 \gamma(\tau)=\tfrac12(1-\tau^2)+i\tau .
\]
Its interior has positive real part and modulus less than one.
Thus the arc avoids $0$, the roots of $g$, and the other twelfth
roots of unity. With $\zeta=e^{2\pi i/3}$, both
$1-\zeta\gamma(\tau)$ and $1-\zeta^2\gamma(\tau)$ have positive
real part, including at the endpoints. Hence one neighborhood
supports the principal branches
\[
 L=\operatorname{Log}(1-\zeta t)+\operatorname{Log}(1-\zeta^2t),
 \qquad
 F=\operatorname{Li}_2(\zeta t)+\operatorname{Li}_2(\zeta^2t).
\]
Here $e^L=g$, $dF=-L\,dt/t$, and
$F(-i)=\overline{F(i)}$.
The absolutely convergent dilogarithm series on the unit circle gives
\[
 \operatorname{Li}_2(w)+\operatorname{Li}_2(\zeta w)
 +\operatorname{Li}_2(\zeta^2w)=\tfrac13\operatorname{Li}_2(w^3).
\]
Taking imaginary parts at $w=i$ therefore gives
$D(\zeta i)+D(\zeta^2i)=-\tfrac43D(i)$, and
\begin{equation}\label{d3:eq:twocyclotomicperiod}
 \int_{\Sigma_\gamma}\Re r_3(\alpha_2)
 =-64\pi D(i)\ne0.
\end{equation}
The number $D(i)=\sum_{k\geq0}(-1)^k/(2k+1)^2$ is positive.

\subsection{Persistence in the arithmetic local ring}

We finish the proof of Theorem~\ref{d3:thm:explicitallprimes}.
For odd $p$, take the coefficient DVR
$\OO=\Z_{(p)}[\zeta_{2p}]$ and use the original coordinate $z=a+1$.
All roots $1+r_i$ of $\Phi_{2p}(z-1)$ lie in its maximal ideal.
At two take $\OO=\Z_{(2)}[i]$ and the coordinate $v=t-1$; its
two roots $\pm i-1$ again lie in the maximal ideal.
The split local rings therefore satisfy \eqref{d3:eq:rootlattice}
and \eqref{d3:eq:generalpiciso}. The Betti and analytic Deligne
persistence argument of Section~\ref{d3:sec:eisenstein} applies to
every permitted finite localization of $Q$.

The open $U$ is one of those permitted localizations:
$a(1-a)$ has nonzero residue $-2$ for odd $p$, and $tg$ has
nonzero residue $1$ at two. On the generic fiber $T$ of the split
arithmetic local ring, the unit triple $c_p$, the Dennis--Stein
product $\alpha_p$, and the lift $\widetilde c_p$ in
\eqref{d3:eq:allprimeslift} have the same image.
Equations \eqref{d3:eq:oddcyclotomicperiod} and
\eqref{d3:eq:twocyclotomicperiod} show that the regulator of
$\widetilde c_p$ is nonzero already on $Q$.
The cited persistence makes it nonzero on every further permitted
finite open. A torsion relation on $c_p$ would, after passage to
$T$, descend to one such finite open by \eqref{d3:eq:continuity},
contradicting that real regulator.
Thus $c_p$ has infinite order, as asserted.


\section{Completion and arithmetic \texorpdfstring{$p$}{p}-adic detection}\label{d3:sec:completed}

There is a separate route to completed local rings. It uses the arithmetic comparison theorem for dagger varieties of Colmez--Nizio\l, rather than an interchange of a localization limit with a $p$-adic inverse limit. We spell out the map and the order of limits, and distinguish actual $K_3$-classes from classes in a completed spectrum.

\subsection{A ruling-cup injection on the completed node}

\begin{lemma}\label{d3:lem:completedcup}
Let $E/\Q_p$ be finite, let $W=\OO_E$, and let $0\ne d\in\mathfrak m_W$. Put
\[
\widehat B_d=W[[x,y,w]]/(xy+w(w+d)),\qquad
T_d=\Spec\widehat B_d[1/d].
\]
The ideal $I=(x,w)$ is invertible on $T_d$, and cup product induces an injection
\[
H^1(E,\Q_p(2))\longrightarrow
H^3_{\mathrm{cont},\mathrm{\acute et}}(T_d,\Q_p(3)),
\qquad b\longmapsto b\smile c_1(I).
\]
Here continuous cohomology is formed by taking the derived inverse limit of finite-coefficient cohomology and then inverting $p$.
\end{lemma}
\begin{proof}
At a prime containing $x,w$, the element $w+d$ is a unit, and the equation gives $I=(x)$. Outside $V(I)$ the ideal is the unit ideal. It is therefore an invertible ideal, generated globally by two elements, and defines a morphism
\[
f_T:T_d\longrightarrow\PP^1_E,\qquad f_T^*\OO(1)=I.
\]

Consider the smooth dagger affinoid
\[
U^\dagger=\operatorname{Sp}^\dagger
E\langle X,Y,Z\rangle^\dagger/(XY+Z(Z+1)).
\]
The two roots $0,-1$ are distinct in every residue characteristic. Fix
\[
\rho=|d|^{-1/2}\in\sqrt{|E^\times|},\qquad 1<\rho<|d|^{-1},
\]
and let $U_\rho$ be the closed radius-$\rho$ domain of the same algebraic quadric. The substitutions
\begin{equation}\label{d3:eq:fixedradius}
x=dX,\qquad y=dY,\qquad w=dZ
\end{equation}
define a homomorphism $\widehat B_d[1/d]\to\OO(U_\rho)$. Indeed every formal series with coefficients in $W$ converges on this \emph{same fixed domain}: terms of total degree $n$ have weighted norm at most $(|d|\rho)^n$, which tends to zero. A fixed power of $d^{-1}$ does not change the convergence radius.

For completeness, a noetherian integral model of this fixed affinoid is
\[
 A_\rho=W\langle a,b,c,r,s,t\rangle/
 (a^2-dr,\ b^2-ds,\ c^2-dt,\ ab+c(c+d)).
\]
Its generic fiber is $U_\rho$ under $X=a/d$, $Y=b/d$, $Z=c/d$:
the last three auxiliary variables impose
$|dX^2|,|dY^2|,|dZ^2|\leq1$.
The elements $a,b,c$ are topologically nilpotent in $A_\rho$,
since their squares lie in $dA_\rho$. Substitution therefore defines
$\widehat B_d\to A_\rho$ before inverting $d$.
The pair $(A_\rho,(p))$ is a noetherian henselian pair; removing its
$W$-torsion, if desired, preserves the generic fiber.

Pull scheme cohomology to $\Spec\OO(U_\rho)$, then to the affinoid \'etale topos. The latter map exists and is compatible with all finite coefficients, including $p$-power coefficients; see the Gabber--Fujiwara comparison in \cite[Theorem~6.2.1]{Achinger}, applied to a complete noetherian integral model of this one affinoid. At this fixed $\rho$, take the derived inverse limit over $p^\nu$ and invert $p$. Continuous \'etale and pro-\'etale cohomology agree for the affinoid, as recalled in \cite[Section~3.2.3]{ColmezNiziol}. Finally map into the homotopy colimit over shrinking strict neighborhoods, which is the dagger pro-\'etale complex in that reference. This constructs
\begin{equation}\label{d3:eq:schemetodagger}
H^3_{\mathrm{cont},\mathrm{\acute et}}(T_d,\Q_p(3))
\longrightarrow H^3_{\mathrm{pro\acute et}}(U^\dagger,\Q_p(3)).
\end{equation}
It respects pullback from $\PP^1_E$ and first Chern classes. We have not exchanged the inverse limit over coefficients with a direct limit over opens. The reverse comparison from dagger \'etale to dagger pro-\'etale cohomology is not assumed.

The ruling morphism $f:U^\dagger\to\PP^1_E$ obtained from $(X,Z)$ has an explicit two-chart description. Above the two standard closed discs of $\PP^1$ its charts are dagger bidiscs
\begin{align*}
U_0: &(X,t),&Z&=Xt,&Y&=-t(1+Xt),\\
U_\infty: &(s,Y),&Z&=-1-sY,&X&=-s(1+sY).
\end{align*}
Their intersection is a dagger unit annulus times a disc, with
\[
s=t^{-1},\qquad Y=-t-Xt^2.
\]
The bidiscs have no positive-degree de Rham cohomology, and the intersection has $H^1=E\,d\log t$. These facts follow by coefficientwise integration of overconvergent power and Laurent series: division by integers costs only a slight reduction of the overconvergence radius, and the $t^{-1}dt$ term remains. Coherent acyclicity for dagger affinoids \cite[Proposition~3.1]{GrosseKlonne} and Mayer--Vietoris give
\[
H^2_{\mathrm{dR}}(U^\dagger)=E[d\log t],\qquad H^3_{\mathrm{dR}}(U^\dagger)=0,
f^*:H^2_{\mathrm{dR}}(\PP^1_E)\xrightarrow{\sim}H^2_{\mathrm{dR}}(U^\dagger).
\]

The arithmetic Theorem~1.3(2), or Theorem~7.9(2) together with Proposition~5.9(2), of \cite{ColmezNiziol} gives a natural boundary
\[
\partial_3:\widetilde H^2(R\Gamma_{\mathrm{dR}}(U^\dagger)/F^3)
\longrightarrow\widetilde H^3_{\mathrm{pro\acute et}}(U^\dagger,\Q_p(3)).
\]
We verify that it is an isomorphism, so no passage from a monomorphism in a topological heart to ordinary cohomology is needed. The space $U^\dagger$ has the smooth weak formal model
$\OO_E\langle X,Y,Z\rangle^\dagger/(XY+Z(Z+1))$; the two factors $Z,Z+1$ generate the unit ideal. The ideal $(X,Z)$ is invertible on this model and extends the ruling morphism to the weak formal projective line over $\OO_E$. Let $F_0=\Frac W(k_E)$. The functorial Hyodo--Kato to de Rham comparison of \cite[equation~(5.4) and Propositions~5.5--5.6]{ColmezNiziol} becomes a strict quasi-isomorphism after tensoring from $F_0$ to $E$. Our de Rham calculation, and its vanishing in degree three, therefore give
\[
f^*:H^2_{\mathrm{HK}}(\PP^1_E)\xrightarrow{\sim}H^2_{\mathrm{HK}}(U^\dagger),
\qquad H^3_{\mathrm{HK}}(U^\dagger)=0.
\]
Indeed scalar extension to the field $E$ is faithful. Both degree-two groups are the rank-one slope-one isocrystal $F_0h$ with $\varphi=p\sigma$ and $N=0$. Its $\varphi=p^2$ eigenspace is zero by valuation. The second exact sequence of Theorem~1.3(2) now makes the term after $\widetilde H^3_{\mathrm{pro\acute et}}$ vanish, on both $U^\dagger$ and $\PP^1_E$. Since $F^3=0$ on this surface, the first exact sequence gives isomorphisms
\begin{align*}
H^2_{\mathrm{dR}}(U^\dagger)&\xrightarrow{\sim}\widetilde H^3_{\mathrm{pro\acute et}}(U^\dagger,\Q_p(3)),\\
H^2_{\mathrm{dR}}(\PP^1_E)&\xrightarrow{\sim}\widetilde H^3_{\mathrm{pro\acute et}}(\PP^1_E,\Q_p(3)).
\end{align*}
Here the sources are classical. Applying the classical-part functor of \cite[Section~3.1.1]{ColmezNiziol} to these isomorphisms gives the corresponding isomorphisms on ordinary cohomology. Naturality and the de Rham ruling isomorphism show that
\[
f^*:H^3_{\mathrm{\acute et}}(\PP^1_E,\Q_p(3))
\longrightarrow H^3_{\mathrm{pro\acute et}}(U^\dagger,\Q_p(3))
\]
is an isomorphism. Proper pro-\'etale and analytic \'etale cohomology
are identified by \cite[Section~7.2.4]{ColmezNiziol}. For the scheme
$\PP^1_E$, the torsion algebraic--analytic comparison of
\cite[Theorem~6.1.1]{BerkovichEtale} gives the same identification,
first for every $p^\nu$ and then by derived inverse limit.
These comparisons preserve pullbacks and Kummer Chern classes.
For finite coefficients, put $\Lambda=\Z/m$, where $m>0$ is
invertible in $E$, and write
$\mathrm{pr}\colon\PP^1_E\to\Spec E$.
The Kummer boundary defines
$\xi_m=c_1(\OO(1))\in H^2_{\mathrm{\acute et}}(\PP^1_E,\Lambda(1))$
\cite[Example~7.9(a)]{MilneEtale}.
For every integer $r$, pullback and cup product give a morphism
of \'etale complexes
\begin{equation}\label{d3:eq:finitePonebundle}
 \Lambda(r)\oplus\Lambda(r-1)[-2]
 \xrightarrow{(\mathrm{pr}^*,\,\mathrm{pr}^*(-)\smile\xi_m)}
 R\mathrm{pr}_*\Lambda(r).
\end{equation}
To check that it is a quasi-isomorphism, take a geometric stalk.
Proper base change identifies the target with the cohomology
complex of a geometric projective line, whose only nonzero
cohomology is in degrees zero and two
\cite[Proposition~14.2 and Theorems~17.7 and~17.10]{MilneEtale}.
The degree-zero map is the identity.
In degree two, the Kummer calculation sends the generator
$\OO(1)$ of $\Pic(\PP^1)=\Z$ to a generator modulo $m$;
this is also the Chern normalization in
\cite[Theorems~23.2 and~23.3(b)]{MilneEtale}.
Thus the map is a quasi-isomorphism on geometric stalks and
hence on $\Spec E$. This argument works over any field in
which $m$ is invertible.

The Kummer classes make these maps compatible for $m=p^\nu$.
Assume $p\ne\operatorname{char}(E)$.
Applying $R\Gamma(E,-)$, then $R\varprojlim_\nu$, and then
inverting $p$ preserves this finite direct sum; all
$\Q_p$-cohomology groups here, including $H^1(E,\Q_p(2))$,
are understood as continuous cohomology obtained by this
construction. Taking $r=3$ gives a split injection
\[
 H^1(E,\Q_p(2))
 \xrightarrow{\ \smile c_1(\OO(1))\ }
 H^3_{\mathrm{cont},\mathrm{\acute et}}(\PP^1_E,\Q_p(3)).
\]
Composing the ruling-cup map on $T_d$ with
\eqref{d3:eq:schemetodagger} is this split injection followed
by the already proved isomorphism $f^*$.
Hence the ruling-cup map on $T_d$ is injective.
\end{proof}

\subsection{Simultaneous ruling classes}
A map from a completed quadratic node isolates any chosen pair of roots.
The resulting detector records their coefficients modulo the diagonal.

\begin{lemma}\label{d3:lem:completedrootcup}\label{d3:lem:completedtamecup}
Let $E/\Q_p$ be finite and let $a_1,\ldots,a_N$ be distinct elements of $\mathfrak m_{\OO_E}$, with $N\geq2$. Put
\[
\widehat B=\OO_E[[x,y,z]]/
\left(xy+\prod_{i=1}^N(z-a_i)\right),
\qquad T=\Spec\widehat B[1/p],
\qquad J_i=(x,z-a_i).
\]
The $J_i$ are invertible ideals on $T$, their product is $(x)$, and the map
\begin{equation}\label{d3:eq:completedtamecup}
\frac{H^1(E,\Q_p(2))^N}
{\operatorname{diag}H^1(E,\Q_p(2))}
\longrightarrow H^3_{\mathrm{cont},\mathrm{\acute et}}(T,\Q_p(3)),
\qquad
[(\alpha_i)]\longmapsto\sum_i\alpha_i\smile c_1(J_i)
\end{equation}
is injective.
\end{lemma}

\begin{proof}
At a point with $x=z-a_i=0$, every other root factor is a unit, since each nonzero $a_i-a_k$ is invertible after inverting $p$. The equation then gives $J_i=(x)$ locally. Away from its vanishing locus $J_i$ is the unit ideal. This proves invertibility and, by the same local calculation, $\prod_iJ_i=(x)$.

Fix distinct indices $i,j$, and set
\[
d=a_j-a_i,\qquad H_{ij}(z)=\prod_{k\ne i,j}(z-a_k),
\]
where the empty product is $1$. Consider the completed quadratic node
\[
C_{ij}=\OO_E[[u,v,w]]/(uv+w(w-d)),\qquad
T_{ij}=\Spec C_{ij}[1/p].
\]
There is a continuous homomorphism of complete local $\OO_E$-algebras
\begin{equation}\label{d3:eq:pairwiseformalmap}
\widehat B\longrightarrow C_{ij},\qquad
x\longmapsto u,\quad
y\longmapsto H_{ij}(a_i+w)v,\quad
z\longmapsto a_i+w.
\end{equation}
All three images lie in the maximal ideal, so substitution defines every formal series. The defining relation maps to
\[
H_{ij}(a_i+w)\bigl(uv+w(w-d)\bigr)=0.
\]
In particular no inverse of the cofactor $H_{ij}$ is used. Inverting $p$ gives a morphism $q_{ij}:T_{ij}\to T$.

The generated ideals obtained from $J_i,J_j$ are
\[
I_i=(u,w),\qquad I_j=(u,w-d).
\]
For $k\ne i,j$, put $b=a_k-a_i$, so $b\ne0,d$. In $C_{ij}[1/p]$ the identity
\[
b(b-d)=-uv-(w-b)(w+b-d)
\]
shows that $(u,w-b)$ is the unit ideal, since $b(b-d)\in E^\times$.

These computations identify the pullbacks of the line bundles even though $q_{ij}$ need not be flat. The pullback of each invertible $J_k$ surjects onto its generated ideal. For $k=i,j$ that image is an invertible ideal by the quadratic-node calculation, and for every other $k$ it is the unit ideal. A surjection between invertible modules is an isomorphism. Hence
\[
q_{ij}^*J_i\cong I_i,\qquad
q_{ij}^*J_j\cong I_j,\qquad
q_{ij}^*J_k\cong\OO\quad(k\ne i,j).
\]
Also $I_iI_j=(u)$, so $c_1(I_j)=-c_1(I_i)$.

Pulling the cup sum to $T_{ij}$ therefore gives
\[
q_{ij}^*\left(\sum_k\alpha_k\smile c_1(J_k)\right)
=(\alpha_i-\alpha_j)\smile c_1(I_i).
\]
Lemma~\ref{d3:lem:completedcup}, with parameter $-d$, makes the right-hand cup map injective. Its prescribed inversion is the same as inversion of $p$, because $d$ is a nonzero element of the maximal ideal of the DVR $\OO_E$. Thus vanishing of the original sum forces $\alpha_i=\alpha_j$. Taking the pairs $(1,j)$ proves that its kernel is the diagonal. The diagonal already vanishes because $\prod_iJ_i=(x)$, proving \eqref{d3:eq:completedtamecup}.
\end{proof}

\subsection{Actual local coefficients and the line-bundle formula}

\begin{lemma}\label{d3:lem:continuouschern}
Let $E/\Q_p$ be finite and let $R$ be an $E$-algebra. The higher Chern maps used below have natural values in continuous \'etale cohomology, and their finite-coefficient line-bundle product identity holds in that continuous target. No vanishing of $\varprojlim^1$ for $\Spec R$ is required.
\end{lemma}
\begin{proof}
Write $R=\colim R_\lambda$ over its finitely generated $E$-subalgebras. For $X_\lambda=\Spec R_\lambda$, geometric \'etale finiteness and finiteness of local Galois cohomology, through Hochschild--Serre, show that every
$H^q_{\mathrm{\acute et}}(X_\lambda,\mu_{p^\nu}^{\otimes i})$ is finite. See \cite[Theorem~19.1(b)]{MilneEtale} and \cite[Chapter~I, Corollary~2.3]{MilneDuality}; the relevant field has characteristic zero, so this includes $p$-power coefficients and $p=2$. Jannsen's exact sequence \cite[equation~(0.2)]{Jannsen} consequently identifies continuous cohomology with the ordinary inverse limit at each such finite-type stage. The compatible finite-coefficient Chern classes therefore have unique continuous lifts there, and their product identities hold there as well.

Now use $K_3(R)=\colim K_3(R_\lambda)$ and pull those continuous classes to $\Spec R$. An element, an equality of elements, and any finite collection of projective modules descend to one common stage. This proves independence of the stage, naturality, and the product identity. It is the finite-type descent construction of \cite[Section~5.1]{LevineIndK3}, applied here to affine $E$-algebras. In particular it applies to the completed generic rings and the affinoid algebras in this section. We do not identify $H^3_{\mathrm{cont}}(\Spec R,\Z_p(3))$ with $\varprojlim H^3_{\mathrm{\acute et}}(\Spec R,\mu_{p^\nu}^{\otimes3})$, and do not interchange these cohomological limits with the filtered limit of rings.
\end{proof}

\begin{lemma}\label{d3:lem:actuallocal}
For a finite extension $E/\Q_p$, the images of actual elements of $K_3(\OO_E)$ under the rational \'etale Chern map
\[
c_{2,1}:K_3(\OO_E)\longrightarrow H^1(E,\Q_p(2))
\]
span the target over $\Q_p$. If $E/F$ is quadratic, actual integral anti-invariant elements span its anti-invariant subspace, whose $\Q_p$-dimension is $[F:\Q_p]$.
\end{lemma}
\begin{proof}
Levine's finite-coefficient theorem \cite[Theorem~4.12]{LevineIndK3}, also stated as the exact sequence in \cite[Section~5]{WeibelChern}, gives
\[
K_3^M(E)\longrightarrow K_3(E;\Z/p^\nu)
\xrightarrow{c_{2,1}}H^1(E,\mu_{p^\nu}^{\otimes2})\longrightarrow0,
\]
with the induced map on the indecomposable quotient an isomorphism, including $p=2$. Milnor $K_3$ of a local field is divisible \cite[Chapter~III, Exercise~7.4]{Weibel}; hence $K_3^M(E)/p^\nu=0$. The first image is therefore zero, and the displayed Chern map is a natural isomorphism.

The coefficient sequence for $G=K_3(E)$ is
\cite[IV, Section~2]{Weibel}
\[
0\longrightarrow G/p^\nu G\longrightarrow K_3(E;\Z/p^\nu)
\longrightarrow K_2(E)[p^\nu]\longrightarrow0.
\]
The transition on the right is multiplication by $p$. By \cite[III, Theorem~6.2.4]{Weibel}, with the characteristic-zero torsion assertion attributed there to \cite{MerkurjevTorsion}, $K_2(E)$ is the direct sum of a uniquely divisible group and a finite cyclic group. Hence that right-hand inverse limit is zero. Each $G/p^\nu G$ is finite, since it injects into the finite middle group, and this inverse system has vanishing $\varprojlim^1$. The finite groups $H^0(E,\mu_{p^\nu}^{\otimes2})$ also have vanishing $\varprojlim^1$, so the inverse limit of the $H^1$ groups is the continuous group $H^1(E,\Z_p(2))$. Taking inverse limits therefore gives the natural identification
\[
\varprojlim_\nu G/p^\nu G\cong H^1(E,\Z_p(2)).
\]
The actual group $G$ is dense in its completion, so its rational Chern image spans $H^1(E,\Q_p(2))$. There is no Tate-module contribution from $K_2$ hidden in this argument.

Localization lifts every element of $K_3(E)$ to $K_3(\OO_E)$, since $K_2(k_E)=0$. The natural Bloch--Kato isomorphism
\[
\exp_E:E\xrightarrow{\sim}H^1(E,\Q_p(2))
\]
identifies the anti-invariant subspace for a quadratic $E/F$ with the trace-zero subspace of $E$, of dimension $[F:\Q_p]$. Apply $1-\sigma$ to actual lifts whose images span that space. This gives integral classes $\beta$ with $\sigma\beta=-\beta$, without dividing by two. Naturality of the Chern maps and the exponential ensures all these assertions are equivariant. For the exponential and its normalization, see \cite[Section~1.3]{HuberKings}.
\end{proof}

We also need the precise line-bundle product. If $L$ is a line bundle and $\beta\in K_3$, then
\begin{equation}\label{d3:eq:etaleproduct}
c_{3,3}\bigl(([L]-1)\beta\bigr)
=-2c_1(L)\smile c_{2,1}(\beta).
\end{equation}
This is \cite[Theorem~3.11(i)]{WeibelChern}. The $K$-degree is three, so the exceptional degree-two correction in that paper does not apply, even at $p=2$. Lemma~\ref{d3:lem:continuouschern} puts this formula in continuous cohomology before passing to $\Q_p$. The nonzero factor $-2$ is harmless for rational detection.

\subsection{Complete quadratic examples over every $p$-adic base}

\begin{theorem}\label{d3:thm:completed}
Let $V=\OO_F$ for a finite extension $F/\Q_p$, and let $f(z)\in V[z]$ be an Eisenstein quadratic polynomial. Then
\[
\widehat A_f=V[[x,y,z]]/(xy+f(z))
\]
is a regular local ring of dimension three, and
\[
\dim_{\Q}\ker\bigl(K_3(\widehat A_f)_{\Q}
\longrightarrow K_3(\Frac\widehat A_f)_{\Q}\bigr)\geq[F:\Q_p].
\]
Taking $V=\Z_p$, this applies to $f=z^2+p$ for every prime $p$, including $p=2$.
\end{theorem}
\begin{proof}
Write $E=F(\theta)$ for a root of $f$, $W=V[\theta]=\OO_E$, and $\sigma$ for the quadratic involution. Eisenstein's condition proves regularity exactly as before, and the divisor $x=0$ has regular coordinate ring $W[[y]]$. Choose actual anti-invariant classes $\beta_j\in K_3(W)$ whose \'etale Chern classes are $\Q_p$-independent; Lemma~\ref{d3:lem:actuallocal} supplies $[F:\Q_p]$ such classes. Push them forward from $x=0$ to obtain $c_j\in K_3(\widehat A_f)$. Each generic image is zero by localization.

After the finite free base change from $V$ to $W$, put $w=z-\theta$ and $d=\theta-\sigma\theta$. Since both roots are topologically nilpotent and distinct, this is a valid formal change of variables and $0\ne d\in\mathfrak m_W$. The equation becomes
\[
xy+w(w+d)=0.
\]
After inverting $d$, the two components of $x=0$ are disjoint. With $I=(x,w)$ and $\eta_+=[\OO_{(x,w)}]=1-[I]$, flat base change gives
\[
c_j|_{T_d}=\beta_j\eta_++\sigma(\beta_j)\eta_-
=2\beta_j\eta_+,
\qquad \eta_++\eta_-=0.
\]
Applying \eqref{d3:eq:etaleproduct} gives
\[
c_{3,3}(c_j|_{T_d})=4c_1(I)\smile c_{2,1}(\beta_j).
\]
These classes are $\Q_p$-independent by Lemma~\ref{d3:lem:completedcup}. Thus the integral classes $c_j$ are $\Q$-independent after rationalization, proving the assertion.
\end{proof}

\begin{remark}
The simplest residue-two complete example in this theorem is $\Z_2[[x,y,z]]/(2+xy+z^2)$. It is different from the noncomplete residue-two example of Section~\ref{d3:sec:two}, which uses the imaginary quadratic number field $\Q(i)$ to obtain an archimedean source. The completed theorem uses local \'etale regulator classes and requires no imaginary number field.
\end{remark}

\subsection{The explicit class also survives completion at three}\label{d3:sec:explicitcomplete}

\begin{proposition}\label{d3:prop:explicitcomplete}
For
\[
\widehat A=\Z_3[[x,y,z]]/(3+xy+z^2),
\]
the explicit class $c_{\mathrm{exp}}$ of \eqref{d3:eq:explicitfront} has infinite order and zero generic image. Consequently the same explicit class has infinite order in the henselization of the original ring $A$ as well.
\end{proposition}
\begin{proof}
Let $E_0=\Q(u)$, $u^2=-3$, and retain the supported lift $\widetilde c\in K_3(Q)$ constructed in Section~\ref{d3:sec:explicit}. Since $D_+$ and $D_-$ are affine lines over $E_0$, homotopy invariance and their pushforwards give
\[
\widetilde c=\gamma_+\eta_++\gamma_-\eta_-
=(\gamma_+-\gamma_-)\eta_+=\gamma\eta_+,
\qquad \gamma\in K_3(E_0).
\]
The positive complex period in \eqref{d3:eq:positiveperiod} implies $\gamma\ne0$ in $K_3(E_0)_{\Q}$. Borel's theorem \cite[Proposition~12.2]{Borel} gives $\dim_{\Q}K_3(E_0)_{\Q}=1$.

Put $\zeta=(1+u)/2$, so $1-\zeta=\zeta^{-1}$. The special-unit Bloch cycle $[\zeta]_2$ gives a rational class $\beta_{\mathrm{cyc}}\in K_3(E_0)_{\Q}$, compatible with its integral model at three and with the embedding $E_0\hookrightarrow E=\Q_3(u)$. Besser--de Jeu's Theorem~1.10(2), in its unconditional weight-two case, computes its local syntomic regulator as
\[
\pm L_{\mathrm{mod},2}(\zeta)=\pm\operatorname{Li}_2(\zeta),
\]
because $\log\zeta=\log(1-\zeta)=0$ \cite{BesserDeJeu}. We verify directly that this value is nonzero.

Normalize $v_3(3)=1$ and put
\[
 \epsilon=1+\zeta=\frac{3+u}{2}.
\]
Then $v_3(\epsilon)=1/2$ and $\epsilon^3=3u$.
Coleman's reflection and inversion identities, and the duplication
formula, give
\[
 \operatorname{Li}_2(\epsilon)
 =-\operatorname{Li}_2(-\zeta),\qquad
 \operatorname{Li}_2(\zeta)+\operatorname{Li}_2(-\zeta)
 =\frac12\operatorname{Li}_2(\zeta^2)
 =-\frac12\operatorname{Li}_2(-\zeta).
\]
Here $1-\epsilon=-\zeta$, $( -\zeta)^{-1}=\zeta^2$,
and all logarithms of roots of unity vanish. See
\cite[equation~(1.2) and (2.4)]{BesserDeJeu}
and \cite[p.~377]{BesserDeJeuCurves}.
Consequently
\[
 \frac{\operatorname{Li}_2(\zeta)}u
 =\frac{3}{2u}\operatorname{Li}_2(\epsilon)
 =\frac{3}{2u}\sum_{n\geq1}\frac{\epsilon^n}{n^2}.
\]
The term $n=3$ is exactly $1/2$. Every other term has valuation
\[
 \frac{n+1}{2}-2v_3(n)\geq1.
\]
For $v_3(n)=0$ this is immediate; for $v_3(n)=1$ and $n\ne3$
use $n\geq6$; for $r=v_3(n)\geq2$ use
$(3^r+1)/2-2r\geq1$. The valuations tend to infinity.
Galois equivariance and inversion show that the quotient
$\operatorname{Li}_2(\zeta)/u$ lies in $\Q_3$.
Intersecting the resulting congruence modulo $3\OO_E$ with
$\Q_3$ therefore gives
\begin{equation}\label{d3:eq:localdilog}
\boxed{\quad
\frac{\operatorname{Li}_2(\zeta)}u\in\frac12+3\Z_3,
\qquad v_3(\operatorname{Li}_2(\zeta))=\frac12.
\quad}
\end{equation}

The syntomic-to-\'etale map for $\Spec\OO_E$ is compatible with Chern classes, and under the identification
\[ H^1_{\mathrm{syn}}(\OO_E,2)\cong E \]
it is the Bloch--Kato exponential. Precise references are \cite[Example~2.2.4 and Propositions~2.2.7, 2.2.9, 2.3.4]{HuberKings}; see also the proof of \cite[Corollary~5.19]{Tamme}. These statements allow ramified $E/\Q_3$. Since the exponential for $\Q_3(2)$ is an isomorphism, \eqref{d3:eq:localdilog} gives
\[
c_{2,1}(\beta_{\mathrm{cyc}}|_E)\ne0.
\]
In particular $\beta_{\mathrm{cyc}}$ is nonzero globally. The one-dimensionality of $K_3(E_0)_{\Q}$ now gives $\gamma=q\beta_{\mathrm{cyc}}$ in $K_3(E_0)_{\Q}$ for some $q\in\Q^\times$. No exact comparison scalar between a Dennis--Stein presentation and a Bloch lift has been assumed or is needed.

On the completed split generic ring $T$ the image of $c_{\mathrm{exp}}$ equals the image of $\widetilde c$: indeed $1-z^2$ is a unit in $\widehat A$, and the equality $\widetilde c|_U=2\alpha$ from Section~\ref{d3:sec:explicit} pulls back there. Thus
\[
c_{\mathrm{exp}}|_T=q\beta_{\mathrm{cyc}}\eta_+\qquad\text{in }K_3(T)_{\Q}.
\]
Its \'etale Chern class is nonzero by \eqref{d3:eq:etaleproduct} and Lemma~\ref{d3:lem:completedcup}, with $w=z-u$, $d=2u$. Therefore $c_{\mathrm{exp}}$ has infinite order in $K_3(\widehat A)$. Its generic image is zero by the integral identity \eqref{d3:eq:explicitgeneric}, which holds after completion. The map from the original local ring to its completion factors through its henselization, so infinite order after completion also proves infinite order there.
\end{proof}

\subsection{The explicit class also survives completion at two}
\label{d3:sec:explicittwocomplete}

\begin{lemma}\label{d3:lem:localdilogtwo}
For the Coleman dilogarithm on $E=\Q_2(i)$, $i^2=-1$, normalized
by its power series at zero and with $\log2=0$, one has
\begin{equation}\label{d3:eq:localdilogtwo}
 \frac{\operatorname{Li}_2(i)}{i}\in\frac12+\Z_2,
 \qquad v_2(\operatorname{Li}_2(i))=-1,
 \qquad v_2(2)=1.
\end{equation}
\end{lemma}
\begin{proof}
The modified dilogarithm
$D(z)=\operatorname{Li}_2(z)+\frac12\log z\log(1-z)$
satisfies $D(1-z)=-D(z)$ and $D(z^{-1})=-D(z)$;
see \cite[equation~(1.8), p.~307, and p.~377]{BesserDeJeuCurves}.
Put $q=1-i$, so $v_2(q)=1/2$.
Since $\log i=0$, reflection gives the convergent expansion
\[
 \operatorname{Li}_2(i)
 =-\operatorname{Li}_2(q)
 =-\sum_{n\geq1}\frac{q^n}{n^2}.
\]
Now $q^2=-2i$, $q^4=-4$, and $q^8=16$. Thus the fourth and eighth
terms cancel exactly, while the second term is $-i/2$.
For $n\notin\{2,4,8\}$,
\[
 v_2(q^n/n^2)=n/2-2v_2(n)\geq0.
\]
Indeed, writing $n=2^rk$ with $k$ odd proves this immediately
for $r\leq3$; for $r\geq4$ use
$2^{r-1}k-2r\geq2^{r-1}-2r\geq0$.
These valuations tend to infinity. Hence
$\operatorname{Li}_2(i)-i/2\in\OO_E$, proving the asserted valuation.
Galois equivariance and inversion give
$\sigma(\operatorname{Li}_2(i))=-\operatorname{Li}_2(i)$
for $\sigma(i)=-i$; see
\cite[Remark~2.3 and equation~(1.2)]{BesserDeJeu}.
Thus $\operatorname{Li}_2(i)/i\in\Q_2$, and intersecting the integral
congruence with $\Q_2$ proves \eqref{d3:eq:localdilogtwo}.
\end{proof}

\begin{proposition}\label{d3:prop:explicittwocomplete}
The explicit triple $c_2$ of \eqref{d3:eq:explicittwoprime} has
infinite order and zero generic image in
\[
 \widehat A_2
 =\Z_2[[x,y,w]]/(xy+w^2+2w+2),\qquad t=1+w.
\]
The same holds in the henselization of its original local ring.
\end{proposition}
\begin{proof}
Let $E_0=\Q(i)$ and let $\widetilde c_2$ be the supported lift
of \eqref{d3:eq:allprimeslift}. Homotopy invariance on the two
affine-line divisors and $\eta_-=-\eta_+$ give
\[
 \widetilde c_2=\gamma\eta_+,\qquad
 \gamma\in K_3(E_0)_{\Q},\qquad
 \eta_+=1-[(x,t-i)].
\]
Its nonzero period \eqref{d3:eq:twocyclotomicperiod} implies
$\gamma\ne0$.

The rational Bloch cycle $[i]_2$ defines
$\beta_i\in K_3(E_0)_{\Q}$. Interpret it over
$\Z[i]_{(1-i)}$ by the rational localization isomorphism.
Theorem~1.12 of \cite{BesserDeJeu} applies to every root of unity,
including $p$-power order, and gives its local syntomic regulator
as $\pm\operatorname{Li}_2(i)$ in $E=\Q_2(i)$.
The theorem has no odd-prime or $p>$weight restriction.
The Chern-compatible syntomic-to-\'etale map at the point
$\Spec\Z_2[i]$ is the Bloch--Kato exponential, also at $p=2$;
see \cite[Example~2.2.4 and Propositions~2.2.7, 2.3.4]{HuberKings}
and \cite[Corollary~5.19 and its proof]{Tamme}.
This exponential is an isomorphism in twist two, so
Lemma~\ref{d3:lem:localdilogtwo} proves
$c_{2,1}(\beta_i|_E)\ne0$.
Borel's equality $\dim_{\Q}K_3(\Q(i))_{\Q}=1$
\cite[Proposition~12.2]{Borel} now gives
$\gamma=a\beta_i$ for some $a\in\Q^\times$.

After finite free coefficient extension to $W=\Z_2[i]$, put
\[
 s=t-i=w-(i-1),\qquad d=2i.
\]
Because $i-1$ is topologically nilpotent, this is a valid formal
coordinate change, and the split completed ring is
\[
 \widehat B=W[[x,y,s]]/(xy+s(s+d)).
\]
Let $T=\Spec\widehat B[1/2]$ and $I=(x,s)$ on $T$.
The open used for the supported lift is $D(tg)$,
where $g=1+t+t^2$. Both $t=1+w$ and
$g=3+3w+w^2$ are units after completion. Thus
\eqref{d3:eq:allprimeslift} pulls back to
$c_2|_T=a\beta_i(1-[I])$.
Equation~\eqref{d3:eq:etaleproduct} gives
\[
 c_{3,3}(c_2|_T)
 =2a\,c_1(I)\smile c_{2,1}(\beta_i|_E)\ne0
\]
by Lemma~\ref{d3:lem:completedcup}, with parameter $d=2i$.
Its scaled model is $XY+Z(Z+1)=0$, obtained from
$x=dX$, $y=dY$, $s=dZ$; hence that lemma applies in residue
characteristic two.
This proves infinite order in $K_3(\widehat A_2)$.
The integral generic-zero identity for $c_2$ remains valid after
completion. Finally the completion map factors through the
henselization, giving the same conclusion there.
\end{proof}


\subsection{Cyclotomic local values and a common normalization}
\label{d3:sec:cyclotomicnormalization}


\begin{lemma}\label{d3:lem:cyclotomicdilognonzero}
Let $p$ be odd, $r\geq1$, and let $\rho$ be a primitive $p^r$-th
root of unity in $E=\Q_p(\rho)$. For the Coleman dilogarithm
normalized by its series at zero and $\log p=0$, one has
\[
 v_p(\operatorname{Li}_2(\rho))
 =\frac p{p-1}-2r,\qquad
 \operatorname{Li}_2(-\rho)\ne0.
\]
\end{lemma}
\begin{proof}
Put $e=p^{r-1}(p-1)$ and $q=1-\rho$, so $v_p(q)=1/e$.
Coleman reflection for
$D(z)=\operatorname{Li}_2(z)+\frac12\log z\log(1-z)$
and $\log\rho=0$ give
\[
 \operatorname{Li}_2(\rho)
 =-\sum_{n\geq1}\frac{q^n}{n^2};
\]
see \cite[p.~377]{BesserDeJeuCurves}.
Writing $n=p^sk$ with $p\nmid k$, its term valuation is
$p^sk/e-2s$. For $k=1$, the consecutive differences are
$p^{s-r+1}-2$, negative for $s<r$ and positive for $s\geq r$.
Thus the unique least term is $n=p^r$, with the asserted
valuation; every other term has valuation at least one more,
and the valuations tend to infinity.

Let $\sigma_2(\rho)=\rho^2$. Galois equivariance and duplication
\cite[Remark~2.3 and equation~(2.4)]{BesserDeJeu} give
\[
 \operatorname{Li}_2(-\rho)
 =\tfrac12\sigma_2(\operatorname{Li}_2(\rho))
       -\operatorname{Li}_2(\rho).
\]
If this vanished, iteration through the finite order $d$ of
$\sigma_2$ would give
$(2^d-1)\operatorname{Li}_2(\rho)=0$, a contradiction.
\end{proof}

\begin{lemma}\label{d3:lem:cyclotomicnormalization}
Fix the Besser--de Jeu weight-two symbol map with one natural
sign convention. For a root of unity $a$ of order $N$ in a
number field $F$, with $a,1-a\in\OO_F^\times$, let
$\beta_{\mathrm{BDJ}}(a)$ denote its rational symbol class and set
\[
 \beta_{\mathrm{BT}}(a)=\frac1N\operatorname{bl}(N[a])
       \in K_3(F)_{\Q}
\]
using \cite[Theorem~4.9]{BunkeTamme}.
There is one $s\in\Q^\times$, independent of $F$ and $a$, such that
\[
 \beta_{\mathrm{BDJ}}(a)=s\,\beta_{\mathrm{BT}}(a).
\]
\end{lemma}
\begin{proof}
The multiple $N[a]$ has zero integral wedge boundary because
$a^N=1$, including with the antisymmetric-tensor convention.
The BT class is unique rationally by Borel injectivity
\cite[Section~6]{BunkeTamme} and has
regulator $(-iD_{\mathrm{BW}}(\sigma a))_\sigma$.

The field map in \cite[equation~(1.1), Theorem~1.10(1)]{BesserDeJeu}
is de Jeu's map. Its complex regulator is computed in
\cite[Proposition~4.1 and Remark~5.2]{DeJeuWedge}, also restated in \cite[Theorem~2.3]{DeJeuCurves}, simultaneously
at all embeddings, by $P_2=D_{\mathrm{BW}}$.
Transporting its degree-dependent target identifications to the
BT target therefore gives one nonzero real constant $C$ with
\[
 r_{\mathrm{BT}}(\beta_{\mathrm{BDJ}}(a))_\sigma
       =C\,iD_{\mathrm{BW}}(\sigma a).
\]
The sign choice is made once in the fixed weight and degree,
as in \cite[Remark~1.7]{BesserDeJeu}.

Calibrate at $\zeta=(1+\sqrt{-3})/2$ in $\Q(\sqrt{-3})$.
Both classes have nonzero regulator and lie in its rank-one
rational $K_3$ \cite[Proposition~12.2]{Borel}, so
$\beta_{\mathrm{BDJ}}(\zeta)=s\beta_{\mathrm{BT}}(\zeta)$
for some $s\in\Q^\times$.
Their regulator values imply $C=-s$.
For every $F,a$, the two classes in the statement now have
the same regulator at all embeddings. Borel injectivity proves
their equality.
\end{proof}


\subsection{The explicit odd-prime triples survive completion}
\label{d3:sec:explicitoddcomplete}

We identify all supported coefficients before applying the local
detector. This avoids using a single complex embedding to identify
a class in a higher-rank number-field $K$-group.

\begin{lemma}\label{d3:lem:explicitoddcoefficients}
Let $p$ be odd, let $E_0=\Q(\mu_{2p})$, and enumerate its primitive
$2p$-th roots as $a_1,\ldots,a_{p-1}$. On
\[
 Q=\Spec E_0[x,y,a]/(xy+\Phi_{2p}(a))
\]
put $D_i=(x,a-a_i)$, $J_i=\OO_Q(-D_i)$ and
$\eta_i=[\OO_{D_i}]=1-[J_i]$.
Normalize the rational Bloch classes $\beta_{\rm BT}(a_i)$ by
\[
 r_2(\beta_{\rm BT}(a_i))_\sigma
   =-iD_{\rm BW}(\sigma(a_i))
\]
in the Bunke--Tamme cone coordinate, at every complex embedding
$\sigma$ of $E_0$.
Then every supported lift $\widetilde c_p$ in
\eqref{d3:eq:allprimeslift} satisfies
\begin{equation}\label{d3:eq:explicitoddcoefficient}
 \widetilde c_p=-2p\sum_{i=1}^{p-1}\beta_{\rm BT}(a_i)\eta_i
 \qquad\text{in }K_3(Q)_{\Q}.
\end{equation}
\end{lemma}

\begin{proof}
The stated normalization defines actual rational classes:
$a_i,1-a_i$ are global units, because $\Phi_{2p}(1)=1$,
and $2p[a_i]$ lies in the kernel of
$[a]\mapsto a\wedge(1-a)$.
One may therefore take
\[
 \beta_{\rm BT}(a_i)=\frac1{2p}\operatorname{bl}(2p[a_i])|_{E_0}
\]
using \cite[Theorem~4.9 and identification~(4.8)]{BunkeTamme}.
Borel injectivity \cite[Section~6]{BunkeTamme} makes this rational class independent of the
integral choice and equivariant under field embeddings.

Homotopy invariance on the affine-line components of $D_x$ and
the projection formula give
$\widetilde c_p=\sum_i\gamma_i\eta_i$, with
$\gamma_i\in K_3(E_0)$.
Fix distinct $i,j$ and an embedding $\sigma:E_0\hookrightarrow\C$.
The chord from $\sigma(a_i)$ to $\sigma(a_j)$ avoids zero:
primitive $2p$-th roots cannot be antipodal when $p$ is odd.
Its interior lies inside the unit circle and meets no other root,
and the whole chord has $\Re a<1$. Thus the principal
$\operatorname{Log}(1-a)$ and $\operatorname{Li}_2(a)$ are defined
on one neighborhood of it.

Orient the associated matching sphere from the $i$th root to
the $j$th root as in \eqref{d3:eq:cyclotomicendpoint}.
That formula gives the real period
\[
 2\pi(2p)\bigl(D_{\rm BW}(\sigma(a_j))
                  -D_{\rm BW}(\sigma(a_i))\bigr).
\]
The sphere meets $D_i$ with oriented intersection $+1$,
$D_j$ with intersection $-1$, and misses the other divisors.
Indeed, near its starting pole $\tau+1$ is a positive smooth
multiple of $|x|^2$, whereas near its ending pole $1-\tau$ is.
Consequently the periods of $\operatorname{ch}_1(\eta_i)$
and $\operatorname{ch}_1(\eta_j)$ are $+2\pi i$ and $-2\pi i$.

Write $r_2(\gamma_k)_\sigma=iq_k(\sigma)$ with $q_k(\sigma)$ real.
Multiplicativity therefore gives
\[
 q_i(\sigma)-q_j(\sigma)
 =2p\bigl(D_{\rm BW}(\sigma(a_i))
                  -D_{\rm BW}(\sigma(a_j))\bigr).
\]
Since this holds at every complex embedding, Borel injectivity
implies
\[
 \gamma_i-\gamma_j
 =-2p\bigl(\beta_{\rm BT}(a_i)-\beta_{\rm BT}(a_j)\bigr)
 \quad\text{in }K_3(E_0)_{\Q}.
\]
Thus $\gamma_i+2p\beta_{\rm BT}(a_i)$ is one common coefficient.
Its contribution is zero because
$\sum_i\eta_i=[\OO_{V(x)}]=0$, proving
\eqref{d3:eq:explicitoddcoefficient}.
\end{proof}

\begin{proposition}\label{d3:prop:explicitoddcomplete}
For every odd prime $p$, the explicit triple $c_p$ in
\eqref{d3:eq:explicitoddprime} has infinite order and zero generic
image in
\[
 \widehat A_p=\Z_p[[x,y,z]]/(xy+\Phi_{2p}(z-1)).
\]
The same assertions hold in the henselization of its original
arithmetic local ring.
\end{proposition}

\begin{proof}
Put $E=\Q_p(\mu_{2p})$ and $W=\OO_E$.
After the finite free coefficient extension to $W$, the completed
ring has the split form
\[
 \widehat B=
 W[[x,y,z]]\Big/\left(xy+\prod_i(z-(1+a_i))\right).
\]
All $1+a_i$ lie in the maximal ideal of $W$.
Let $T=\Spec\widehat B[1/p]$ and still write
$J_i=(x,z-(1+a_i))$ on $T$.
The finite open $U=D(a(1-a))$, with $a=z-1$, contains this generic
scheme, since $a$ and $1-a$ are units in $\widehat B$.
Thus \eqref{d3:eq:explicitoddcoefficient} pulls back to $c_p|_T$.

By Lemma~\ref{d3:lem:cyclotomicnormalization}, the classes
$\beta_{\rm BT}(a_i)$ are one common nonzero rational multiple of
the de Jeu classes represented by $[a_i]_2$.
Theorem~1.12 of \cite{BesserDeJeu}, followed by the point
syntomic-to-\'etale comparison and the Bloch--Kato isomorphism,
therefore gives
\[
 c_{2,1}(\beta_{\rm BT}(a_i)|_E)
   =\kappa\,\exp_E(\operatorname{Li}_2(a_i))
\]
with one $\kappa\in\Q_p^\times$, independent of $i$.
Lemma~\ref{d3:lem:cyclotomicdilognonzero} proves
$\operatorname{Li}_2(a_i)\ne0$, and inversion gives
$\operatorname{Li}_2(a_i^{-1})=-\operatorname{Li}_2(a_i)$.
The displayed coefficient vector is consequently non-diagonal.

The line-bundle formula \eqref{d3:eq:etaleproduct} gives
\[
 c_{3,3}(c_p|_T)
 =-4p\sum_i c_1(J_i)\smile
                c_{2,1}(\beta_{\rm BT}(a_i)|_E)\ne0
\]
by Lemma~\ref{d3:lem:completedrootcup}.
This proves infinite order after completion.
The integral vanishing identity \eqref{d3:eq:allprimesDS} remains
valid after every base change, so its generic image is zero.
Finally, the completion map factors through the henselization,
which proves both assertions there as well.
\end{proof}

\section{Completed Eisenstein and wild cyclotomic families}\label{d3:sec:completedfamilies}
Pairwise ruling maps detect all root coefficients modulo the diagonal.
We apply this detector first to arbitrary Eisenstein polynomials and
then to explicit cyclotomic coefficients with independent wild conjugates.

\subsection{Eisenstein polynomials and unbounded rank}
\label{d3:sec:completedarbitraryrank}\label{d3:sec:completedtamerank}

Lemma~\ref{d3:lem:completedrootcup} applies to every Eisenstein polynomial, including wild degrees. We combine it with actual local coefficients to obtain the following rank bound.


\begin{theorem}
\label{d3:thm:completedeisensteinrank}\label{d3:thm:completedtamerank}
Let $F/\Q_p$ be finite, let $V=\OO_F$, and let $f(z)\in V[z]$ be an Eisenstein polynomial of degree $N\geq2$. Then
\[
A_f=V[[x,y,z]]/(xy+f(z))
\]
is a complete regular local ring of dimension three with residue field $k_F$, and
\[
\dim_{\Q}\ker\left(K_3(A_f)_{\Q}
\longrightarrow K_3(\Frac A_f)_{\Q}\right)
\geq [F:\Q_p](N-1).
\]
The lower bound is witnessed by actual integral classes whose generic images vanish integrally.
\end{theorem}

\begin{proof}
In the regular local ring $V[[x,y,z]]$, the equation $xy+f(z)$ has a nonzero linear term in a uniformizer of $V$. Indeed the constant coefficient of $f$ has valuation one, while all other terms lie in the square of the maximal ideal. Its quotient is therefore regular of dimension three, with maximal ideal $(x,y,z)$ and residue field $k_F$.

Choose a root $\theta$ of $f$ and put
\[
K=F(\theta),\qquad W=V[\theta]=\OO_K.
\]
Eisenstein's criterion gives $[K:F]=N$, and the divisor $x=0$ has regular coordinate ring
\[
A_f/(x)\cong W[[y]].
\]
For an actual $\beta\in K_3(W)$ define
\[
c(\beta)=i_*(\beta|_{W[[y]]})\in K_3(A_f).
\]
Quillen localization \cite{Quillen} makes its image in $K_3(A_f[1/x])$, hence in $K_3(\Frac A_f)$, zero.

Let $E/F$ be a splitting field of $f$, let $\sigma_1,\ldots,\sigma_N:K\hookrightarrow E$ be its $F$-embeddings, and write $a_i=\sigma_i(\theta)$. Characteristic zero makes the roots distinct; Eisenstein's condition puts every $a_i$ in $\mathfrak m_{\OO_E}$. Finite free coefficient extension gives
\[
A_f\otimes_V\OO_E
\cong
\widehat B=\OO_E[[x,y,z]]/
\left(xy+\prod_i(z-a_i)\right).
\]
Set $T=\Spec\widehat B[1/p]$. This scheme is regular: it is the finite \'etale base change of $\Spec A_f[1/p]$ along the separable field extension $E/F$.

After this flat base change and inversion of $p$, the divisor $W[[y]]$ splits into the disjoint regular Cartier divisors
\[
D_i=V(x,z-a_i)\subset T,\qquad
D_i\cong\Spec\OO_E[[y]][1/p].
\]
Let $J_i=(x,z-a_i)$ and $\eta_i=[\OO_{D_i}]=1-[J_i]$. Flat base change and the projection formula give
\begin{equation}\label{d3:eq:tamegysin}
c(\beta)|_T=\sum_{i=1}^N\sigma_i(\beta)\eta_i.
\end{equation}
Here $\sigma_i(\beta)$ denotes restriction to $K_3(K)$ followed by extension along $\sigma_i$, and then the coefficient pullback from $K_3(E)$ to $K_3(T)$. The displayed formal divisor rings are localized power-series rings; they are not being replaced by $E[[y]]$.

Put
\[
\alpha=c_{2,1}(\beta|_K)\in H^1(K,\Q_p(2)),
\qquad a=\exp_K^{-1}(\alpha)\in K.
\]
Naturality of the Bloch--Kato exponential gives
\begin{equation}\label{d3:eq:tamecoefficientvector}
\bigl(c_{2,1}(\sigma_i(\beta))\bigr)_i
=\bigl(\exp_E(\sigma_i(a))\bigr)_i.
\end{equation}
The continuous Chern maps of Lemma~\ref{d3:lem:continuouschern} and the product formula \eqref{d3:eq:etaleproduct} turn \eqref{d3:eq:tamegysin} into
\[
c_{3,3}(c(\beta)|_T)
=2\sum_i c_{2,1}(\sigma_i(\beta))\smile c_1(J_i).
\]
The factor $2$ remains nonzero over $\Q_p$ for every $p$.

By Lemma~\ref{d3:lem:completedrootcup}, the cup detector kills exactly the diagonal coefficient vectors. By \eqref{d3:eq:tamecoefficientvector}, this means that all the $\sigma_i(a)$ are equal. This happens exactly when $a\in F$: every element of $\Gal(E/F)$ restricts to one of the embeddings $\sigma_i$, so equality of all conjugates makes $a$ fixed by $\Gal(E/F)$. Consequently the map
\[
g:K\longrightarrow E^N/E(1,\ldots,1),
\qquad a\longmapsto[(\sigma_i(a))_i]
\]
has kernel $F$. Its image has $\Q_p$-dimension
$[K:\Q_p]-[F:\Q_p]=[F:\Q_p](N-1)$.

By Lemma~\ref{d3:lem:actuallocal}, the values
$\exp_K^{-1}c_{2,1}(\beta)$ of actual $\beta\in K_3(W)$ span $K$ over $\Q_p$. Choose $[F:\Q_p](N-1)$ actual integral coefficient classes whose images under $g$ are $\Q_p$-independent. The displayed detector makes their pushforwards $\Q$-independent after rationalization: a rational relation would give a $\Q_p$-linear relation among those detected coefficient vectors. Each generic image already vanishes integrally by localization, proving the theorem.
\end{proof}

In particular, for every fixed prime $p$ and every integer $N\geq2$, with no restriction on divisibility by $p$,
\[
A_N=\Z_p[[x,y,z]]/(p+xy+z^N)
\]
has residue field $\F_p$, dimension three, and a rational generic kernel in $K_3$ of dimension at least $N-1$. Thus these dimensions are unbounded with the prime, residue field, and local dimension all fixed.


\subsection{Explicit wild-orbit families at every residue prime}
\label{d3:sec:explicitwildorbit}

We prove a local dilogarithm independence statement before applying the
geometric detector. It gives the whole minus space at the primes two
and three. For larger primes the assertion concerns one wild orbit;
no nonvanishing assertion for the other odd tame characters is needed.
Throughout, the Coleman dilogarithm has its usual series at zero and
the logarithm is chosen with $\log p=0$.

\begin{lemma}\label{d3:lem:wilddepletion}
Let $\rho_r$ be a primitive $p^r$-th root with
$\rho_r^p=\rho_{r-1}$, put $R_r=\operatorname{Li}_2(\rho_r)$,
and normalize $v_p(p)=1$.
For odd $p$ and $r\geq2$, write $e=p^{r-1}(p-1)$. Then
\[
 v_p(R_r-p^{-2}R_{r-1})=
 \begin{cases}
 -3/e,&p=3,\\
 -(p-2)/e,&p\geq5.
 \end{cases}
\]
For $p=2$ and $r\geq3$ one has
\[
 v_2(R_r-\tfrac14R_{r-1})=-2^{-(r-2)}.
\]
\end{lemma}
\begin{proof}
Put $F_p(z)=\operatorname{Li}_2(z)-p^{-2}\operatorname{Li}_2(z^p)$.
Theorem~2.7(3) and equation~(3.3) of
\cite{BesserDeJeuAlgorithm} identify it with the analytic continuation
of $\sum_{p\nmid n}z^n/n^2$ to $|1-z|>p^{-1/(p-1)}$.
These statements apply to every prime, including two.
For $j\geq0$ put
\[
 S_j(t)=\sum_{m\geq0}m^jt^m=\frac{P_j(t)}{(1-t)^{j+1}},
 \quad P_j\in\Z[t],\quad P_0=1,
 \qquad
 A_j(z)=\sum_{a=1}^{p-1}\frac{z^a}{a^{j+2}},
\]
with $m^0=1$ also at $m=0$.
Expanding $(a+pm)^{-2}$ first on $|z|<1$ gives
\begin{equation}\label{d3:eq:wilddepletedexpansion}
 F_p(z)=\sum_{j\geq0}(-1)^j(j+1)p^j A_j(z)S_j(z^p).
\end{equation}
For any rational $0<a<1$, the right side converges uniformly on
\[
 |z|\leq1,\qquad v_p(1-z^p)\leq a,
\]
since its $j$th summand has valuation at least $j-(j+1)a$.
This domain is the connected annulus
$p^{-a/p}\leq|1-z|\leq1$, and lies in the quoted domain of $F_p$.
The rigid analytic identity principle proves
\eqref{d3:eq:wilddepletedexpansion} there. For all the roots under
consideration one may take $a=2/3$.

At $p=2$ one has $A_j(z)=z$. Write
$d=v_2(1-\rho_r^2)=2^{-(r-2)}\leq1/2$.
The $j=0$ term is $\rho_r/(1-\rho_r^2)$, of valuation $-d$;
every $j\geq1$ term has valuation at least
$j-(j+1)d=-d+j(1-d)>-d$. This proves the assertion at two.

At $p=3$ the numerator of the $j=0$ term is
$\rho_r+\rho_r^2/4$, which reduces to $2$.
With $d=v_3(1-\rho_r^3)=3/e\leq1/2$, the same argument gives $-3/e$.

Suppose now that $p\geq5$. For $A=A_0$, reduction modulo $p$ gives
\[
 A(1)=A'(1)=0,\qquad A''(1)/2=-1/2\quad\text{in }\F_p.
\]
Indeed these are the usual sums of inverse squares, inverses,
and $(a-1)/(2a)$ over $a\in\F_p^\times$.
Since $q=1-\rho_r$ has valuation $1/e$ and $2/e<1$,
the unique least term of $A(1-q)$ is its quadratic term.
Thus $v_p(A(\rho_r))=2/e$ and the $j=0$ term has valuation $(2-p)/e$.
For $j\geq1$ the difference between its lower valuation bound
$j-(j+1)p/e$ and $(2-p)/e$ is
\[
 j(1-p/e)-2/e\geq1-(p+2)/e>0,
\]
because $e\geq p(p-1)>p+2$.
Again the $j=0$ term is uniquely minimal.
\end{proof}

\begin{proposition}\label{d3:prop:wildorbit}
For odd $p$, let $E_r=\Q_p(\rho_r)$ and
$\Gamma_r=\Gal(E_r/E_1)$.
The $p^{r-1}$ elements
\[
 \{\sigma\operatorname{Li}_2(-\rho_r):\sigma\in\Gamma_r\}
\]
are linearly independent over $\Q_p$.
For $p=3$ they form a basis of $E_r^-$.

At $p=2$, let $r\geq2$, $E_r=\Q_2(\rho_r)$, and
$\Gamma_r=\Gal(E_r/\Q_2(i))$.
Put
\[
 \Psi_2(\rho_r)=\tfrac13\operatorname{Li}_2(\rho_r^3)
                              -\operatorname{Li}_2(\rho_r).
\]
Its $2^{r-2}$ conjugates under $\Gamma_r$ form a $\Q_2$-basis of
$E_r^-$.
\end{proposition}
\begin{proof}
For odd $p$, write
$\Gal(E_r/\Q_p)=\Delta_r\times\Gamma_r$, with $|\Delta_r|=p-1$,
and let $\omega:\Delta_r\to\Q_p^\times$ be the Teichm\"uller
character: $\omega(\sigma_a)\equiv a\pmod p$ when
$\sigma_a(\rho_r)=\rho_r^a$.
Use its idempotent
\[
 e_\omega=\frac1{p-1}\sum_{\sigma\in\Delta_r}
                        \omega(\sigma)^{-1}\sigma.
\]
If $y\ne0$ has valuation $m/e$ with $m\equiv1\pmod{p-1}$,
then $e_\omega y$ has the same valuation as $y$.
In fact, write $y=q^m u$ with $u$ a unit. Since the extension is
totally ramified, $\sigma_a(u)/u$ reduces to $1$, while
$\sigma_a(q)/q$ reduces to $a$. Each summand of $e_\omega y/y$ reduces to $1$, as does their average.

The integer numerators $-3$ and $-(p-2)$ in
Lemma~\ref{d3:lem:wilddepletion} are congruent to $1$ modulo $p-1$.
Consequently the $\omega$ projection of every depleted value is nonzero.
The base value satisfies
$v_p(R_1)=-(p-2)/(p-1)$ by
Lemma~\ref{d3:lem:cyclotomicdilognonzero}, and its $\omega$ projection
is nonzero for the same reason.

Put $C_r=p^{2r-1}R_r$.
Distribution \cite[Proposition~2.10(1)]{BesserDeJeuAlgorithm} gives
\[
 \operatorname{Tr}_{E_r/E_{r-1}}R_r=p^{-1}R_{r-1},
 \qquad p^{-1}\operatorname{Tr}_{E_r/E_{r-1}}C_r=C_{r-1}.
\]
Hence the successive rational isotypic projections of
$e_\omega C_r$ under the cyclic group $\Gamma_r$ are
\[
 e_\omega C_1,\quad e_\omega(C_2-C_1),\quad\ldots,
 \quad e_\omega(C_r-C_{r-1}).
\]
They are all nonzero. The normal basis theorem identifies
$e_\omega E_r$ with the regular $\Q_p[\Gamma_r]$-module
\cite[Chapter~I, p.~34]{MilneDuality}.
Its irreducible constituents have minimal polynomials
$T-1,\Phi_p(T),\ldots,\Phi_{p^{r-1}}(T)$, each once;
these are irreducible over $\Q_p$.
Thus $e_\omega R_r$ is cyclic and its orbit is a basis.
Projecting any relation shows that the orbit of $R_r$ itself
is independent. Duplication gives
\[
 \operatorname{Li}_2(-\rho_r)=(\tfrac12\sigma_2-1)R_r.
\]
The displayed operator is invertible, commutes with $\Gamma_r$,
and preserves independence: $\sigma_2$ has finite order whereas
$2$ is not a root of unity.
Inversion makes every value anti-invariant. At $p=3$,
$\Delta_r=\{1,\text{inversion}\}$, so this is the whole minus space.

At $p=2$ use the base level $r=2$ and
$\Gal(E_r/\Q_2)=\{1,\text{inversion}\}\times\Gamma_r$.
The group $\Gamma_r$ is cyclic of order $2^{r-2}$, generated by
$\sigma_5$. Its minus space is its regular representation
\cite[Chapter~I, p.~34]{MilneDuality}.
Set again $C_r=2^{2r-1}R_r$; distribution gives average trace
$C_r\mapsto C_{r-1}$. The base $C_2$ is nonzero by
Lemma~\ref{d3:lem:localdilogtwo}, and every
$C_j-C_{j-1}$ for $j\geq3$ is nonzero by
Lemma~\ref{d3:lem:wilddepletion}.
The same multiplicity-one argument, now using
$T-1,\Phi_2(T),\ldots,\Phi_{2^{r-2}}(T)$ over $\Q_2$,
shows that $R_r$ is a cyclic vector in $E_r^-$.
Finally $\Psi_2(\rho_r)=(\sigma_3/3-1)R_r$, an invertible
operator commuting with $\Gamma_r$, so it too is cyclic.
\end{proof}

\begin{theorem}
\label{d3:thm:explicitwildtriples}
For odd $p$ and $r\geq1$ put
\[
 A_{p,r}=\left(\Z_{(p)}[x,y,z]/
       (xy+\Phi_{2p^r}(z-1))\right)_{(p,x,y,z)}.
\]
Write $a=z-1$, $n=2p^r$, and
\[
 H=\frac{a^n-1}{\Phi_{2p^r}(a)},\qquad
 \lambda=1-x,\qquad \mu=1+\lambda yH.
\]
For $0\leq j<p^{r-1}$, let $k_j$ be the unique odd integer with $1\leq k_j<2p^r$ and
$k_j\equiv(1+p)^j\pmod{p^r}$.
Then the explicit classes
\[
 c_j=k_j\{1-a^{k_j},\lambda,\mu\}\in K_3(A_{p,r})
\]
are linearly independent over $\Q$ after tensoring with $\Q$,
and each has zero generic image integrally.
Their images retain these properties in the henselization and
completion of $A_{p,r}$.

At $p=2$, for $r\geq2$ put
\[
 A_{2,r}=\left(\Z_{(2)}[x,y,t]/
                 (xy+\Phi_{2^r}(t))\right)_{(2,x,y,t-1)},
 \qquad n=3\cdot2^r,
\]
and set $H=(t^n-1)/\Phi_{2^r}(t)$,
$\lambda=1-x$, $\mu=1+\lambda yH$.
Let $k_j$ be the least positive representative of $5^j$ modulo
$2^r$, for $0\leq j<2^{r-2}$.
The same assertions hold for the $2^{r-2}$ classes
\[
 c_j=k_j\{1+t^{k_j}+t^{2k_j},\lambda,\mu\}.
\]
All the displayed local rings are regular of dimension three.
\end{theorem}
\begin{proof}
The polynomials $\Phi_{2p^r}(z-1)$ for odd $p$ and
$\Phi_{2^r}(1+z)$ at two are Eisenstein.
Their constant terms are respectively $p$ and $2$; modulo that
prime they are powers of $z$. This proves regularity and dimension
as in Theorem~\ref{d3:thm:eisenstein}.
Every entry of each displayed triple is a unit: at the closed
point $a=-1$ for odd $p$, and $t=1$ at two.

For either prime let $u$ denote $a$ or $t$, and put $b=yH$.
Then $1-xb=u^n$ and the Dennis--Stein identity
\eqref{d3:eq:allprimesDS} applies to
$\delta=\langle x,b\rangle$.
In the odd case, for any allowed $k$,
\[
 k\{1-a^k,a^n\}=\{1-a^k,(a^k)^n\}=0.
\]
At two, use the finite \'etale quadratic extension adjoining
$\zeta_3$. Steinberg and the transfer projection formula give
\[
 \begin{split}
 0&=\operatorname{Tr}\{1-\zeta_3t^k,(\zeta_3t^k)^n\}\\
  &=\{1+t^k+t^{2k},t^{kn}\}
   =k\{1+t^k+t^{2k},t^n\}.
 \end{split}
\]
Here $3\mid n$, and the norm of $1-\zeta_3t^k$ is exactly
$1+t^k+t^{2k}$, with no extra factor of two.
Only oddness of $k$ is needed; $k$ may be divisible by $3$.
If $g_k$ denotes the first entry of the corresponding triple, then
\[
 k[g_k]\delta=k\{g_k,\lambda,\mu\},\qquad
 k[g_k]\delta|_{D(x)}=0.
\]
This proves integral generic vanishing before and after every base change.

We specify the coefficient calculation needed for independence.
On the split finite surface
$Q=\Spec F[x,y,u]/(xy+P(u))$, with $F=\Q(\mu_{2p^r})$ for
odd $p$ and $F=\Q(\mu_{2^r})$ at two, write the roots as $u_i$,
and put $\eta_i=1-[(x,u-u_i)]$.
All roots lie in the finite open $D(ug_k)$.
Localization and d\'evissage give a supported lift of
$k[g_k]\delta$ to $K_3(Q)$.
The matching-sphere calculation \eqref{d3:eq:cyclotomicendpoint}
identifies any such lift, rationally, as
\begin{equation}\label{d3:eq:wildsupportedcoefficients}
 \widetilde c_k=-n\sum_i b_k(u_i)\eta_i,
\end{equation}
where
\[
 b_k(u)=
 \begin{cases}
 \beta_{\rm BT}(u^k),&p\text{ odd},\\
 \operatorname{Tr}^{K}_{F(\zeta_3)/F}
       \beta_{\rm BT}(\zeta_3u^k),&p=2.
 \end{cases}
\]
For completeness, choose an arc between each pair of roots whose
interior lies in the punctured open unit disc. There
$\operatorname{Log}(1-u^k)$, or the sum
$\operatorname{Log}(1-\zeta_3u^k)+
 \operatorname{Log}(1-\zeta_3^2u^k)$, has a holomorphic branch.
Its primitive for $-L\,du/u$ is respectively
\[
 \frac1k\operatorname{Li}_2(u^k),\qquad
 \frac1k\bigl(\operatorname{Li}_2(\zeta_3u^k)
                  +\operatorname{Li}_2(\zeta_3^2u^k)\bigr).
\]
The factor $k$ in $c_k$ cancels this denominator.
The endpoint periods therefore have the common factor $n$,
independent of $k$. Checking the periods at every complex embedding,
and then applying Borel injectivity as in
Lemma~\ref{d3:lem:explicitoddcoefficients}, determines all coefficient
differences. The remaining diagonal coefficient vanishes because
$\sum_i\eta_i=0$. This proves
\eqref{d3:eq:wildsupportedcoefficients}.

The common normalization of
Lemma~\ref{d3:lem:cyclotomicnormalization}, the root-of-unity formula
\cite[Theorem~1.12]{BesserDeJeu}, and the point comparison give,
with one nonzero scalar for the whole vector, the local coefficient
values
\[
 c_{2,1}(b_k(u_i))=
 \begin{cases}
 \kappa\exp_E(\operatorname{Li}_2(u_i^k)),&p\text{ odd},\\
 \kappa\exp_E(\Psi_2(u_i^k)),&p=2.
 \end{cases}
\]
At two this follows by restriction-transfer and distribution:
$\operatorname{Li}_2(\zeta_3u^k)+
\operatorname{Li}_2(\zeta_3^2u^k)=\Psi_2(u^k)$.
All these roots and their complements used in $\beta_{\rm BT}$
are global units. No root-dependent normalization is introduced.

After completion and the finite free extension to
$\OO_E$, every shifted root lies in its maximal ideal, so
Lemma~\ref{d3:lem:completedrootcup} and the line-bundle Chern formula
apply to \eqref{d3:eq:wildsupportedcoefficients}.
They kill exactly the diagonal coefficient vectors.
Every vector displayed above changes sign when $u_i$ is replaced
by $u_i^{-1}$. A diagonal linear combination must therefore be zero.
Looking at a single fixed primitive root now turns a relation
among the $c_j$ into a relation among the wild-orbit values in
Proposition~\ref{d3:prop:wildorbit}. Its independence forces every
coefficient to vanish, even over $\Q_p$ at the regulator level.
Thus the actual $K_3$ classes are $\Q$-linearly independent after
completion. The arithmetic local ring and its henselization map
to that completion, proving the corresponding assertions there.
\end{proof}

\begin{remark}
The explicit lower bounds are $p^{r-1}$ for odd $p$ and
$2^{r-2}$ at two. They equal the cyclotomic minus dimension at
$p=3$ and $p=2$. For $p\geq5$ this argument proves independence
through the Teichm\"uller component; it does not assert full minus
rank $(p-1)p^{r-1}/2$.
\end{remark}


\section{Henselization and infinite rank}\label{d3:sec:henselian}
Actual local coefficients descend to henselian arithmetic valuation rings.
Their finite families give the same rank bounds before completion and
produce infinite rank after strict henselization.

\subsection{Actual coefficients over a henselian valuation ring}
\label{d3:sec:henselianeisenstein}

The coefficient classes needed by the completed argument can already be
chosen over henselian number-field valuation rings.  The finite-quotient
assertion in Levine's theorem gives this directly.

\begin{lemma}
\label{d3:lem:henseliancoefficients}
Let $W_0$ be a discrete valuation ring whose fraction field is a number
field and whose residue field has characteristic $p$.  Put
\[
H=\Frac(W_0^h),\qquad K=\Frac(\widehat W_0).
\]
For every $\nu\geq1$, the natural map
\[
K_3(W_0^h)\longrightarrow K_3(K)/p^\nu
\]
is surjective.  In particular the rational \'etale Chern classes of
actual elements of $K_3(W_0^h)$ span $H^1(K,\Q_p(2))$ over $\Q_p$.
\end{lemma}

\begin{proof}
The henselization $W_0^h$ is a discrete valuation ring with the same
residue field and completion as $W_0$; see
\cite[Section~15.46, Lemmas~15.46.1, 15.46.3 and 15.46.11]{StacksHenselization}.
Its fraction field embeds densely in $K$, and
\[
H=\{a\in K:a\text{ is algebraic over }\Q\}.
\]
Indeed $H$ is algebraic over $\Frac(W_0)$.  Conversely a henselian
rank-one valued field is separably algebraically closed in its completion:
approximate a separable algebraic root by an element of the field closely
enough to apply Hensel's lemma, and isolate that root from the finitely
many other roots.  The Henselian root then equals the original root.
Characteristic zero makes every algebraic extension separable.

The second assertion of \cite[Theorem~4.13, p.~336]{LevineIndK3},
applied to $K$ and its field of constants $H$, gives
\[
K_3(H)^{\mathrm{ind}}/p^\nu
\xrightarrow{\sim}K_3(K)^{\mathrm{ind}}/p^\nu.
\]
Milnor $K_3(K)$ is $p$-divisible, as used in
Lemma~\ref{d3:lem:actuallocal}.  Its image in $K_3(K)$ is therefore
contained in $p^\nu K_3(K)$, so
$K_3(K)^{\mathrm{ind}}/p^\nu=K_3(K)/p^\nu$.
Consequently $K_3(H)\to K_3(K)/p^\nu$ is surjective.
Quillen localization lifts every element of $K_3(H)$ to $K_3(W_0^h)$,
because $K_2(k_{W_0})=0$.

The proof of Lemma~\ref{d3:lem:actuallocal} identifies
$\varprojlim_\nu K_3(K)/p^\nu$ with $H^1(K,\Z_p(2))$.
The surjectivity just proved makes the image of the actual group
$K_3(W_0^h)$ dense in this inverse limit.  Its rational Chern image
therefore spans $H^1(K,\Q_p(2))$.
\end{proof}

\begin{theorem}
\label{d3:thm:henselianeisensteinrank}
Let $F_0$ be a number field, let $V_0=(\OO_{F_0})_{\mathfrak p}$ for a
prime above $p$, and let $F$ be its completion.  Let $f(z)\in V_0[z]$
be Eisenstein of degree $N\geq2$, and put
\[
R_0=\left(V_0[x,y,z]/(xy+f(z))\right)_{(\mathfrak p,x,y,z)},
\qquad R=R_0^h.
\]
Then $R$ is a henselian regular local ring of dimension three with
residue field $k_F$, and
\[
\dim_{\Q}\ker\left(K_3(R)_{\Q}\longrightarrow
K_3(\Frac R)_{\Q}\right)\geq [F:\Q_p](N-1).
\]
The lower bound is witnessed by actual integral classes whose generic
images vanish integrally.
\end{theorem}

\begin{proof}
Choose a root $\theta$ of $f$ and put
\[
W_0=V_0[\theta],\qquad K=F(\theta),\qquad W=\OO_K.
\]
Eisenstein's condition makes $W_0$ a discrete valuation ring with
uniformizer $\theta$, residue field $k_F$, and completion $W$.
It also makes $R_0$ regular of dimension three, by the same linear-term
calculation as in Theorem~\ref{d3:thm:completedeisensteinrank}.
Regularity, dimension and the residue field are preserved by
henselization, and
\[
\widehat R=\OO_F[[x,y,z]]/(xy+f(z)).
\]
These assertions, including the faithful flatness of $R\to\widehat R$,
follow from \cite[Section~15.46]{StacksHenselization}.

Henselization commutes with the quotient by $x$, so the regular divisor
$D=\Spec(R/(x))$ has coordinate ring
\[
R/(x)\cong \left(W_0[y]_{(\theta,y)}\right)^h.
\]
Here we use \cite[Section~10.156, Lemma~10.156.2]{StacksHenselization}.
The local map $W_0\to R/(x)$ extends to $W_0^h\to R/(x)$.
On completion this map is compatible with
$W\to W[[y]]=\widehat R/(x)$, by the universal property of
henselization.

For $\beta\in K_3(W_0^h)$ define
\[
c^h(\beta)=i_*(\beta|_D)\in K_3(R),
\qquad i:D\hookrightarrow\Spec R.
\]
Localization makes its image in $K_3(R[1/x])$, hence in
$K_3(\Frac R)$, zero integrally.  Flat base change along
$R\to\widehat R$ sends $c^h(\beta)$ to exactly the class $c(\widehat
\beta)$ constructed in Theorem~\ref{d3:thm:completedeisensteinrank},
where $\widehat\beta$ is the image in $K_3(W)$.

By Lemma~\ref{d3:lem:henseliancoefficients}, the values
\[
a(\beta)=\exp_K^{-1}c_{2,1}(\beta|_K)
\]
of actual $\beta\in K_3(W_0^h)$ span $K$ over $\Q_p$.
The completed detector in
Theorem~\ref{d3:thm:completedeisensteinrank} has kernel precisely $F$
on this coefficient space.  Choose $[F:\Q_p](N-1)$ actual coefficient
classes whose values are independent in $K/F$ over $\Q_p$.
Their completed pushforwards have independent detected classes.
Any rational relation among their henselian pushforwards would remain
a relation after completion and would contradict that independence.
This proves the asserted lower bound.
\end{proof}

In particular, for every prime $p$ and every $N\geq2$, the henselization
\[
\left(\left(\Z_{(p)}[x,y,z]/(p+xy+z^N)\right)_{(x,y,z)}\right)^h
\]
has residue field $\F_p$, dimension three, and rational generic
$K_3$-kernel of dimension at least $N-1$.


\subsection{One complete strictly henselian ring with infinite rational rank}

\begin{lemma}\label{d3:lem:completedcupperfect}
Lemma~\ref{d3:lem:completedcup} remains valid when $E$ is a complete
discretely valued field of mixed characteristic $(0,p)$ with perfect
residue field, containing a fixed finite extension of $\Q_p$ as a valued
subfield.
\end{lemma}
\begin{proof}
We check the points in that proof where the base could matter.
The valuation ring $\OO_E$ is still a noetherian complete DVR.
The fixed rational-radius models used in \eqref{d3:eq:fixedradius}
are therefore noetherian henselian models to which
\cite[Theorem~6.2.1]{Achinger} applies, including for $p$-power coefficients.
The order of operations in \eqref{d3:eq:schemetodagger} is unchanged:
finite coefficients at one fixed radius, derived inverse limit,
inversion of $p$, and finally the dagger homotopy colimit.
The two-chart calculation by integration of overconvergent series gives
the same isomorphism
\[
 f^*:H^2_{\mathrm{dR}}(\PP^1_E)
       \xrightarrow{\sim}H^2_{\mathrm{dR}}(U^\dagger).
\]

The standing hypotheses of \cite{ColmezNiziol}, stated at the beginning
of its introduction, allow any complete mixed-characteristic DVR with
perfect residue field. Thus its Theorem~1.3(2) still gives the injective
boundary $\partial_3$ on $U^\dagger$.
For $\PP^1_E$ the boundary is an isomorphism:
$H^2_{\mathrm{HK}}(\PP^1)=F_0h$ with
$\varphi(ah)=p\sigma(a)h$, where
$F_0=\Frac W(k_E)$, while $H^3_{\mathrm{HK}}(\PP^1)=0$.
The eigenspace $\varphi=p^2$ is zero, since
$p\sigma(a)=p^2a$ for $a\ne0$ contradicts
$v_p(\sigma(a))=v_p(a)$.
This argument does not require finite order of $\sigma$.
Proper \'etale and pro-\'etale cohomology agree by
\cite[Section~7.2.4]{ColmezNiziol}.

We also remove the distinction between topological and ordinary
cohomology in the target. The unit quadric has the smooth weak formal
model
$\operatorname{Spwf}\OO_E\langle X,Y,Z\rangle^\dagger/
(XY+Z(Z+1))$, and its ruling extends to the weak formal projective
line: the ideal $(X,Z)$ is already invertible on the integral model.
The natural Hyodo--Kato comparison in
\cite[equation~(5.4) and the following strict quasi-isomorphism,
Propositions~5.5--5.6]{ColmezNiziol}, together with faithful extension
from $F_0$ to $E$, transports the de Rham degree-two isomorphism to
\[
 H^2_{\mathrm{HK}}(\PP^1_E)
       \xrightarrow{\sim}H^2_{\mathrm{HK}}(U^\dagger);
 \qquad H^3_{\mathrm{HK}}(U^\dagger)=0.
\]
Thus the same Hyodo--Kato term vanishes for $U^\dagger$, and
$\partial_{3,U}$ is an isomorphism.
Its de Rham source is classical by
\cite[Section~5.1]{ColmezNiziol}, even if $E$ has infinite dimension
over $\Q_p$. Hence its target is classical too.
The classical-part functor of \cite[Section~3.1.1]{ColmezNiziol}
sends $\widetilde H^i$ to ordinary $H^i$; applying it to this
isomorphism gives an isomorphism on ordinary cohomology.
Naturality now proves injectivity of pullback from
$H^3_{\mathrm{cont},\mathrm{\acute et}}(\PP^1_E,\Q_p(3))$ on the
ordinary cohomology used by the Chern maps.

For the last step, the splitting \eqref{d3:eq:finitePonebundle}
with $r=3$ gives the split injection
\[
 H^1(E,\Q_p(2))
 \xrightarrow{\ \smile c_1(\OO(1))\ }
 H^3_{\mathrm{cont},\mathrm{\acute et}}(\PP^1_E,\Q_p(3)).
\]
Its finite-coefficient splitting survives derived inverse limits and
inversion of $p$. No vanishing assertion about
$H^3(E,\Q_p(3))$ is required.
Its pullback is the ruling-cup map, which is therefore
injective.
\end{proof}

\begin{theorem}
\label{d3:thm:stricthenselianinfinite}
Set
\[
 V_\infty=W(\overline{\mathbf F}_3),\qquad
 A_\infty=V_\infty[[x,y,z]]/(3+xy+z^2).
\]
Then $A_\infty$ is a complete strictly henselian regular local ring of
dimension three, and
\[
 \dim_{\Q}\ker\bigl(
   K_3(A_\infty)_{\Q}
       \longrightarrow K_3(\Frac A_\infty)_{\Q}\bigr)=\infty.
\]
For every $m\geq1$ this kernel contains $m$ independent classes
represented by actual integral elements of $K_3(A_\infty)$.
\end{theorem}
\begin{proof}
The Witt ring $V_\infty$ is a complete DVR with uniformizer $3$.
The equation has nonzero linear term $3$ in the cotangent space of
$V_\infty[[x,y,z]]$, so its quotient is regular local of dimension three,
with maximal ideal $(x,y,z)$. Completeness and the algebraically closed
residue field imply strict henselianity.

For each positive integer $m$, put
\[
 V_m=W(\mathbf F_{3^m}),\quad F_m=\Frac V_m,\quad
 A_m=V_m[[x,y,z]]/(3+xy+z^2).
\]
The canonical local map $A_m\to A_\infty$ is flat by miracle flatness:
both rings are regular of dimension three and the closed fiber has
dimension zero \cite[Tag~00R4]{StacksFlatness}.

Choose $u^2=-3$ and write
\[
 E_m=F_m(u),\quad W_m=V_m[u],\qquad
 E_\infty=(\Frac V_\infty)(u),\quad W_\infty=V_\infty[u].
\]
The field $E_\infty$ is the completion of $E_m^{\mathrm{ur}}$ for every
$m$. Indeed its Witt coefficients can be approximated to every finite
precision in finite residue subfields. Completion of a henselian
discretely valued field preserves its category of finite separable
extensions \cite[Lemma~4.2.1]{SaitoCompletion}; consequently
\[
 G_{E_\infty}\simeq I_{E_m}.
\]
Restriction therefore induces an injection
\begin{equation}\label{d3:eq:inertiarestriction}
 H^1(E_m,\Q_3(2))\hookrightarrow H^1(E_\infty,\Q_3(2)).
\end{equation}
Here is a direct verification. For $G=G_{E_m}$ and $I=I_{E_m}$,
$\Q_3(2)^I=0$ because the cyclotomic character squared is nontrivial
on inertia. A cocycle trivial in $H^1(I,\Q_3(2))$ can be changed by a
coboundary to vanish on $I$. For such a cocycle $b$,
$i b(g)=b(ig)=b(g(g^{-1}ig))=b(g)$, so every $b(g)$ belongs to
$\Q_3(2)^I$ and is zero.

By Lemma~\ref{d3:lem:actuallocal}, choose actual integral
anti-invariant classes
$\beta_{m,1},\ldots,\beta_{m,m}\in K_3(W_m)$ whose Chern classes
\[
 \alpha_{m,j}=c_{2,1}(\beta_{m,j})
       \in H^1(E_m,\Q_3(2))
\]
are $\Q_3$-independent. Pull these classes to the regular divisor
$A_m/(x)\simeq W_m[[y]]$ and push forward to obtain
$c_{m,j}\in K_3(A_m)$. Localization makes them zero after inverting
$x$. Their images $c^\infty_{m,j}\in K_3(A_\infty)$ are therefore
also generically zero.

Make the finite free base change $V_\infty\to W_\infty$ and put
\[
 T_\infty=\Spec
 \bigl(W_\infty[[x,y,z]]/(xy+z^2-u^2)\bigr)[1/u],
 \qquad I_+=(x,z-u),\qquad \eta_+=1-[I_+].
\]
Flat base change and the projection formula give
\[
 c^\infty_{m,j}|_{T_\infty}
      =2\,\beta_{m,j}|_{E_\infty}\,\eta_+.
\]
All rings here are algebras over the fixed finite extension
$\Q_3(u)$, so Lemma~\ref{d3:lem:continuouschern} applies without any
new finiteness assumption over $E_\infty$.
Using \eqref{d3:eq:etaleproduct} gives
\[
 c_{3,3}(c^\infty_{m,j}|_{T_\infty})
   =4\,c_1(I_+)\smile
       \operatorname{res}_{E_\infty/E_m}(\alpha_{m,j}).
\]
The coefficient classes remain $\Q_3$-independent by
\eqref{d3:eq:inertiarestriction}. Lemma~\ref{d3:lem:completedcupperfect},
with $w=z-u$ and $d=2u$, proves independence of these Chern images.
Hence the actual integral classes $c^\infty_{m,1},\ldots,
c^\infty_{m,m}$ are $\Q$-independent after rationalization.
Since $m$ is arbitrary and $A_\infty$ is fixed, the asserted kernel
is infinite dimensional.
\end{proof}

\begin{remark}
This proof uses arbitrarily large finite families from finite
unramified coefficient fields. It neither requires that actual
regulators exhaust $E_\infty$ nor identifies $A_\infty$ with a
filtered colimit of the rings $A_m$.
\end{remark}


\subsection{Infinite rank already in an arithmetic strict henselization}
\label{d3:sec:stricthenseldescent}

\begin{lemma}
\label{d3:lem:stricthenselcompletion}
For a prime $p$, set
\[
 A_{0,p}=\left(\Z_{(p)}[x,y,z]/(p+xy+z^2)\right)_{(x,y,z)},
 \qquad S_p=A_{0,p}^{sh},
\]
using the geometric residue field $\overline{\F}_p$.
There is a canonical compatible local map
\[
 S_p\longrightarrow
 A_{\infty,p}\coloneqq W(\overline{\F}_p)[[x,y,z]]/(p+xy+z^2),
\]
and it identifies $A_{\infty,p}$ with the completion of $S_p$.
In particular $S_p$ is noetherian regular local of dimension three
with residue field $\overline{\F}_p$.
\end{lemma}
\begin{proof}
The target is a complete local ring with algebraically closed residue
field, hence strictly henselian.  The local map from $A_{0,p}$ extends
to $S_p$ by the universal property with the prescribed geometric
residue field.  The assertions about $S_p$ follow from
\cite[Section~15.46]{StacksHenselization}.

Write $\m=(x,y,z)$.  Fix $r\geq1$ and choose $s$ with $2s\geq r$.
Since $p\in\m^2$, the Artinian local ring
$R_r=A_{0,p}/\m^r$ is a $\Z/p^s$-algebra, and
\[
 A_{\infty,p}/\m^r
 \cong R_r\otimes_{\Z/p^s}W_s(\overline{\F}_p).
\]
Here $W_s$ denotes Witt vectors of length $s$.  The equality
\[
 W_s(\overline{\F}_p)=\colim_m W_s(\F_{p^m}),
\]
with $m$ ordered by divisibility, follows because a finite Witt vector
has only finitely many residue-field coordinates.  Each
$R_r\otimes_{\Z/p^s}W_s(\F_{p^m})$ is the finite \'etale local extension
of $R_r$ with residue field $\F_{p^m}$.
Finite \'etale extensions of a henselian local ring are determined by
their residue extensions \cite[Tag~04GK]{StacksHenselization}.
Thus $A_{\infty,p}/\m^r$ is the strict henselization of $R_r$.
By quotient compatibility
\cite[Section~10.156, Lemma~10.156.4]{StacksHenselization},
\[
 A_{\infty,p}/\m^r\cong S_p/\m^rS_p.
\]
These identifications respect the transition maps and the chosen
residue field.  Taking inverse limits proves
$\widehat S_p\cong A_{\infty,p}$.
\end{proof}

\begin{corollary}
\label{d3:cor:stricthenselizationinfinite}
Let
\[
 S=\left(\left(\Z_{(3)}[x,y,z]/(3+xy+z^2)\right)_{(x,y,z)}\right)^{sh}.
\]
Then
\[
 \dim_{\Q}\ker\left(K_3(S)_{\Q}
      \longrightarrow K_3(\Frac S)_{\Q}\right)=\infty.
\]
For every $m\geq1$ there are $m$ independent classes represented by
actual integral elements of $K_3(S)$ whose generic images vanish
integrally.
\end{corollary}
\begin{proof}
For each $m$, choose a monic polynomial $g_m\in\Z[T]$ whose reduction
modulo $3$ is irreducible of degree $m$, and choose a residue root in
$\F_{3^m}\subset\overline{\F}_3$.
The ring
\[
 V_{m,0}=\Z_{(3)}[T]/(g_m)
\]
is a finite \'etale local discrete valuation ring with number-field
fraction field, residue field $\F_{3^m}$, and completion
$V_m=W(\F_{3^m})$.  The chosen residue root gives compatible embeddings
of $V_{m,0}$ into $\Z_{(3)}^{sh}$ and into $V_m$.
Put
\[
 R_{m,0}=A_{0,3}\otimes_{\Z_{(3)}}V_{m,0},\qquad R_m=R_{m,0}^h.
\]
The finite \'etale algebra $R_{m,0}$ is local, because its closed fiber
is the field $\F_{3^m}$.  There are compatible local maps
\[
 R_m\longrightarrow S\longrightarrow A_{\infty,3},
 \qquad
 \widehat R_m=A_m=V_m[[x,y,z]]/(3+xy+z^2)
       \longrightarrow A_{\infty,3}.
\]
The two maps from $R_m$ to $A_{\infty,3}$ agree by the universal
property of henselization.  The second row is the usual inclusion of
Witt coefficients after completion.

Set $u^2=-3$ and $W_{m,0}=V_{m,0}[u]$.
Lemma~\ref{d3:lem:henseliancoefficients}, followed by $1-\sigma$ for
$\sigma(u)=-u$, supplies actual anti-invariant classes
\[
 \beta_{m,1}^h,\ldots,\beta_{m,m}^h\in K_3(W_{m,0}^h)
\]
whose Chern classes in $H^1(\Frac(V_m)(u),\Q_3(2))$ are
$\Q_3$-independent.
Pull them to $R_m/(x)$ and push forward along $x=0$ as in
Theorem~\ref{d3:thm:henselianeisensteinrank}.  Their resulting classes in
$K_3(R_m)$ vanish after inverting $x$, and so do their images in
$K_3(S)$.

After mapping to $A_{\infty,3}$ these are exactly the completed
finite-level classes detected in
Theorem~\ref{d3:thm:stricthenselianinfinite}.  Its proof applies to
arbitrary independent anti-invariant coefficient classes from each
finite unramified level.  It therefore proves that these $m$ images
are $\Q$-independent.  Any rational relation in $K_3(S)$ would remain
a relation in $K_3(A_{\infty,3})$, so the classes in $K_3(S)$ are
independent as well.  Since $m$ is arbitrary, the result follows.
\end{proof}


\section{Cyclotomic coefficients in higher odd degrees}
\label{d3:sec:completedhigherodd}

The following construction supplies actual higher $K$-classes without
requiring a general surjectivity assertion from actual $K$-theory to
its $p$-completed realization. Fix a weight $n\geq2$. We write
$c_{n,1}:K_{2n-1}\to H^1(-,\Q_p(n))$ and
$c_{n+1,3}:K_{2n-1}\to H^3(-,\Q_p(n+1))$; the second subscript
denotes the cohomological degree.

\begin{lemma}
\label{d3:lem:highercoefficients}
Let $E/\Q_p$ be finite and $n\geq2$.
For a root of unity $\rho\in E$, the number-field cyclotomic class
$b_n(\rho)\in K_{2n-1}(\OO_E)_{\Q}$ has
\begin{equation}\label{d3:eq:highercyclotomicregulator}
 c_{n,1}(b_n(\rho))
 =\exp_{E,n}\bigl(\kappa_n\operatorname{Li}_n(\rho)\bigr),
 \qquad \kappa_n\in\Q^\times,
\end{equation}
where $\kappa_n$ is common to all roots and embeddings under one
fixed symbol convention. The exponential is an isomorphism for
$n\geq2$. For an invertible module $L$ and an actual class $\beta$
of degree $2n-1$, one has
\begin{equation}\label{d3:eq:higherlineproduct}
 c_{n+1,3}\bigl(\beta(1-[L])\bigr)
 =n\,c_{n,1}(\beta)\cup c_1(L).
\end{equation}
Finally, Lemmas~\ref{d3:lem:completedcup} and
\ref{d3:lem:completedrootcup} hold with coefficient group
$H^1(E,\Q_p(n))$ and target twist $n+1$.
\end{lemma}
\begin{proof}
Theorem~1.12 of \cite{BesserDeJeu} constructs the actual rational
cyclotomic class and includes roots of $p$-power order.
Its normalized syntomic value is
$\pm(n-1)!\operatorname{Li}_n(\rho)$.
To compare the normalizations, put $q=p^f=|k_E|$, let $E_0$ be the
maximal unramified subfield of $E$, and write $\sigma$ for its Frobenius.
Set
\[
 A=1-p^{-n}\sigma,\qquad
 N=\sum_{j=0}^{f-1}p^{-nj}\sigma^j,\qquad P=1-q^{-n}.
\]
Then $NA=AN=P$.  The point syntomic complex of
\cite[Example~2.2.4]{HuberKings} is
$[E_0\longrightarrow E\oplus E_0]$, with
$c\mapsto(c,Ac)$, and its canonical $E$-coordinate is
$\eta([a,b])=a-A^{-1}b$.
The comparison of \cite[Proposition~8.6(2)--(3), its proof and equation~(8.1)]{BesserRegulators}
uses $N$ to pass from the semilinear $p$-Frobenius to the linear
$q$-Frobenius.  On the point the latter acts as the identity on $E$,
and the raw modified-cone coordinate is
\[
 \epsilon=Pa-Nb=P\eta([a,b]).
\]
Definition~4.6 of \cite{BesserDeJeu} sends this coordinate to
\begin{equation}\label{d3:eq:pointnormalization}
 (1-\varphi_q^*/q^n)^{-1}\epsilon
   =P^{-1}\epsilon=\eta([a,b]).
\end{equation}
Thus \cite[Corollary~9.10 and Proposition~9.11]{BesserRegulators},
or equivalently \cite[Proposition~2.3.4]{HuberKings}, identifies the
normalized regulator with the Chern-compatible Bloch--Kato
exponential.  This proves \eqref{d3:eq:highercyclotomicregulator};
the only remaining constant is the fixed rational symbol normalization
in Theorem~1.12 and Remark~1.7 of \cite{BesserDeJeu}.
The calculation applies to arbitrary ramification and to $p=2$.

Formula~(3.10) and Theorem~3.11(i) of \cite{WeibelChern}, applied
to $[L]-1$ in degree zero and $\beta$ in degree $2n-1$, have only
the term $i_1=1,i_2=n$ in weight $n+1$.
Its coefficient is $-n$; negating $[L]-1$ gives
\eqref{d3:eq:higherlineproduct}. The exceptional $K$-degree two
case does not occur. The finite-type descent proof of
Lemma~\ref{d3:lem:continuouschern} applies unchanged in degree
$2n-1\geq3$, so the identity holds in continuous cohomology.

For the cup assertion use the same fixed-radius map
\eqref{d3:eq:schemetodagger} and the same ruling morphism
$f:U^\dagger\to\PP^1_E$. The boundary with twist $n+1$ in
\cite[Theorem~1.3]{ColmezNiziol} is injective in degree three
when $n=2$ and is an isomorphism when $n\geq3$, since
$3\leq(n+1)-1$. On $\PP^1_E$ the boundary is an isomorphism
also for $n=2$, as in Lemma~\ref{d3:lem:completedcup}.
Its naturality and the de Rham ruling isomorphism make
$f^*$ injective on $H^3_{\mathrm{pro\acute et}}(-,\Q_p(n+1))$.
The splitting \eqref{d3:eq:finitePonebundle}, with $r=n+1$,
identifies cup product with $c_1(\OO(1))$ as a split injection
from $H^1(E,\Q_p(n))$ into its source.
Composing with the injective $f^*$ proves the quadratic assertion.
The pairwise formal maps in Lemma~\ref{d3:lem:completedrootcup}
then give the simultaneous assertion with the same diagonal kernel.
\end{proof}

\begin{lemma}
\label{d3:lem:evenweightibase}
For every even integer $n\geq2$ and $i^2=-1$,
\[
 \frac{\operatorname{Li}_n(i)}{i}\in\frac12+\Z_2.
\]
In particular $v_2(\operatorname{Li}_n(i))=-1$.
\end{lemma}
\begin{proof}
Choose a primitive cube root $\zeta$ in an unramified quadratic
extension and put
$D_n=\operatorname{Li}_n(\zeta i)+\operatorname{Li}_n(\zeta^2i)$.
Write $F_{2,n}(z)=\operatorname{Li}_n(z)-2^{-n}\operatorname{Li}_n(z^2)$.
Distribution and inversion give
\[
 D_n=-(1+3^{1-n})\operatorname{Li}_n(i),\qquad
 \operatorname{Li}_n(-\zeta)+\operatorname{Li}_n(-\zeta^2)=0.
\]
Thus $D_n=F_{2,n}(\zeta i)+F_{2,n}(\zeta^2i)$.
For $S_j(t)=\sum_{m\geq0}m^jt^m
=P_j(t)/(1-t)^{j+1}$, expand the odd denominators by the binomial
series first on $|z|<1$. The terms converge uniformly when
$|z|\leq1$ and $v_2(1-z^2)\leq2/3$, with valuation at least
$j-(j+1)2/3$. The rigid identity principle and
\cite[Theorem~2.7(3)]{BesserDeJeuAlgorithm} therefore give
\[
 \frac{D_n}{i}
 =\sum_{j\geq0}(-1)^j\binom{n+j-1}{j}2^j A_j,
 \qquad A_j=\sum_{a=1}^2\zeta^a S_j(-\zeta^{2a}).
\]
It applies because $1+\zeta^{2a}$ is a unit.
Every $A_j$ belongs to $\Z[\zeta]$ and is fixed by
$\zeta\mapsto\zeta^2$, hence lies in $\Z$.
Directly, $A_0=-2$ and $A_1=1$.
Consequently
\[
 D_n/i\in-2-2n+4\Z_2=-2+4\Z_2.
\]
Since $n-1$ is odd, $3^{n-1}\equiv3\pmod8$ and
$v_2(1+3^{1-n})=2$.
The first displayed identity now gives
$\operatorname{Li}_n(i)/i\in\Q_2$ of valuation $-1$,
which is exactly the asserted coset.
\end{proof}

\begin{proposition}
\label{d3:prop:higherwildaugmentation}
Let $r\geq3$, let $\rho_r$ be a primitive $2^r$-th root with
$\rho_r^2=\rho_{r-1}$, and put
$E_r=\Q_2(\rho_r)$ and $D=2^{r-2}$.
For every $n\geq2$,
\begin{equation}\label{d3:eq:higherwilddepletion}
 v_2\left(\operatorname{Li}_n(\rho_r)
       -2^{-n}\operatorname{Li}_n(\rho_{r-1})\right)
 =-2^{-(r-2)}.
\end{equation}
The $D-1$ elements
\[
 \operatorname{Li}_n(\rho_r^{5^{j+1}})
 -\operatorname{Li}_n(\rho_r^{5^j}),
 \qquad 0\leq j\leq D-2,
\]
are $\Q_2$-linearly independent, and their span intersects
$\Q_2\subset E_r$ only in zero.
If $n$ is even, the $D$ elements
\[
 \operatorname{Li}_n(\rho_r^{5^j}),\qquad 0\leq j\leq D-1,
\]
form a $\Q_2$-basis of the inversion-odd subspace $E_r^-$.
\end{proposition}
\begin{proof}
Put $F_{2,n}(z)=\operatorname{Li}_n(z)-2^{-n}\operatorname{Li}_n(z^2)$.
With $S_j(t)=P_j(t)/(1-t)^{j+1}$ as above, expansion on $|z|<1$ gives
\[
 F_{2,n}(z)=z\sum_{j\geq0}(-1)^j
       \binom{n+j-1}{j}2^j S_j(z^2).
\]
The sum converges uniformly on
$|z|\leq1$, $v_2(1-z^2)\leq2/3$, since its $j$th term has
valuation at least $j-(j+1)2/3\to\infty$.
This is a connected annulus inside $|1-z|>1/2$;
\cite[Theorem~2.7(3)]{BesserDeJeuAlgorithm} and the rigid identity
principle extend the equality there. At $\rho_r$ put
$d=v_2(1-\rho_r^2)=2^{-(r-2)}\leq1/2$.
The $j=0$ term has valuation $-d$, while every $j\geq1$ term
has valuation at least $j-(j+1)d>-d$.
This proves \eqref{d3:eq:higherwilddepletion}.

Write $R_s=\operatorname{Li}_n(\rho_s)$ and
$\Gamma_r=\Gal(E_r/\Q_2(i))=\langle\gamma\rangle$, with
$\gamma(\rho_r)=\rho_r^5$ and $|\Gamma_r|=D$.
Inversion acts on $R_r$ by $\epsilon=(-1)^{n+1}$.
The normal-basis theorem makes $E_r^\epsilon$ the regular
$\Q_2[\Gamma_r]$-module \cite[Chapter~I, p.~34]{MilneDuality}. Distribution gives average trace
$R_s\mapsto2^{-n}R_{s-1}$.
For every $3\leq s\leq r$, the rational conductor
$2^{s-2}$ projection of $R_r$ is therefore
\[
 2^{-n(r-s)}(R_s-2^{-n}R_{s-1})\ne0.
\]
Each nontrivial rational irreducible of the cyclic group occurs
once, and its polynomial $\Phi_{2^a}$ is irreducible over $\Q_2$.
Hence the augmentation ideal acts injectively on $R_r$.
Its basis $\gamma^{j+1}-\gamma^j$, $0\leq j\leq D-2$, gives
the asserted independent values. Every such value has
$\Gamma_r$-average zero, so their span meets $\Q_2$ only in zero.
If $n$ is even, the trivial $\Gamma_r$ projection is
$2^{-n(r-2)}\operatorname{Li}_n(i)\ne0$ by
Lemma~\ref{d3:lem:evenweightibase}. All rational blocks then occur,
so the whole orbit is a basis of $E_r^-$.
\end{proof}

\begin{theorem}
\label{d3:thm:completedhigherodd}
For $r\geq3$ put $D=2^{r-2}$ and
\[
 A_r=\Z_2[[x,y,z]]/
 \bigl(xy+\Phi_{2^r}(1+z)\bigr).
\]
For every $n\geq2$,
\[
 \dim_{\Q}\ker\left(K_{2n-1}(A_r)_{\Q}
       \longrightarrow K_{2n-1}(\Frac A_r)_{\Q}\right)
 \geq
 \begin{cases}
 D,& n\text{ even},\\
 D-1,& n\text{ odd}.
 \end{cases}
\]
More precisely, let $W_r=\Z_2[\rho_r]$, identify
$A_r/(x)=W_r[[y]]$, and let $i$ be this closed immersion.
For even $n$, use the $D$ classes
\[
 i_*b_n(\rho_r^{5^j}),\qquad 0\leq j\leq D-1.
\]
For odd $n$, use the $D-1$ classes
\[
 i_*\bigl(b_n(\rho_r^{5^{j+1}})-b_n(\rho_r^{5^j})\bigr),
 \qquad 0\leq j\leq D-2.
\]
The coefficients are first pulled to $W_r[[y]]$.
The indicated classes are independent in the displayed kernel.
\end{theorem}
\begin{proof}
The polynomial $\Phi_{2^r}(1+z)$ is Eisenstein at two.
Its linear uniformizer term makes $A_r$ regular local of
dimension three, with residue field $\F_2$.
Localization makes each supported class generically zero.
After the finite free coefficient extension from $\Z_2$ to $W_r$,
the roots are $a_\sigma=\sigma(\rho_r)-1$.
On the regular generic open, flat base change and projection give
\[
 i_*(\beta)|_T
 =\sum_{\sigma\in\Gal(E_r/\Q_2)}
      \sigma(\beta)(1-[J_\sigma]).
\]
Base change may be performed in $G$-theory before restricting
to this regular open, so regularity of the split integral model
is not required. By Lemma~\ref{d3:lem:highercoefficients}, its Chern
class is
\[
 n\sum_\sigma
 \exp_{E_r,n}(\kappa_n\sigma(R))\cup c_1(J_\sigma),
\]
where $R$ is the corresponding difference in
Proposition~\ref{d3:prop:higherwildaugmentation}, or the individual
value there when $n$ is even.
The simultaneous cup detector kills a coefficient vector exactly
when it is diagonal. Since the exponential is injective, a
relation among these detectors would put the corresponding
$\Q_2$-linear combination of the $R$'s in $\Q_2$.
Proposition~\ref{d3:prop:higherwildaugmentation} makes that combination
zero (using $E_r^-\cap\Q_2=0$ in even weight), and then makes
all its coefficients zero.
The detectors are $\Q_2$-independent, hence the actual rational
$K$-classes are $\Q$-independent. A common integer clears the
denominators of this finite family and gives integral classes
whose generic images are zero integrally.
\end{proof}

For odd weight the trivial wild-character block lies in the
constant field and is deliberately removed by taking differences.
No assertion about surjectivity of the actual local higher
$K$-regulator onto its completed target is needed.


\begin{corollary}
For the single regular local ring
\[
 R=\Z_2[[x,y,z]]/
 \bigl(xy+z^4+4z^3+6z^2+4z+2\bigr),
\]
the generic restriction has nonzero rational kernel in every
odd degree at least three. In degree $4j-1$ its kernel has
$\Q$-dimension at least two, and in degree $4j+1$, $j\geq1$,
at least one.
\end{corollary}
\begin{proof}
Take $r=3$ in Theorem~\ref{d3:thm:completedhigherodd} and use
$\Phi_8(1+z)=(1+z)^4+1$.
\end{proof}


\subsection{Higher polylogarithms on odd-prime wild orbits}
\label{d3:sec:oddprimehigherwild}

Fix an odd prime $p$, put $M=p-1$, and choose compatible primitive
$p^r$-th roots $\rho_r$. Write
\[
 E_r=\Q_p(\rho_r),\qquad e_r=Mp^{r-1},\qquad
 R_r=\operatorname{Li}_n(\rho_r),\qquad n\geq2.
\]
Let $\Delta$ be the tame subgroup and $\omega:\Delta\to\mu_M$
its Teichm\"uller character. Put $\chi_n=\omega^{n-1}$ and
$\Gamma_r=\Gal(E_r/E_1)$, generated by
$\gamma(\rho_r)=\rho_r^{1+p}$.

\begin{proposition}
\label{d3:prop:oddprimehigherwild}
For even $n$, let $\ell\in\{1,3,\ldots,p-2\}$ satisfy
$\ell\equiv n-1\pmod M$. Then
\begin{equation}\label{d3:eq:oddprimehigherbase}
 v_p(R_1)=\frac{\ell}{p-1}-1-v_p(n-1).
\end{equation}
For either parity of $n$, let $q\in\{0,\ldots,p-2\}$ satisfy
$q\equiv n\pmod M$. For $r\geq2$,
\begin{equation}\label{d3:eq:oddprimehighernew}
 v_p(R_r-p^{-n}R_{r-1})=\frac{q-p}{e_r}.
\end{equation}
The $\chi_n$-projection preserves the leading term in both formulas.
Writing $D=p^{r-1}$, the following families are $\Q_p$-linearly
independent:
\[
 \{\gamma^jR_r:0\leq j<D\}\quad(n\text{ even}),
\]
and
\[
 \{\gamma^{j+1}R_r-\gamma^jR_r:0\leq j<D-1\}
                                  \quad(n\text{ odd}).
\]
The span of either displayed family has zero intersection with
$\Q_p\subset E_r$.
\end{proposition}

\begin{proof}
Set $F_{p,n}(z)=\operatorname{Li}_n(z)-p^{-n}
\operatorname{Li}_n(z^p)$ and
$S_j(t)=P_j(t)/(1-t)^{j+1}$, $P_j\in\Z[t]$, where
$S_j(t)=\sum_{h\geq0}h^jt^h$ and $S_0(t)=1/(1-t)$.
Grouping the prime-to-$p$ denominators as $a+ph$ gives
\begin{equation}\label{d3:eq:oddprimehigherdepleted}
 F_{p,n}(z)=\sum_{j\geq0}(-1)^j\binom{n+j-1}{j}p^j
       \sum_{a=1}^{p-1}\frac{z^a}{a^{n+j}}S_j(z^p).
\end{equation}
On $|z|\leq1$, $v_p(1-z^p)\leq b<1$, the $j$th term has
valuation at least $j-(j+1)b$, so the series converges uniformly.
This domain is the annulus
$p^{-b/p}\leq|1-z|\leq1$ and lies in the continuation region
of \cite[Theorem~2.7(3)]{BesserDeJeuAlgorithm}.
Both sides agree on $|z|<1$, hence on the annulus. Taking
$b=2/3$ includes all points evaluated below.

Suppose first that $n$ is even, and put $\rho=\rho_1$.
Distribution and inversion
\cite[Proposition~2.10]{BesserDeJeuAlgorithm} give
\[
 \sum_{\xi\in\mu_M\setminus\{1\}}\operatorname{Li}_n(\xi\rho)
       =-(1+M^{1-n})\operatorname{Li}_n(\rho).
\]
Since $(\xi\rho)^p=\xi$ and
$\sum_{\xi\ne1}\operatorname{Li}_n(\xi)=0$ by inversion,
the sum on the left can be replaced by
$\sum_{\xi\ne1}F_{p,n}(\xi\rho)$.
All terms of \eqref{d3:eq:oddprimehigherdepleted} with $j\geq1$
then belong to $p\mathcal O_{E_1}$. Its $j=0$ sum is
\[
 P_n(\rho)=\sum_{a=1}^{p-1}\frac{C_a\rho^a}{a^n},\qquad
 C_a=\sum_{\xi\ne1}\frac{\xi^a}{1-\xi}=a-\frac p2.
\]
Indeed, pairing inverse roots gives
$\sum_{\xi\ne1}(1-\xi)^{-1}=(M-1)/2$; hence
$C_1=1-p/2$, and $C_{a+1}-C_a=1$ for $1\leq a<M$.

For any integer $b$, the coefficient of $(Z-1)^k$ in
$\sum_{a=1}^{p-1}a^bZ^a\in\F_p[Z]$ is
$\sum_a a^b\binom ak$.
If $q$ is the least nonnegative residue of $-b$ modulo $M$,
the first nonzero coefficient is at $k=q$ and equals $-1/q!$:
for $k<q$ no power occurring in the degree-$k$ polynomial
$a^b\binom ak$ has exponent divisible by $M$, while at $q$
only its leading coefficient contributes.
Applying this with $b=1-n$ shows that $\overline P_n$
has a zero of exact order $\ell$ at $1$. Thus
\[
 v_pP_n(\rho)=\ell/M<1.
\]
The lower Taylor coefficients are divisible by $p$, and the
higher terms have strictly greater valuation. They and the
$j\geq1$ remainder cannot cancel this term.
Since $n-1$ is odd, the elementary lifting-the-exponent identity gives
\[
 v_p(1+M^{1-n})=1+v_p(n-1),
\]
proving \eqref{d3:eq:oddprimehigherbase}.

For $r\geq2$ and either parity, the $j=0$ term in
\eqref{d3:eq:oddprimehigherdepleted} is
\[
 \frac{\sum_{a=1}^{p-1}a^{-n}\rho_r^a}{1-\rho_r^p}.
\]
The same Taylor calculation gives numerator valuation $q/e_r$,
including $q=0$. The denominator has valuation $p/e_r$.
For $j\geq1$ the valuation is at least
\[
 j-(j+1)p/e_r\ \geq\ j-(j+1)/(p-1)\ \geq\ 0,
\]
and this bound tends to infinity. The $j=0$ term has negative
valuation, so it is the unique leading contribution. This proves
\eqref{d3:eq:oddprimehighernew}.

It remains to verify the character assertion, rather than merely
nonvanishing of the total values. Put $\pi_r=\rho_r-1$.
If $\tau\in\Delta$ has $\overline{\omega(\tau)}=a$, then
$\tau(\pi_r)/\pi_r\equiv a\pmod{\pi_r}$.
For $0\ne x\in E_r$ with $v_{\pi_r}(x)=m$, the leading term
therefore transforms by $a^m$. The idempotent
\[
 e_{\chi_n}=\frac1M\sum_{\tau\in\Delta}\chi_n(\tau)^{-1}\tau
\]
preserves it if $m\equiv n-1\pmod M$.
For the base value this congruence follows from
$m=\ell-M(1+v_p(n-1))$; for the wild difference it follows
from $m=q-p$.

The normal basis theorem identifies $E_r^{\chi_n}$ with the
regular $\Q_p[\Gamma_r]$-module
\cite[Chapter~I, p.~34]{MilneDuality}. Its rational character blocks
occur once, with polynomials $\Phi_{p^j}$, $0\leq j\leq r-1$.
If $H_r=\Gal(E_r/E_{r-1})$, distribution gives
\[
 \frac1p\sum_{h\in H_r}hR_r=p^{-n}R_{r-1}.
\]
Thus \eqref{d3:eq:oddprimehighernew} supplies a nonzero projection
to the new block $\Phi_{p^{r-1}}$, and the averaging identity
preserves every previously nonzero block.
For even $n$, \eqref{d3:eq:oddprimehigherbase} starts this induction
with a nonzero trivial wild block, giving a cyclic vector of
the full regular module. For odd $n$, all nontrivial blocks
occur, so the augmentation ideal maps injectively.
Applying $e_{\chi_n}$ to a relation now proves the two asserted
independence statements for the unprojected values themselves.
Finally, for even $n$ inversion is odd, excluding $\Q_p$;
for odd $n$ every displayed difference has zero $\Gamma_r$-average,
which likewise excludes $\Q_p$.
\end{proof}

\begin{remark}
The tame idempotent is used only on the regulator target.
Actual coefficients may remain the root classes
$b_n(\rho_r^{(1+p)^j})$ or their adjacent differences; no
$\Q_p$-linear projector is applied inside actual rational
$K$-theory. With the common normalization
\eqref{d3:eq:highercyclotomicregulator}, their usual local Chern
values are the corresponding displayed values followed by
multiplication by $\kappa_n\ne0$ and the Bloch--Kato exponential.
Since these operations are $\Q_p$-linear isomorphisms, the
actual root family has $\Q$-span of dimension $D$ in even
weight, and the actual difference family has $\Q$-span of
dimension $D-1$ in odd weight: any rational relation would
give a forbidden $\Q_p$-linear relation of regulator values.
The proposition detects one tame character and
does not assert nonvanishing of all tame character components.
\end{remark}

\begin{corollary}
\label{d3:cor:oddprimecompletedhigher}
For an odd prime $p$ and $r\geq2$ put $D=p^{r-1}$ and
\[
 A_{p,r}=\Z_p[[x,y,z]]/(xy+\Phi_{p^r}(1+z)).
\]
For every $n\geq2$, the rational generic kernel in degree $2n-1$
has dimension at least $D$ when $n$ is even and at least $D-1$
when $n$ is odd. The classes are the pushforwards from $x=0$ of
$b_n(\rho_r^{(1+p)^j})$ in even weight, or their consecutive
differences in odd weight, after one common denominator is cleared.
They vanish after inverting $x$ integrally.
In particular, fixing $r=2$ gives one complete regular local ring
with residue field $\F_p$ for which rational generic injectivity
fails in every odd degree at least three.
\end{corollary}
\begin{proof}
The polynomial is Eisenstein and the equation has nonzero linear
term $p$, so the displayed ring is regular local of dimension three.
Apply the proof of Theorem~\ref{d3:thm:completedhigherodd}, with
Proposition~\ref{d3:prop:oddprimehigherwild} in place of
Proposition~\ref{d3:prop:higherwildaugmentation}.
The simultaneous ruling-cup detector in
Lemma~\ref{d3:lem:highercoefficients} kills exactly the diagonal
coefficient vectors. The regulator spans just computed meet the
constant field only in zero, so the same proof gives precisely
these ranks and the stated integral supported representatives.
\end{proof}

\section{Continuous cyclic characters and explicit families}\label{d3:sec:directcyclic}

This section supplies an independent proof for the explicit completed
three-adic class, and a direct endpoint criterion for all the explicit
prime-power families. The first proof uses an elementary convergent
series. The second uses Coleman functions only as analytic primitives.
The coefficient-cup argument at the end independently supplies actual
local coefficients for existence and rank constructions.

\subsection{A direct detector in continuous cyclic homology}
\label{d3:sec:directtcdetector}

For a spectrum $M$, write $M_{\mathbf Q_p}=M^{\wedge}_p[1/p]$.
Thus $K(R;\mathbf Q_p)$ denotes a rationalized completed spectrum;
it does not denote the tensor product of each actual $K$-group with
$\mathbf Q_p$. Products of units below are products of actual
Quillen $K_1$-classes.

Let $W=\mathcal O_E$, where $E/\mathbf Q_p$ is finite, and let $\pi$
be a uniformizer. If $\mathfrak X=\operatorname{Spf}B$ is smooth and
topologically of finite presentation over $\operatorname{Spf}W$, put
\[
 \Omega_B^j=\widehat\Omega^j_{B/W}[1/p].
\]
The hat denotes continuous differentials, completed before inverting
$p$. Cohomology is ordinary cohomology in $D(E)$, with differential images taken without closure.

\begin{lemma}
\label{d3:lem:relativelogproduct}
Suppose that $\mathfrak X=\operatorname{Spf}B$ is smooth, affine,
topologically of finite presentation, and of relative dimension at
most two over $\operatorname{Spf}W$. There is a natural relative map
\[
 \chi_B:K(B,\pi B)\longrightarrow
                  \Sigma HC^{\mathrm{cts}}(B/W;\mathbf Q_p).
\]
For $v\in1+\pi B$ and $b_1,b_2\in B^\times$, the top HKR component
of its degree-three value is
\begin{equation}\label{d3:eq:tcreviewlogproduct}
 \chi_B([v][b_1][b_2])_{\mathrm{top}}
   =\log(v)\,d\log b_1\wedge d\log b_2
       \quad\text{in }\Omega_B^2/d\Omega_B^1,
\end{equation}
up to one universal sign determined by the norm and suspension
conventions. Here $[v]$ has its relative lift furnished by its reduction
to the identity unit.
\end{lemma}

\begin{proof}
The fiber square of \cite[Proposition~4.10]{AMMN} has rows
\[
\begin{matrix}
 K^{\mathrm{cts}}(\mathfrak X;\mathbf Q_p)
       &\longrightarrow&K(B/\pi;\mathbf Q_p)\\
 \downarrow&&\downarrow\\
 TC^{\mathrm{cts}}(\mathfrak X;\mathbf Q_p)
       &\longrightarrow&TC(B/\pi;\mathbf Q_p)\\
 \downarrow&&\downarrow\\
 HC^{-,\mathrm{cts}}(B/W;\mathbf Q_p)
       &\longrightarrow&HP^{\mathrm{cts}}(B/W;\mathbf Q_p).
\end{matrix}
\]
The last horizontal fiber is $\Sigma HC^{\mathrm{cts}}(B/W;\mathbf Q_p)$,
not negative cyclic homology. Mapping the actual relative spectrum to
these fibers gives $\chi_B$. Consequently relative $K_3$ has target
$HC_2$, and its top component is $\Omega_B^2/d\Omega_B^1$.

We record why the fiber comparison respects multiplication by absolute
$K$-classes. Set $H=THH(B)$, a commutative algebra in cyclotomic spectra.
If $TC_{\mathrm{cyc}}$ denotes the functor on that category, then
$TC_{\mathrm{cyc}}(H)=TC(B)$; no second $THH$ is intended.
All objects $H\otimes D$ and maps $\mathrm{id}_H\otimes f$ in
\cite[Proposition~2.8, Corollary~2.10, Theorem~2.12]{AMMN} are
$H$-modules and module maps. The fixed-point and Tate functors, their
canonical maps, and the norm triangle preserve the induced $TC(B)$-actions
\cite[Corollary~I.4.3]{NikolausScholze}. Thus the inverses of the
equivalences used there are module inverses. The comparisons from
derived to ordinary reduction and then to $B/\pi$ are induced by
$B$-algebra maps. The projection $HH(B)\to HH(B/W)$ in the proof of
Proposition~4.10 also preserves the action. These inverses and fiber
equivalences are used after $p$-completion and rationalization.

The multiplicative cyclotomic trace
\cite[Theorem~7.4]{BGTTrace} now gives $K(B)$-linearity, with target
action through the canonical map $K(B)\to HC^{-,\mathrm{cts}}(B/W;\mathbf Q_p)$.
At the cyclic-complex level, the compatibility of the norm boundary
and its suspension with the $HC^-$-module action is the Connes-sequence
compatibility of \cite[Proposition~2.3]{BokstedtOttosen}.
The identification of the spectral norm with the Connes operator is \cite[Theorem~2.1, Lemma~2.2, Theorem~2.3]{HoyoisCircle};
this includes
the usual graded sign of a connecting map. In particular, no assumption
that an unspecified additive fiber equivalence is multiplicative is
being made.

To normalize a relative unit, first take $B=W\langle T\rangle$ and
$v=1+\pi T$. Its relative character is a function
$f\in HC_0^{\mathrm{cts}}(B/W;\mathbf Q_p)=E\langle T\rangle$.
Choose the norm convention in which its image in $HC^-_1$ is $df$.
The left vertical arrow in \cite[Theorem~2.12]{AMMN} is canonical;
its projection to Hochschild homology takes a unit to the Dennis trace
$v^{-1}\otimes v$. Hence $df=d\log v$. Evaluation at $T=0$ gives
$f(0)=0$. Continuous differentiation on $E\langle T\rangle$ has
only constants in its kernel, and therefore
\[
 f=\log(1+\pi T)=\sum_{n\ge1}(-1)^{n+1}\pi^nT^n/n.
\]
The series converges beyond the closed unit disc. Substitution
$T=(v-1)/\pi$ proves the assertion for every relative unit under
consideration.

Finally use the normalized cyclic complex
$C^+=E[t,t^{-1}]/tE[t]\otimes C_*$, with $|t|=-2$.
The $HC^-\times HC$ action is shuffle plus $t$ times cyclic shuffle
\cite[Definition~2.1 and Remark~2.2]{BokstedtOttosen}.
Starting with the column-zero representative $\log v$, all positive
columns and cyclic-shuffle corrections vanish in $C^+$. After each
of the two unit multiplications only the ordinary shuffle remains.
Under
\[
 b_0\otimes\cdots\otimes b_n\longmapsto
       b_0\,db_1\wedge\cdots\wedge db_n/n!,
\]
the two shuffles in degree two cancel the factor $2!$, giving
\eqref{d3:eq:tcreviewlogproduct}, with the suspension signs just specified.
This calculation involves only degrees through two and commutes with
the continuous, relative-to-$W$ projection. It does not assume a
multiplicative splitting of an infinite HKR filtration.
\end{proof}

\begin{lemma}
\label{d3:lem:tcreviewresidue}
Let $a_0,a_1\in W$ have unit difference, and put
\[
 B=W\langle X,Y,Z\rangle/
           \bigl(XY+(Z-a_0)(Z-a_1)\bigr).
\]
For $F\in E\langle Z\rangle$, the form
$\omega_F=(dX/X)\wedge dF$ is globally regular. If it is exact in
$\Omega_B^2/d\Omega_B^1$, then $F(a_0)=F(a_1)$.
\end{lemma}
\begin{proof}
The two root-deletion affine opens have completed coordinates
\[
 Z=a_i+Xt_i,\qquad
 Y=-t_i(a_i-a_{1-i}+Xt_i)\quad(i=0,1).
\]
They are the completions of $D(X)\cup D(Z-a_{1-i})$, cover the formal
scheme, and have intersection $W\langle X,X^{-1},Z\rangle$.
The forms
\[
 q_i=-(F(Z)-F(a_i))\,dX/X
\]
are regular on the respective bidiscs, satisfy $dq_i=\omega_F$, and
have difference $(F(a_1)-F(a_0))dX/X$.

A closed bidisc form $A(X,t)dX+B(X,t)dt$ has zero Laurent residue
after $t=c/X$, for every $|c|\le1$. Indeed its residue is the
convergent series
\[
 \sum_{m\ge0}(a_{m,m+1}-b_{m+1,m})c^{m+1}=0,
\]
because closedness equates $(m+1)a_{m,m+1}$ and
$(m+1)b_{m+1,m}$. Division by $m+1$ is used only in this individual
coefficient equality, not to construct a bounded primitive.
If $\omega_F=dq$ globally, apply this to $q_i-q$ after restricting
the overlap to $Z=0$, where $c=-a_i$. The difference then has residue
zero, which is the required endpoint equality. This argument includes
$c=0$ and works over the residue field $\mathbf F_2$.
\end{proof}

\begin{proposition}
\label{d3:prop:directtcoriginal}
Let $A=\mathbf Z_3[[x,y,z]]/(3+xy+z^2)$. The actual element
\[
 c=2\left\{\frac{1+z}{2},\,1+x,\,
                        1+\frac{(1+x)y}{4}\right\}\in K_3(A)
\]
has infinite order and has zero image in $K_3(A[1/x])$.
The ring $A$ is regular local of dimension three.
\end{proposition}
\begin{proof}
This is the completed case of Theorem~\ref{thm:explicit-three-dimensional}.
The proof there constructs the relative class on the smooth formal quadric,
uses the actual vanishing $K_4(C)=0$ to pass to absolute $K$-theory,
and detects the resulting class by an endpoint difference in the ordinary
continuous de Rham quotient.
\end{proof}

\paragraph{The finite continuous HKR calculation.}
For the three-adic quadric, put
\[
 B_0=W[X,Y,Z]/(XY+Z^2-1),\qquad B=\widehat{B_0}_3,
 \qquad \Omega_B^i=\widehat\Omega^i_{B/W}[1/3].
\]
The algebra $B_0$ is smooth over $W$ of relative dimension two.
The normalized HKR map on its Hochschild complex is
\[
 b_0\otimes\cdots\otimes b_n\longmapsto
 \frac{1}{n!}b_0\,db_1\wedge\cdots\wedge db_n
 \quad(0\le n\le2),
\]
and is zero in higher degrees.  It is defined over $W$, since $2$ is
a unit, and intertwines the Hochschild differential with zero and the
Connes operator with $d$.  Smoothness makes it a quasi-isomorphism of
mixed complexes: the Hochschild homology is $\Omega^n_{B_0/W}$ for
$0\le n\le2$ and vanishes above degree two.
The comparison with the continuous invariant is supplied by
\cite[Proposition~4.2]{AMMN}.  In the cyclic, negative cyclic, and
periodic total complexes of this finite mixed complex, each total
degree has only finitely many differential-form terms.  Their
reduction towers are surjective, so derived three-adic completion is
computed termwise.  After inverting $3$, this gives in particular
\[
 HC^{\mathrm{cts}}_0(B/W;\mathbf Q_3)=\Omega_B^0,
 \qquad
 HC^{\mathrm{cts}}_2(B/W;\mathbf Q_3)
   =\frac{\Omega_B^2}{d\Omega_B^1}
     \oplus\ker(d:\Omega_B^0\to\Omega_B^1).
\]
The norm map from the degree-zero cyclic term to the degree-one
negative cyclic term is $f\mapsto df$, with the chosen global
suspension sign.  In particular, the displayed quotient is the
ordinary quotient by the image of $d$; its image is not replaced by
its closure.  The same calculation in relative dimension one applies
to the universal unit in $W\langle T\rangle$.


\paragraph{The ordinary quotient is essential.}
The target in \cite[Constructions~4.6 and~4.11]{AMMN} uses the
ordinary derived category. For example, in the preceding quadric put
$\omega_0=(dX/X)\wedge dZ$. The globally regular primitives
$-(Z-Z^{3^j})dX/X$ satisfy
\[
 d\!\left(-(Z-Z^{3^j})\frac{dX}{X}\right)
   =\omega_0-3^jZ^{3^j-1}\omega_0\longrightarrow\omega_0.
\]
The ruling form is therefore in the closure of exact forms although
the residue lemma shows it is not exact. No assertion about
Hausdorff-reduced cohomology or the order at each finite Artin stage
is made here.

\subsubsection{The same computation for the explicit families}

For odd $p$ and $r\ge1$, put $M=p^r$, $n=2M$, and
\[
 A_{p,r}=\mathbf Z_p[[x,y,z]]/(xy+\Phi_{2M}(z-1)),
 \qquad u=z-1,\qquad P(u)=\Phi_{2M}(u).
\]
For $p=2$, take $r\ge2$, $M=2^r$, $n=3M$, and
\[
 A_{2,r}=\mathbf Z_2[[x,y,w]]/(xy+\Phi_M(1+w)),
 \qquad u=1+w,\qquad P(u)=\Phi_M(u).
\]
Let $H(u)=(u^n-1)/P(u)$, $b=yH(u)$,
$\lambda=1-x$, and $\mu=1+\lambda b$. For positive odd $k$ prime
to $p$, put
\[
 g_k=\begin{cases}1-u^k,&p\text{ odd},\\
                  1+u^k+u^{2k},&p=2,
       \end{cases}
 \qquad c_k=k\{g_k,\lambda,\mu\}\in K_3(A_{p,r}).
\]
Every displayed entry is a unit. We use Coleman $\operatorname{Li}_2$
with its series at zero and $\log p=0$, and at roots set
\[
 D_k(a)=\begin{cases}
       \operatorname{Li}_2(a^k),&p\text{ odd},\\
       \tfrac13\operatorname{Li}_2(a^{3k})-
                        \operatorname{Li}_2(a^k),&p=2.
       \end{cases}
\]

\begin{proposition}
\label{d3:prop:directtcendpoints}
The classes $c_k$ vanish integrally after inverting $x$.
If a rational relation $\sum_k q_kc_k=0$ holds in
$K_3(A_{p,r})\otimes\mathbf Q$, then
\[
       \sum_k q_kD_k(a)=0
\]
for every root $a$ of $P$.
\end{proposition}
\begin{proof}
Write $s=1-xb=u^n$. The Dennis--Stein identity gives
$\delta=\langle x,b\rangle=\{\lambda s,\mu\}$.
At odd $p$, Steinberg applied to $u^k$ gives $k\{g_k,s\}=0$.
At two, apply Steinberg to $\zeta_3u^k$ after the finite \'etale
extension adjoining $\zeta_3$ and transfer: since $3\mid n$, the
second entry comes from $u^{kn}$ and the norm of the first is exactly
$g_k$. Thus the same relation holds, without an extra factor two.
It follows that $c_k=k[g_k]\delta$, and the unit expression for
$\delta$ on $D(x)$ makes its generic image zero.

Fix distinct roots $a,a'$ in $E=\mathbf Q_p(\mu_M)$, set
$W=\mathcal O_E$, $d=a'-a$, and
$C(u)=P(u)/((u-a)(u-a'))$. There is an actual convergent map to the
smooth good reduction model
\[
 B=W\langle X,Y,Z\rangle/(XY+Z(Z+1)),\qquad
 u\mapsto a-dZ,\quad x\mapsto dX,\quad
 y\mapsto dC(a-dZ)Y.
\]
All roots reduce to $-1$ at odd $p$ and to $1$ at two, so $d$ and
the images of the original formal parameters are topologically
nilpotent. No inverse of $C$ is used. The special fiber is again an
affine-line torsor over $\mathbf P^1$, so its $K_3$ and $K_4$
vanish after $p$-completion and rationalization.

Set $m=p-1$ at odd $p$ and $m=1$ at two. Then $g_k^m\in1+\pi B$,
and $k[g_k^m][\lambda][\mu]$ is an actual relative lift of $mc_k$.
Let $L_k=\log g_k$ and choose $F_k$ with
$dF_k=L_k\,d\log s=nL_k\,du/u$. These functions converge beyond
the closed $Z$-disc. The differential identity
\[
 d\log\lambda\wedge d\log\mu
 =d\log x\wedge d\log s-d\log s\wedge d\log\mu
 =-dx\wedge db/\mu
\]
gives the following global representative of the relative character:
\[
 \tau_k=mk\frac{dX}{X}\wedge dF_k
                         -d(mkF_k\,d\log\mu).
\]

At odd $p$, $u^k$ lies in the nonsingular residue disc of $-1$;
the analytic differential equation gives
$F_k=-(n/k)\operatorname{Li}_2(u^k)+\mathrm{constant}$.
At two, use the rigid analytic sum
\[
 \operatorname{Li}_2(\zeta_3u^k)+
                   \operatorname{Li}_2(\zeta_3^2u^k).
\]
Its two arguments lie in nonsingular residue discs, its derivative is
$-kL_k\,du/u$, and it descends from the unramified quadratic
extension. Distribution gives its endpoint values $D_k(a)$.
These are only identities of analytic functions; the necessary
analyticity, distribution and inversion are
\cite[Theorem~2.7(1), Proposition~4.4, Proposition~2.10]{BesserDeJeuAlgorithm}.
No assertion of analyticity for the individual $\operatorname{Li}_2$
terms across the residue disc of $1$ is used at two.

Lemma~\ref{d3:lem:tcreviewresidue}, with roots $0,-1$, now computes the
endpoint obstruction of $mc_k$ as
\[
          \pm mn\bigl(D_k(a')-D_k(a)\bigr).
\]
The prefactor $k$ has canceled the $1/k$ in the primitive, leaving
one common scalar. A rational relation therefore makes
$\sum_kq_kD_k(a)$ independent of $a$. Inversion sends roots to
roots and changes this value's sign, since logarithms vanish at
roots of unity. The common value must consequently be zero.
\end{proof}

\begin{corollary}\label{d3:cor:directtcwildorbit}
Using the purely analytic wild-orbit independence statement of
Proposition~\ref{d3:prop:wildorbit}, choose odd representatives
$k_j\equiv(1+p)^j\pmod{p^r}$ for $0\le j<p^{r-1}$ at odd $p$,
and $k_j\equiv5^j\pmod{2^r}$ for $0\le j<2^{r-2}$ at two.
The corresponding actual classes $c_{k_j}$ are rationally independent
in $K_3(A_{p,r})$ and have integral generic image zero.
The same holds over the polynomial quotients with coefficient ring
$\mathbf Z_{(p)}$, localized at the ideal generated by $p$ and the
displayed formal parameters, and over their henselizations.
\end{corollary}
\begin{proof}
Apply Proposition~\ref{d3:prop:directtcendpoints} at one fixed primitive
root and then the stated analytic independence. A rational relation
in a preceding local ring or its henselization would map to one in
the completion. The generic-zero Dennis--Stein relation holds already
in those local rings.
\end{proof}


\section{Finite mixed HKR and classes in every higher degree}
\label{d3:sec:allhigherdegreelaurent}

\begin{lemma}
\label{d3:lem:finitehkrgeneral}
Let $W$ be a complete mixed-characteristic DVR with perfect residue
field, and let $B$ be the $p$-adic completion of a smooth affine
$W$-algebra $B_0$ of relative dimension $D<\infty$.
The continuous cyclic, negative-cyclic, and periodic complexes of
$B/W$, completed before inverting $p$, are computed by the finite
mixed complex of continuous forms $(\Omega_B^{\leq D},0,d)$.
In particular the highest form component of $HC_D^{\mathrm{cts}}$
is $\Omega_B^D/d\Omega_B^{D-1}$.
For $v\in1+\pi B$ and $b_1,\ldots,b_D\in B^\times$, the character
of Lemma~\ref{d3:lem:relativelogproduct} satisfies
\[
 \chi_B\bigl([v]_{\mathrm{rel}}[b_1]\cdots[b_D]\bigr)_{\mathrm{top}}
 =\pm\log(v)\,d\log b_1\wedge\cdots\wedge d\log b_D.
\]
The sign is the same norm and suspension convention as before.
\end{lemma}

\begin{proof}
An integral comparison suffices to control all completion operations.
In Hochschild degree $j\leq D$ use the map
\[
 a_0\otimes\cdots\otimes a_j\longmapsto
       \frac{D!}{j!}a_0\,da_1\wedge\cdots\wedge da_j,
\]
and use zero in higher degrees. It intertwines $b$ with zero and
the Connes operator with $d$. Smooth HKR identifies its map on every
nonzero Hochschild homology group with multiplication by $D!$.
The cone thus has homology in degrees $0,\ldots,D$, all killed by
$p^a$, where $a=v_p(D!)$; flatness of the form modules is used here.

This error can be killed before taking any circle construction.
Its Postnikov tower lies in the category
$\operatorname{Fun}(BS^1,D(W))$ of complexes with circle action.
Each graded piece has trivial circle action: the automorphism
space of a module in the heart is discrete, while $BS^1$ is simply
connected. Hence $p^a$ kills that piece in the equivariant category.
In a triangle $A\to M\to C$, if $p^r$ kills $A$ and $p^s$ kills
$C$, then $p^s\mathrm{id}_M$ factors through $A$, so $p^{r+s}$
kills $M$. Induction on the finite Postnikov tower shows that
$p^{a(D+1)}$ kills the whole equivariant cone.
The exact orbit, fixed-point, Tate, and derived completion functors
preserve this null map. Their comparison errors therefore vanish
after inverting $p$.

In every total degree the cyclic complexes of the finite form
model contain only finitely many form modules. These modules are
$W$-flat and their reduction towers are surjective. Derived
completion is consequently termwise completion, also in the
unbounded total complexes: the cone-of-$1-\mathrm{shift}$ model
for the inverse limit is degreewise exact for these surjections.
Proposition~4.2 of \cite{AMMN}, including its relative version,
identifies this with the asserted continuous invariant. No exchange
of rationalization with an uncontrolled infinite product is used.

The module comparison and the one-variable normalization of a
relative unit in Lemma~\ref{d3:lem:relativelogproduct} do not depend
on the relative dimension. Its ordinary cyclic representative
$\log v$ starts in column zero. At every unit multiplication the
positive columns and the extra $t$ in the cyclic-shuffle correction
vanish in $C^+$. Only the ordinary shuffles of the lowest Dennis
unit traces remain. Their $D!$ permutations cancel the $1/D!$
normalized HKR factor, proving the formula.
\end{proof}

\subsection{The comparison for completed essentially smooth models}

\begin{lemma}
\label{d3:lem:essentiallysmoothAMMN}
Let $W_0=W(\overline{\mathbf F}_p)$ and let $W$
be a finite totally ramified extension, with uniformizer $\pi$
and fraction field $E$. Let $B_0$ be a localization of a smooth
$W$-algebra, with $\Omega^1_{B_0/W}$ finite projective of constant
rank $D$, and put $B=\widehat{B_0}_p$.
There is a natural actual relative character
\[
 \chi_B:K(B,\pi B)\longrightarrow
                 \Sigma HC^{\mathrm{cts}}(B/W;\mathbf Q_p),
\]
which factors through an equivalence from
$K(B,\pi B)^{\wedge}_p[1/p]$. It is $K(B)$-linear through the negative-cyclic
character, and the target has the finite continuous mixed-HKR
model, with ordinary de Rham quotients.
\end{lemma}

\begin{proof}
We justify the extension beyond topological finite presentation.
For each $r$, the ring $W_0/p^r$ is the filtered union of the
finite \emph{\'{e}tale} $\mathbf Z/p^r$-algebras
$W_r(\mathbf F_{p^m})$. Flat base change for cotangent complexes
therefore gives $\widehat L_{W_0/\mathbf Z,p}=0$.
Write $W=W_0[\pi]$ and let $f$ be the Eisenstein polynomial of
$\pi$. Then
\[
 \widehat L_{W/\mathbf Z,p}
   \simeq[W\xrightarrow{f'(\pi)}W]
\]
in cohomological degrees $-1,0$. Choose $a$ with
$p^a\in f'(\pi)W$; multiplication by $p^a$ on this complex is
null-homotopic.

Cotangent transitivity, after derived $p$-completion, gives
\[
 M\longrightarrow\widehat L_{B_0/\mathbf Z,p}
       \longrightarrow P,
 \qquad M=[B\xrightarrow{f'(\pi)}B],\quad
 P=\widehat\Omega^1_{B_0/W}.
\]
The target $P$ is finite projective in degree zero, so this
triangle splits as a triangle of underlying $B$-module
complexes. The error in the $i$th derived exterior power of its
projection is a sum of terms
$(\mathbf L\Lambda^rM)\otimes\Lambda^{i-r}P$, $1\le r\le i$.
The $r$th such term is killed by $p^{ar}$ as an object: derived
exterior power sends the null-homotopy of $p^a\operatorname{id}_M$
to one for $p^{ar}$.

The connective HKR filtration on Hochschild homology has graded
terms $(\mathbf L\Lambda^iL)[i]$. Any fixed homotopy range
involves only finitely many of these terms, even after completion.
Indeed derived $p$-completion preserves connectivity: for a
$c$-connective spectrum $X$, its only possible new bottom group
is $\lim^1_r(\pi_cX/p^r)=0$, because the quotient tower is
surjective. Thus the increasingly connective HKR tails remain
increasingly connective. The preceding bounds
and a finite extension argument therefore make
\[
 HH(B_0;\mathbf Z_p)\longrightarrow HH(B_0/W;\mathbf Z_p)
\]
a quasi-isogeny in each finite range. Its rationalization is an
equivalence. The bounds use the base different and the range;
they are unaffected by the localization set. No interchange of
an infinite localization colimit with completion is used.

To pass to cyclic homology, we control the comparison cone in the
equivariant derived category. Let $C$ be the cone of the map of completed Hochschild
complexes. For every integer $N$, the preceding argument gives an
integer $a_N$ killing each of the finitely many nonzero homotopy
modules of $\tau_{\leq N}C$. The pointwise $t$-structure on
$\operatorname{Fun}(BS^1,D(\mathbf Z_p))$ gives its equivariant
Postnikov tower. Each heart object has trivial circle action:
the automorphism space of a module in the heart is discrete and
$BS^1$ is simply connected. A finite extension argument therefore
gives an integer $b_N$ such that $p^{b_N}$ annihilates
$\tau_{\leq N}C$ as an equivariant object, not merely on its
underlying homotopy groups.

Taking homotopy orbits preserves this scalar null-homotopy.
The omitted Postnikov tail remains $(N+1)$-connective after
homotopy orbits and derived $p$-completion. Indeed, homotopy
orbits of a connective spectrum are connective, and in derived
$p$-completion the potentially new bottom $\lim^1$ term is zero
because the bottom coefficient tower is the surjective tower
$\pi_{N+1}(-)/p^r$. Thus the cone of the corresponding map of
completed cyclic complexes has bounded $p$-torsion in every fixed
degree, and vanishes after inverting $p$.

Equivalently, one can take a sufficiently high finite bar skeleton
for the circle-orbit construction in each degree. Its building blocks
are finite colimits of $(S^1)^j_+\wedge X$, so the comparison commutes
with derived completion. This argument requires no uniform bound
$b_N$ as $N$ varies and no interchange of rationalization with an
infinite product. The norm fiber sequence then makes the square
between absolute and relative $HC^-\to HP$ cartesian. It does not
require that the two vertical maps of that square be equivalences
individually.

By \cite[Proposition~4.2]{AMMN}, including the relative variant
stated just before Proposition~4.10, both $B_0$ and $B$ compute
the same continuous $p$-completed HH and cyclic theories.
The ordinary ring $B$ is henselian along $(p)$, so
\cite[Theorem~A]{AMMN} applies. The map $B/p\to B/\pi$ has
nilpotent kernel, and \cite[Corollary~3.8]{AMMN} identifies its
K-theory and TC after rationalized $p$-completion. Composing the
actual-to-completed map on relative K-theory with the resulting
fiber equivalence gives $\chi_B$.

The module argument is unchanged: for $H=THH(B)$, the objects
$H\otimes D$ and maps $\operatorname{id}_H\otimes f$ in
\cite[Proposition~2.8, Corollary~2.10, and Theorem~2.12]{AMMN}
form a diagram of $H$-modules in cyclotomic spectra.
The lax monoidal TC and circle functors give module diagrams,
their equivalences can be inverted there, and the relative HH
projection preserves the action. Together with the
multiplicative cyclotomic trace this gives the asserted actual
$K(B)$-linearity. This construction has no finite presentation
requirement.

Finally, the integral mixed HKR map in degree $j\le D$ has
coefficient $D!/j!$ and induces multiplication by $D!$ on
$HH_j(B_0/W)=\Omega^j_{B_0/W}$. Its equivariant cone has a
Postnikov tower of length $D+1$ whose graded objects are killed
by $p^{v_p(D!)}$. These heart objects have trivial circle action;
the extension argument kills the entire equivariant cone by
$p^{(D+1)v_p(D!)}$. Thus circle orbits, fixed points, Tate
constructions, and derived completion retain only bounded
torsion error. The rational mixed-HKR comparison follows.
Each total degree has finitely many flat form modules, with
surjective reduction towers, so completion is termwise and
cohomology uses the ordinary image of $d$.
\end{proof}

\subsection{Continuous residues and higher unit products}

\begin{lemma}
\label{d3:lem:continuouslaurentresidue}
Let $W$ be a complete DVR of mixed characteristic with uniformizer
$\pi$ and fraction field $E$. Let $B$ be a smooth formal $W$-algebra
of topologically finite type. Take a nonzero
$g\in W[T_1,\ldots,T_s]$ whose nonzero coefficients are units, and set
\[
 B_g=\bigl(B[T_1^{\pm1},\ldots,T_s^{\pm1},1/g]\bigr)^{\wedge}_{\pi}.
\]
There are $E$-linear maps on continuous forms, with base forms
placed before the logarithmic differentials,
\[
 \operatorname{Res}_s:
 \Omega^{q+s}_{B_g[1/p]/E,\mathrm{cts}}
       \longrightarrow\Omega^q_{B[1/p]/E,\mathrm{cts}},
\]
such that
\[
 \operatorname{Res}_s(d\eta)=d\operatorname{Res}_s(\eta),\qquad
 \operatorname{Res}_s\left(
 \omega\wedge d\log T_1\wedge\cdots\wedge d\log T_s\right)=\omega.
\]
Consequently wedging with these logarithmic differentials gives
a split injection on ordinary continuous de Rham cohomology.
There is also a degree-zero projection $P_0$ which takes the
constant coefficient of the summand containing no $d\log T_i$.
It commutes with $d$ and is the identity on continuous base forms.
\end{lemma}

\begin{proof}
Put $R_0=W[T_1^{\pm1},\ldots,T_s^{\pm1},1/g]$.
There is an integral formal expansion
\[
 R_0\longrightarrow W((T_1))\cdots((T_s)).
\]
To construct it, write $g=T_s^a(g_a+T_sg_{a+1}+\cdots)$ with
$a$ minimal. By induction $g_a$ is a unit in the preceding
iterated Laurent ring: the recursion ends at one of the unit
coefficients of $g$. Expand the remaining inverse as a formal
geometric series in $T_s$. This construction divides by no power
of $\pi$.
Taking the coefficient of $T_s^0$, then $T_{s-1}^0$, and so on,
defines a $W$-linear map
\[
 \operatorname{CT}:R_0\longrightarrow W,
 \qquad \operatorname{CT}(1)=1,\quad
 \operatorname{CT}(T_i\partial_{T_i}h)=0.
\]
It is uniformly $\pi$-adically continuous, since
$\operatorname{CT}(\pi^aR_0)\subseteq\pi^aW$ for every $a$.
The $B$-linear extension therefore extends uniquely to $B_g\to B$.
The same construction works for every finite projective module
of continuous base forms, by taking a direct summand of a finite
free module.

Decompose forms into base forms and the ordered $d\log T_i$.
Retain only the summand containing all $s$ logarithmic factors,
and apply this extended constant-term map to its coefficient.
This defines $\operatorname{Res}_s$ integrally before inverting $p$.
The base differential commutes with constant-term extraction;
every term from a Laurent-variable derivative has zero constant
term. The two displayed identities follow, first on the dense
uncompleted modules and then by continuity. They show directly
that boundaries map to boundaries and that the stated composite
is the identity on forms and cohomology.
For $P_0$, a Laurent-variable derivative always introduces a
logarithmic factor and is discarded, proving its chain identity
by the same calculation.

No embedding of the completed ring into the ordinary Laurent ring
is asserted. For example $\sum_{a\geq0}\pi^aT_s^{-a}$ has no lower
bound on its Laurent exponents. If a completed series target is
desired, it is
$\varprojlim_a(B/\pi^a)((T_1))\cdots((T_s))$.
Nor need the formal expansion of $1/g$ converge analytically
when every $T_i$ has norm one.
\end{proof}


\begin{lemma}
\label{d3:lem:laurentspecialvanishing}
Let $X$ be a smooth variety of dimension $d$ over a perfect field
of characteristic $p$. For every $j>d$,
\[
 \pi_jK(X)^{\wedge}_p=\pi_jK(X;\Q_p)=0.
\]
In particular, if $s\geq1$, $0\ne\overline g\in
\F_3[T_1^{\pm1},\ldots,T_s^{\pm1}]$, and
\[
 C_s=\Spec\!\left(
 \F_3[X,Y,Z,T_1^{\pm1},\ldots,T_s^{\pm1},1/\overline g]
                         /(XY+Z^2-1)\right),
\]
then
\[
 \pi_{s+3}K(C_s;\Q_3)=\pi_{s+4}K(C_s;\Q_3)=0.
\]
\end{lemma}
\begin{proof}
Geisser--Levine prove
\[
 K_j(X;\Z/p^a)=0\qquad(a\geq1,\ j>d)
\]
in \cite[Theorem~8.4, p.~491]{GeisserLevineCharP}.
These are the finite-coefficient $K$-groups, with the usual
coefficient long exact sequence; see their Section~2, p.~464.
In particular the assertion includes the torsion contribution
from the preceding actual $K$-group.
The homotopy-limit exact sequence is
\[
 0\longrightarrow\varprojlim_a^{1}K_{j+1}(X;\Z/p^a)
 \longrightarrow\pi_jK(X)^{\wedge}_p
 \longrightarrow\varprojlim_a K_j(X;\Z/p^a)
 \longrightarrow0.
\]
Both outside terms vanish when $j>d$. Inverting $p$ proves the
first assertion. The displayed $C_s$ is smooth of dimension
$s+2$, so its two asserted groups lie in this range.
\end{proof}


\begin{theorem}
\label{d3:thm:everydegreelaurent}
Put
\[
 A_0=\left(\Z_{(3)}[x,y,z]/(3+xy+z^2)\right)_{(x,y,z)},
 \qquad \mathfrak m_0=(x,y,z)A_0.
\]
For $s\geq0$, let
\[
 R_s=A_0[T_1,\ldots,T_s]_{\mathfrak m_0A_0[T_1,\ldots,T_s]}.
\]
This is a regular local ring of dimension three and residue field
$\F_3(T_1,\ldots,T_s)$. The actual class
\[
 c_s=2\left\{\frac{1+z}{2},\,1+x,\,
       1+\frac{(1+x)y}{4},\,T_1,\ldots,T_s\right\}
                    \in K_{s+3}(R_s)
\]
has infinite order and restricts to zero in $K_{s+3}(R_s[1/x])$
integrally. Consequently rational Gersten injectivity fails in
each degree $j\geq3$ on an explicitly specified ring of dimension
three.
\end{theorem}

\begin{proof}
Polynomial extension and localization preserve regularity; the
chosen prime has height three and the displayed residue field.
Every $T_i$ is a unit. The Dennis--Stein calculation of
Proposition~\ref{d3:prop:directtcoriginal}, followed by multiplication
by these units, proves integral vanishing after inverting $x$.
The case $s=0$ has already been proved.

Suppose that $N c_s=0$ for some positive integer $N$ and $s\geq1$.
Continuity of actual $K$-theory makes this relation hold over
$A_0[T_i^{\pm1},1/f]$ for one finite denominator with
$\bar f\ne0$ in $\F_3[T_1,\ldots,T_s]$.
Use $W=\Z_3[u]/(u^2+3)$ and the original good quadric
$B=W\langle X,Y,Z\rangle/(XY+Z^2-1)$.
Choose $g\in W[T_1,\ldots,T_s]$ by lifting exactly the nonzero
coefficients of $\bar f$ to units, and form $B_g$ as in
Lemma~\ref{d3:lem:continuouslaurentresidue}.
Under $(x,y,z)=(uX,uY,uZ)$, the image of $f$ is congruent to $g$
modulo $u$. Thus $f/g\in1+uB_g$ is a unit, and the finite
relation has an actual image in $K_{s+3}(B_g)$.

The specified image has the actual relative lift
\[
 \xi_s=[a^2]_{\mathrm{rel}}[\lambda][\mu]
                                  [T_1]\cdots[T_s].
\]
The formal algebra $B_g$ has relative dimension $D=s+2$ over $W$.
Lemma~\ref{d3:lem:finitehkrgeneral} gives its top character as
\[
 \chi_{B_g}(\xi_s)_{\mathrm{top}}
       =\pm\tau_c\wedge d\log T_1\wedge\cdots\wedge d\log T_s,
\]
where $[\tau_c]\ne0$ was calculated in
Proposition~\ref{d3:prop:directtcoriginal}.
The residue of Lemma~\ref{d3:lem:continuouslaurentresidue} sends this
class back to $\pm[\tau_c]$, so it is nonzero in the ordinary
quotient by exact forms.

By Lemma~\ref{d3:lem:laurentspecialvanishing},
$\pi_{s+4}K(B_g/u;\Q_3)=0$. The long exact sequence of the
completed relative spectrum therefore makes
\[
 \pi_{s+3}K(B_g,uB_g;\Q_3)
       \longrightarrow\pi_{s+3}K(B_g;\Q_3)
\]
injective. The nonzero character of $\xi_s$ shows that its absolute
image has infinite order, contrary to the image of $Nc_s=0$.
This proves the assertion for the actual class in $R_s$.
Only the auxiliary special fiber is subject to the finite-coefficient
vanishing theorem; its perfect ground field is $\F_3$ even though
the source ring has the stated imperfect residue field.
\end{proof}


\section{Coefficient products and infinite rank}
The continuous cyclic character also detects products of actual point
classes with a ruling class. This gives an independent construction
of the coefficient classes used in the rank arguments.

\subsection{Coefficient products in cyclic homology}
\label{d3:sec:directtccoefficients}

Write $K(R;\mathbf Q_p)=K(R)^{\wedge}_p[1/p]$; this notation denotes
a rationalized completed spectrum. Let $E/\mathbf Q_p$ be finite, let
$W=\mathcal O_E$, and let $k$ be its residue field. By
\cite[Proposition~4.10 and Example~4.3]{AMMN}, the point has a natural
relative character
\[
 r_W:K_3(W)\longrightarrow\pi_3K(W;\mathbf Q_p)
 \xrightarrow{\sim}HC_2^{\mathrm{cts}}(W/W;\mathbf Q_p)=E.
\]
The isomorphism uses
$\pi_3K(k;\mathbf Q_p)=\pi_4K(k;\mathbf Q_p)=0$.

\begin{lemma}
\label{d3:lem:cherncompletionfactor}
Let $E/\Q_p$ be finite, let $W=\mathcal O_E$, and let
$n\geq2$ and $s=2n-1$. Write
\[
 K(R)^{\wedge}_p=\operatorname*{holim}_{\nu}K(R)/p^\nu,
 \qquad K(R;\Q_p)=K(R)^{\wedge}_p[1/p].
\]
There is a natural $\Q_p$-linear map
\[
 \widetilde c_{n,E}:\pi_sK(W;\Q_p)
             \longrightarrow H^1(E,\Q_p(n))
\]
such that the usual \'etale Chern class on actual $K$-theory is
the composite
\[
 K_s(W)_{\Q}\longrightarrow\pi_sK(W;\Q_p)
       \xrightarrow{\widetilde c_{n,E}}H^1(E,\Q_p(n)).
\]
These maps are natural under finite extensions of $p$-adic fields.
Consequently, actual classes with $\Q_p$-linearly independent
\'etale Chern values have $\Q_p$-linearly independent images in
$\pi_sK(W;\Q_p)$.
\end{lemma}

\begin{proof}
First construct the factor over the field $E$, where $p$ is invertible.
For every $\nu\geq1$ the finite-coefficient Chern class is an
additive map
\[
 c^{(\nu)}_{n,1}:K_s(E;\Z/p^\nu)
       \longrightarrow H^1(E,\mu_{p^\nu}^{\otimes n}).
\]
This includes $p=2$: the exceptional $K$-degree in
\cite[Proposition~2.4, p.~10]{WeibelChern} is two, whereas $s\geq3$.
The maps are compatible with $p^{\nu+1}\to p^\nu$ reduction.
Indeed, take $q=p$ and $r=p^\nu$ in the middle square of
\cite[Proposition~2.8, p.~12]{WeibelChern}:
\[
\begin{array}{ccc}
 K_s(E;\Z/p^{\nu+1})&\longrightarrow&K_s(E;\Z/p^\nu)\\
 {\scriptstyle c^{(\nu+1)}_{n,1}}\downarrow&&
             \downarrow{\scriptstyle c^{(\nu)}_{n,1}}\\
 H^1(E,\mu_{p^{\nu+1}}^{\otimes n})&\longrightarrow&
                         H^1(E,\mu_{p^\nu}^{\otimes n}).
\end{array}
\]
The homotopy-limit projections therefore define a homomorphism
\begin{equation}\label{d3:eq:cherncompletionintegral}
 \pi_sK(E)^{\wedge}_p
 \longrightarrow\varprojlim_{\nu} K_s(E;\Z/p^\nu)
 \xrightarrow{\varprojlim c^{(\nu)}_{n,1}}
 \varprojlim_{\nu}H^1(E,\mu_{p^\nu}^{\otimes n}).
\end{equation}
No assertion about a source $\varprojlim^{1}$ term is required:
\eqref{d3:eq:cherncompletionintegral} uses only the canonical projections.
On the target, the groups $H^0(E,\mu_{p^\nu}^{\otimes n})$ are
finite. Their inverse system satisfies the Mittag--Leffler condition,
so \cite[equation~(0.2), p.~208]{Jannsen} gives
\[
 \varprojlim_{\nu}H^1(E,\mu_{p^\nu}^{\otimes n})
       \cong H^1_{\mathrm{cont}}(E,\Z_p(n)).
\]

The resulting map is $\Z_p$-linear. For clarity, the source homotopy
group has its natural $\Z_p=\pi_0\mathbb S^{\wedge}_p$ action.
Each composite to the $\nu$th cohomology group is additive and has
target annihilated by $p^\nu$. If $\lambda\in\Z_p$ and
$\lambda=a+p^\nu b$ with $a\in\Z$ and $b\in\Z_p$, that composite
sends $\lambda x$ to $a$ times the image of $x$, which is also
$\lambda$ times that image. Passing to the inverse limit proves
linearity. Inverting $p$ and using the definition of rational
continuous cohomology \cite[equation~(0.4)]{Jannsen} gives
\[
 \overline c_{n,E}:\pi_sK(E;\Q_p)
                 \longrightarrow H^1(E,\Q_p(n)).
\]
Define $\widetilde c_{n,E}$ by preceding this map with the natural
map $\pi_sK(W;\Q_p)\to\pi_sK(E;\Q_p)$.

By \cite[Lemma~2.3, p.~10]{WeibelChern}, at each finite level the
usual Chern class on $K_s(E)$ is the finite-coefficient class after
reduction. The target inverse-limit identification then gives the
asserted equality on actual classes. Naturality follows from
\cite[Proposition~2.1.1]{WeibelChern} and the naturality of every
map above. Finally, applying the $\Q_p$-linear map
$\widetilde c_{n,E}$ to a relation proves the independence assertion.
\end{proof}

\begin{remark}
This lemma does not assert that actual $K_s(W)$ surjects onto the
homotopy group of its completed spectrum. It also does not identify
any cyclic or syntomic coordinate with a polylogarithm. If an
independently constructed $\Q_p$-linear isomorphism
$\pi_sK(W;\Q_p)\cong E$ is available, the lemma only transports
independence to that coordinate.
\end{remark}

\begin{lemma}\label{d3:lem:directtcactualcoefficients}
The values of $r_W$ on actual classes span $E$ over $\mathbf Q_p$.
For a quadratic extension $E/F$, actual integral anti-invariant
classes span $E^{-}$ under this map.
\end{lemma}
\begin{proof}
Let $d=[E:\mathbf Q_p]$ and put
$V=\pi_3K(\mathcal O_E;\mathbf Q_p)$.
The point comparison of \cite[Proposition~4.10]{AMMN} identifies $V$ with $E$, so
$\dim_{\mathbf Q_p}V=d$. Independently, the Bloch--Kato exponential
identifies $H^1(E,\mathbf Q_p(2))$ with $E$, also of dimension $d$.
The finite-coefficient Chern construction gives a natural linear map
\[
 \widetilde c_{2,E}:V\longrightarrow H^1(E,\mathbf Q_p(2))
\]
through which the Chern map on actual $K_3(\mathcal O_E)$ factors.
This is the case $n=2$ of Lemma~\ref{d3:lem:cherncompletionfactor};
its proof uses only finite-coefficient Chern additivity,
coefficient compatibility, and the canonical map from the homotopy
limit to the limit of homotopy groups.

By Lemma~\ref{d3:lem:actuallocal}, whose proof uses Levine's
finite-coefficient theorem, Milnor $K_3$-divisibility, the
Moore--Merkurjev description of $K_2(E)$, and localization,
the actual Chern values span $H^1(E,\mathbf Q_p(2))$.
Let $V_{\rm act}$ be the span in $V$ of the images of actual classes.
Then $\widetilde c_{2,E}(V_{\rm act})=H^1(E,\mathbf Q_p(2))$.
Both ambient vector spaces have dimension $d$, whence
$V_{\rm act}=V$ and $\widetilde c_{2,E}$ is an isomorphism.
Thus the actual AMMN point values span $E$.

All maps are natural for finite extensions of local fields.
For a quadratic $E/F$, apply $1-\sigma$ to actual integral classes.
Their images span $E^-$ since $(1-\sigma)E=E^-$ as
$\mathbf Q_p$-vector spaces. This argument divides by no integer
inside actual integral $K$-theory, and applies equally at $p=2$.
\end{proof}

\begin{proposition}\label{d3:prop:directtccoefficientcup}
Let $0\ne d\in\mathfrak m_W$, and set
\[
 T=\bigl(W[[x,y,w]]/(xy+w(w+d))\bigr)[1/p],\qquad J=(x,w)\mathcal O_T.
\]
If $\beta\in K_3(W)$ satisfies $r_W(\beta)\ne0$, then
$\beta(1-[J])$ has infinite order in $K_3(T)$. Actual coefficients
with $\mathbf Q_p$-independent $r_W$-values give rationally independent
products.
\end{proposition}
\begin{proof}
Put
\[
 B=W\langle X,Y,Z\rangle/(XY+Z(Z+1)),\qquad I=(X,Z),\qquad C=B/\pi.
\]
The substitution $(x,y,w)=(dX,dY,dZ)$ converges for every formal
series and defines a ring map to $B$. On the generic fiber it takes
$J$ to a line bundle isomorphic to $I$. The special fiber $C$ is an
affine-line torsor over $\mathbf P^1_k$, so
$K(C)\simeq K(k)\oplus K(k)$ and
$\pi_3K(C;\mathbf Q_p)=\pi_4K(C;\mathbf Q_p)=0$.

Let $b=r_W(\beta)$ and $\eta=1-[I]$. The module compatibility of the
AMMN fiber character, verified in Lemma~\ref{d3:lem:relativelogproduct},
identifies the relative image of $\beta\eta$ with the action of
$\operatorname{ch}^{-}(\eta)$ on $t^{-1}b\in HC_2$, where $|t|=-2$.
Write a normalized representative of the former as
$c_0+tc_2+t^2c_4+\cdots$. Its rank term $c_0$ is zero. In the
shuffle-plus-cyclic-shuffle product
\cite[Section~2, Remark~2.2]{BokstedtOttosen}, every cyclic-shuffle
correction puts $b\in E$ in a normalized bar slot and hence is zero.
The remaining product in ordinary cyclic homology is exactly $bc_2$:
terms of positive $t$-degree vanish. Its two-form component is
\[
 -\lambda_E b\,c_1^{\mathrm{dR}}(I),\qquad\lambda_E\in E^\times,
\]
where the nonzero normalization follows from
\cite[Proposition~4.12]{AMMN}. In weight one, the map
$HC^-_0\to HP_0$ induces
$H^2(\Omega_B^{\geq1})\to H^2(\Omega_B^\bullet)$; both groups
are $\Omega_B^2/d\Omega_B^1$, so this is an isomorphism.
Consequently the periodic Chern component determines precisely the
two-form component used above. Thus no comparison with an \'etale
or syntomic regulator, and no additional exact normalization of
$K_0\to HC^{-}$, is needed.

The class $c_1^{\mathrm{dR}}(I)$ is represented, up to sign, by the
globally regular form
\[
 \omega_0=\frac{dX}{X}\wedge dZ
 =\operatorname{Tr}(P\,dP\wedge dP),\qquad
 P=\begin{pmatrix}-Z&X\\Y&Z+1\end{pmatrix}.
\]
To verify its algebraic nonvanishing, use the two root-deletion
bidiscs with $Z=Xt_0$ and $Z=-1+Xt_{-1}$. Their intersection is
$E\langle X,X^{-1},Z\rangle$. Local primitives are
\[
 q_0=-Z\,dX/X,\qquad q_{-1}=-(Z+1)dX/X,
 \qquad q_{-1}-q_0=-dX/X.
\]
If a global primitive existed, subtract it from these primitives,
restrict to $Z=0$, and take the Laurent residue in $X$. A closed
bidisc form $A(X,t)dX+B(X,t)dt$ has zero residue on $t=c/X$ for
every $|c|\le1$: its residue is
\[
 \sum_{m\ge0}(a_{m,m+1}-b_{m+1,m})c^{m+1}=0,
\]
by closedness and convergence of the Tate coefficients. Here the
two values of $c$ are $0$ and $1$, so the argument also applies at
$p=2$. The displayed difference has residue $-1$, a contradiction.
This also proves that $b\mapsto b[\omega_0]$ is injective.

The fiber square therefore detects $\beta\eta$ nontrivially in
$\pi_3K(B;\mathbf Q_p)$. Finally, localization preserves its rational
nonvanishing in $K_3(B[1/p])$: its possible kernel comes from
$K_3(C)$, a finite group. This is the image of $\beta(1-[J])$.
The argument is linear in the coefficient and proves the independence
assertion as well.
\end{proof}

\begin{remark}
The target in this proof is the ordinary de Rham quotient
$\Omega^2/d\Omega^1$ of the termwise completed complex, as in
\cite[Constructions~4.6 and~4.11]{AMMN}. It is not the quotient by
the closure of $d\Omega^1$. Indeed, with
$h_j(Z)=-(-Z)^{p^j}$, the regular primitives
$-(Z-h_j(Z))dX/X$ have derivatives converging to $\omega_0$, although
$\omega_0$ is not exact. This proof gives nonzero images of actual
$K$-classes in a rationalized completed spectrum; it does not assert
a prescribed order in every finite Artin approximation.
\end{remark}


\subsection{Infinite rank over a strict henselian base at every prime}
\label{d3:sec:directtcstrictallprimes}

We first extend the coefficient-cup calculation to the infinite
unramified base. Throughout this subsection, $K(-;\mathbf Q_p)$ means
the rationalization of the $p$-completed spectrum.

\begin{lemma}\label{d3:lem:directtcperfectcoeffcup}
Let $E$ be a complete discretely valued field of mixed characteristic
$(0,p)$ whose residue field $k$ is algebraic over $\mathbf F_p$, and
let $W=\mathcal O_E$. The AMMN point character
\[
 r_W:K_3(W)\longrightarrow\pi_3K(W;\mathbf Q_p)
       \xrightarrow{\sim}HC_2^{\mathrm{cts}}(W/W;\mathbf Q_p)=E
\]
is defined as before and is natural for continuous local maps of
these coefficient DVRs. Proposition~\ref{d3:prop:directtccoefficientcup}
holds for this $W$: any finite set of actual coefficients with
$\mathbf Q_p$-independent $r_W$-values gives rationally independent
products with the ruling class on the completed split node.
\end{lemma}
\begin{proof}
The standing hypotheses of \cite[Section~4.1]{AMMN} require a perfect
residue field, and Proposition~4.10 requires a smooth formal scheme;
neither requires a finite residue field. Example~4.3 applies to the
affine $p$-complete rings in question. Naturality when the coefficient
ring changes follows by keeping the canonical maps used in the proof
of Proposition~4.10: a map of pairs $(W,R)\to(W',R')$ gives a
commutative square
\[
\begin{matrix}
 HH(R;\mathbf Q_p)&\longrightarrow&HH(R/W;\mathbf Q_p)\\
 \downarrow&&\downarrow\\
 HH(R';\mathbf Q_p)&\longrightarrow&HH(R'/W';\mathbf Q_p).
\end{matrix}
\]
The absolute fiber square, cyclic constructions, and horizontal
fibers are natural as well. At the point, the cyclic scalar is sent
by $E\to E'$. Thus $r_{W'}(\beta|_{W'})=r_W(\beta)|_{E'}$.

Quillen's finite-field calculation and filtered-colimit continuity
show that $K_i(k)$ is uniquely $p$-divisible for $i>0$, and that
$K_3(k)$ is torsion. The coefficient exact sequence gives
$K(k)/p^a\simeq H\mathbf Z/p^a$ directly, hence
$K(k;\mathbf Z_p)\simeq H\mathbf Z_p$. In particular the two adjacent
positive homotopy groups needed to define the point fiber lift vanish.

For $B=W\langle X,Y,Z\rangle/(XY+Z(Z+1))$, its reduced special fiber
$C$ is an affine-line torsor over $\mathbf P^1_k$. Therefore
$K(C)\simeq K(k)\oplus K(k)$. The relative-to-absolute step in
Proposition~\ref{d3:prop:directtccoefficientcup} still has vanishing
$\pi_3$ and $\pi_4$ on this special fiber. Its ordinary de Rham
calculation and two-chart residue proof work over the complete field
$E$ verbatim. In particular $b\mapsto b[\omega_0]$ is injective,
where $\omega_0=(dX/X)\wedge dZ$ is globally regular. Finally
$K_3(C)$ is torsion, which is enough for localization to give
$K_3(B)_{\mathbf Q}\hookrightarrow K_3(B[1/p])_{\mathbf Q}$.
The module calculation and the nonzero scalar of
\cite[Proposition~4.12]{AMMN} therefore detect exactly the same
coefficient products. No Hausdorff quotient of de Rham cohomology
is taken.
\end{proof}

\begin{theorem}\label{d3:thm:directtcstrictinfiniteallp}
For every prime $p$, the complete strictly henselian regular local ring
\[
 A_{\infty,p}=W(\overline{\mathbf F}_p)[[x,y,z]]/(p+xy+z^2)
\]
has dimension three and infinite-dimensional rational generic
$K_3$-kernel. More precisely, for every $m\ge1$ there are $m$ actual
integral classes with zero generic image whose rational classes are
independent.
\end{theorem}
\begin{proof}
Put $V_m=W(\mathbf F_{p^m})$, $V_\infty=W(\overline{\mathbf F}_p)$,
$F_m=\operatorname{Frac}V_m$, and choose $u^2=-p$. Write
$W_m=V_m[u]$, $E_m=F_m(u)$, $W_\infty=V_\infty[u]$, and
$E_\infty=\operatorname{Frac}W_\infty$. Eisenstein's criterion makes
each $W_m$ and $W_\infty$ a complete DVR with uniformizer $u$.
The linear term $p$ of the equation in the regular ambient
power-series ring proves regularity and dimension three for
$A_{\infty,p}$. Its residue field is algebraically closed.

By Lemma~\ref{d3:lem:directtcactualcoefficients}, actual $r_{W_m}$-values
span $E_m$. Apply $1-\sigma$, where $\sigma(u)=-u$, and choose
$\beta_{m,1},\ldots,\beta_{m,m}\in K_3(W_m)$ with
$\sigma\beta_{m,j}=-\beta_{m,j}$ whose values are independent in
$E_m^-=uF_m$. This construction involves no division by two in
actual $K$-theory. Lemma~\ref{d3:lem:directtcperfectcoeffcup} carries
these values by the field inclusion into $E_\infty$, preserving
their $\mathbf Q_p$-independence.

Set $A_m=V_m[[x,y,z]]/(p+xy+z^2)$. Pull each coefficient to
$A_m/(x)=W_m[[y]]$ and push forward along $x=0$, obtaining
$c_{m,j}\in K_3(A_m)$. These vanish on $A_m[1/x]$ integrally.
The local map $A_m\to A_{\infty,p}$ is flat by miracle flatness:
both rings are regular of dimension three, and its closed fiber
is a field \cite[Tag~00R4]{StacksFlatness}. Thus their images
are the corresponding supported pushforwards on $A_{\infty,p}$.

After the finite free quadratic coefficient extension and inversion
of $p$, put
\[
 T=\bigl(W_\infty[[x,y,z]]/(xy+(z-u)(z+u))\bigr)[1/p],
 \qquad J_+=(x,z-u),\quad\eta_+=1-[J_+].
\]
The two components of $x=0$ are disjoint Cartier divisors. Their
structure-sheaf classes sum to zero because their union is the
principal divisor $x$. Flat base change and the projection formula
therefore give exactly
\[
 c_{m,j}|_T=\beta_{m,j}\eta_+
          +\sigma\beta_{m,j}\eta_-
          =2\beta_{m,j}\eta_+.
\]
Use $w=z-u$ and $d=2u$ in
Lemma~\ref{d3:lem:directtcperfectcoeffcup}. The substitution into the
good quadric is
$x=dX$, $y=dY$, $z=u+dZ$; all original formal variables have norm
at most $|u|<1$, including at $p=2$. The coefficient detector gives
$m$ independent images, so the original rational classes are
independent. Their generic images are zero by localization at $x$.
Since $m$ is arbitrary, the asserted kernel is infinite dimensional.
\end{proof}

\begin{corollary}\label{d3:cor:directtcstrictarithmeticallp}
For every prime $p$, the arithmetic strict henselization
\[
 S_p=\left(\left(\mathbf Z_{(p)}[x,y,z]/(p+xy+z^2)\right)
                         _{(x,y,z)}\right)^{sh}
\]
also has infinite-dimensional rational generic $K_3$-kernel,
represented by actual integral classes with integral generic vanishing.
\end{corollary}
\begin{proof}
Choose a finite \'etale local DVR $V_{m,0}/\mathbf Z_{(p)}$ with
number-field fraction field, residue $\mathbf F_{p^m}$, completion
$V_m$, and a compatible residue embedding. Set $W_{m,0}=V_{m,0}[u]$,
where $u^2=-p$. The fraction field of $W_{m,0}^h$ is the field of
constants in $E_m$, by Krasner's lemma.

The second assertion of \cite[Theorem~4.13]{LevineIndK3} and
$K_3^M(E_m)/p^a=0$ imply that
$K_3(\operatorname{Frac}W_{m,0}^h)$ surjects onto $K_3(E_m)/p^a$.
The Milnor vanishing follows from norm-residue
\cite[VI, Theorem~4.1]{Weibel} and local
$p$-cohomological dimension two, also at $p=2$
\cite[Chapter~I, discussion preceding Theorem~2.8]{MilneDuality}. Localization, using
$K_2(\mathbf F_{p^m})=0$, then gives
\[
 K_3(W_{m,0}^h)\twoheadrightarrow K_3(W_m)/p^a
 \qquad(a\ge1).
\]
Here the passage from $E_m$ back to $W_m$ does not change the finite
quotient because its kernel comes from prime-to-$p$ torsion
$K_3(\mathbf F_{p^m})$. The actual-completion argument of
Lemma~\ref{d3:lem:directtcactualcoefficients} now shows directly that
the henselian values of the AMMN point map span $E_m$. Applying
$1-\sigma$ supplies $m$ actual anti-invariant coefficients with
independent AMMN values. No comparison with an \'etale-selected
coefficient class is used.

Let $R_m=(A_{0,p}\otimes_{\mathbf Z_{(p)}}V_{m,0})^h$.
Its divisor quotient by $x$ receives $W_{m,0}^h$, its completion is
$A_m$, and it maps compatibly to $S_p$. Push the chosen coefficients
forward along $x=0$ and then map to $S_p$. Their images under
$S_p\to A_{\infty,p}$ are precisely the completed classes detected
in Theorem~\ref{d3:thm:directtcstrictinfiniteallp}; the map and completion
identification are supplied uniformly in $p$ by
Lemma~\ref{d3:lem:stricthenselcompletion}. Hence these $m$ rational
classes on $S_p$ are independent. They vanish after inverting $x$
integrally. Letting $m$ grow proves the corollary.
\end{proof}

\section{Higher point coordinates and unramified coefficients}
\label{d3:sec:higherpointcoordinates}


\begin{lemma}
\label{d3:lem:formalhigherTCcup}
Let $W$ be a complete mixed-characteristic discrete valuation ring
with uniformizer $\pi$, fraction field $E$, and residue field $k$
algebraic over $\mathbf F_p$. Let $B$ be a smooth affine formal
$W$-algebra, topologically of finite presentation and of relative
dimension two, and put $C=B/\pi B$.
Fix $n\geq2$ and suppose that
\[
 K_{2n}(C)=0.
\]
Let $\eta\in K_0(B)$ have rank zero. Write
\[
 \operatorname{ch}^-(\eta)=t\gamma,
 \qquad
 \gamma\in\widehat\Omega^2_{B/W}[1/p]
               /d\widehat\Omega^1_{B/W}[1/p],
 \qquad |t|=-2,
\]
for its negative-cyclic character in the finite mixed-HKR model,
and assume $\gamma\ne0$.

The AMMN point comparison gives a $\mathbf Q_p$-linear isomorphism
\[
 \pi_{2n-1}K(W;\mathbf Q_p)
       \xrightarrow{\ \sim\ }
 HC_{2n-2}^{\mathrm{cts}}(W/W;\mathbf Q_p)
       =t^{1-n}E.
\]
For an actual class $\beta\in K_{2n-1}(W)$, denote its value by
$t^{1-n}r_n(\beta)$. If $r_n(\beta)\ne0$, then the actual class
\[
                  \eta\,\beta_B\in K_{2n-1}(B)
\]
has infinite order. More generally, if a finite family of actual
classes $\beta_j$ has $\mathbf Q_p$-linearly independent values
$r_n(\beta_j)$, then the classes $\eta(\beta_j)_B$ are linearly
independent in $K_{2n-1}(B)\otimes_{\mathbf Z}\mathbf Q_p$.
\end{lemma}

\begin{proof}
All occurrences of $K(-;\mathbf Q_p)$ denote rationalized
spectral $p$-completion, and write
$\Omega_B^j=\widehat\Omega^j_{B/W}[1/p]$.
Since $k$ is algebraic over $\mathbf F_p$,
Quillen's finite-field computation and filtered colimits give
\[
 K_{2m}(k)=0\quad(m>0),
 \qquad K_{2m-1}(k)\text{ is prime-to-}p\text{ torsion}.
\]
The coefficient exact sequences also show that $K(k)^{\wedge}_p$
has no positive homotopy groups. The point isomorphism in the
statement therefore follows from the fiber square of
\cite[Proposition~4.10]{AMMN};
\cite[Example~4.3]{AMMN} identifies the continuous K-theory of
$\operatorname{Spf}W$ with the spectral completion used here.

First handle one actual class $\beta$. Its reduction in
$K_{2n-1}(k)$ need not be zero. Choose an integer $N$, prime to $p$,
which kills that reduction. The actual relative exact sequence,
together with $K_{2n}(k)=0$, gives a unique actual lift
\[
 \kappa\in K_{2n-1}(W,\pi W)
       \quad\text{of}\quad N\beta.
\]
By compatibility with completion, its relative AMMN character is
\[
          \chi_W(\kappa)=N t^{1-n}r_n(\beta).
\]
This constructs an actual relative class before using its
completed character; no completed relative lift is being treated
as an actual lift.

The AMMN fiber comparison is $K(B)$-linear, by the module
construction through
\cite[Proposition~2.8, Corollary~2.10, and Theorem~2.12]{AMMN},
the relative-to-$W$ projection in Proposition~4.10, and the
multiplicative cyclotomic trace
\cite[Theorem~7.4]{BGTTrace}.
More explicitly, with $H=THH(B)$, the objects $H\otimes D$ and
arrows $\mathrm{id}_H\otimes f$ in that construction are
$H$-modules and module maps in cyclotomic spectra. The lax
monoidal circle and TC functors preserve the induced $TC(B)$
actions, and the equivalences used there are inverted in the
module category. The relative Hochschild projection is also
$H$-linear. Thus this assertion uses the actual module maps in
the construction, rather than the naturality of an unspecified
additive equivalence.
Its target action is the usual action of negative cyclic homology
on the suspended ordinary cyclic fiber. The norm compatibility is
\cite[Proposition~2.3]{BokstedtOttosen}, with the spectral norm
identified with the Connes operator by
\cite[Theorem~2.1, Lemma~2.2, and Theorem~2.3]{HoyoisCircle}.

The finite mixed-HKR model in relative dimension two gives
\[
 HC^-_0(B/W;\mathbf Q_p)
    =H^0_{\mathrm{dR}}(B/W)\oplus tH^2_{\mathrm{dR}}(B/W),
\]
where $H^2_{\mathrm{dR}}=\Omega_B^2/d\Omega_B^1$ is the ordinary
quotient of continuous forms. Rank zero removes the first
component of $\operatorname{ch}^-(\eta)$, and higher form degrees
vanish. Thus this character is exactly $t\gamma$.

The pullback of $\chi_W(\kappa)$ is the scalar cyclic class
$Nt^{1-n}r_n(\beta)$ on $B$. Scalar multiplication has no cyclic
shuffle correction: in the normalized relative bar complex any
such correction contains a bar occupied by this scalar in $E$,
and hence is zero. Consequently the actual relative product
\[
       \xi=\eta\,\kappa_B\in K_{2n-1}(B,\pi B)
\]
has character
\begin{equation}\label{d3:eq:formalhigherTCcupvalue}
       \chi_B(\xi)=N t^{2-n}r_n(\beta)\gamma
       \quad\text{in }HC_{2n-2}^{\mathrm{cts}}(B/W;\mathbf Q_p).
\end{equation}
There is no constant component. For $n\geq2$, the displayed
two-form component is an ordinary de Rham cohomology summand;
since $r_n(\beta)\ne0$ and $E$ is a field, it is nonzero.
Thus $\xi$ has infinite order.

Now use the actual relative exact sequence
\[
 K_{2n}(C)=0\longrightarrow K_{2n-1}(B,\pi B)
                    \longrightarrow K_{2n-1}(B).
\]
It is injective at the middle term, and the absolute image of
$\xi$ is exactly $N\eta\beta_B$. Hence $\eta\beta_B$ has infinite
order. No inverse to multiplication by $N$ on actual K-groups is used.

For a finite family, choose one common prime-to-$p$ integer $N$
killing all reductions, and form the unique actual lifts
$\kappa_j$. Tensoring the last injection with the flat
$\mathbf Z$-module $\mathbf Q_p$ preserves injectivity.
Applying the $\mathbf Q_p$-linear extension of $\chi_B$ to a
relation between the relative products gives
\[
       N\left(\sum_j q_jr_n(\beta_j)\right)\gamma=0.
\]
The assumed independence forces every $q_j=0$.
\end{proof}

\begin{remark}
For the good all-prime quadric
\[
 B=W\langle X,Y,Z\rangle/(XY+Z(Z+1)),
 \qquad I=(X,Z),\qquad\eta=1-[I],
\]
the special fiber is an affine-line torsor over $\mathbf P^1_k$.
Homotopy invariance, descent, and the projective-bundle formula
give $K(C)\simeq K(k)\oplus K(k)$, so the required actual even
K-group vanishes. The idempotent
\[
 e=\begin{pmatrix}-Z&X\\Y&1+Z\end{pmatrix}
\]
has image $I$ and satisfies
\[
 \operatorname{Tr}(e\,de\wedge de)
                  =\frac{dX}{X}\wedge dZ.
\]
Thus $\gamma$ is, up to the consistent first-Chern sign, the
nonzero ruling form detected by the root-chart residue lemma.
The same formal coefficient calculation works in every fixed
degree $2n-1$: changing $n$ only changes the cyclic column.
For $n=1$, by contrast, the prospective term $t\gamma$ is killed
in ordinary cyclic homology, so this argument gives no
degree-one nonvanishing statement.
\end{remark}

\begin{remark}
Let $B_g$ be the completed Laurent test algebra obtained from a
smooth algebraic model of $B$ by adjoining
$T_1^{\pm1},\ldots,T_s^{\pm1}$ and inverting a polynomial $g$
in those variables with coefficients in $W$. Its relative
dimension is $D=s+2$. For the actual relative lift $\kappa$ above,
the same module calculation gives
\[
 \chi_{B_g}\!\left(\eta\,\kappa_{B_g}
                         [T_1]\cdots[T_s]\right)_{\mathrm{top}}
   =N t^{2-n}r_n(\beta)\,
           \gamma\wedge d\log T_1\wedge\cdots\wedge d\log T_s
\]
in the $t^{2-n}H^D_{\mathrm{dR}}$ component of
$HC_{2n+s-2}^{\mathrm{cts}}(B_g/W;\mathbf Q_p)$.

When $n\geq3$, the starting cyclic column is negative, so the
positive columns cannot simply be discarded on the basis of the
ordinary cyclic quotient. Instead use the following chain-degree
calculation. Write a negative-cyclic cycle for $\eta$ as
$\sum_{r\geq1}t^r e_{2r}$ and one for each unit as
$\sum_{a_i\geq0}t^{a_i}u_{2a_i+1}$.
The first sum starts at $r=1$: its degree-zero term represents the
rank in $HH_0(B_g/W)=B_g$, and $b:C_1\to C_0$ is zero for a
commutative algebra. Hence rank zero forces that term itself to
vanish.

The scalar point cup has zero cyclic-shuffle correction. After
the remaining $s$ unit multiplications, a term using indices
$r,a_1,\ldots,a_s$ and $c$ cyclic-shuffle corrections has
Hochschild degree
\[
                   2r+s+2\sum_i a_i+2c.
\]
The finite normalized HKR map in dimension $D=s+2$ kills this
term unless $r=1$, all $a_i=0$, and $c=0$. Thus only the first
Chern component, the lowest unit Dennis traces, and ordinary
shuffles survive. Their normalized image is exactly the
displayed wedge, with no additional degree-dependent factor.
This argument allows the unwanted terms to survive in the
ordinary cyclic complex; it kills them by their final form degree.

For a test polynomial $g$ chosen as the same-support unit lift of
its residual polynomial, the integral continuous Laurent residue
maps this wedge to $\gamma$. It is therefore nonzero whenever
$\gamma\ne0$. If the special fiber of $B_g$ is smooth of dimension
$D$ over $k$, Geisser--Levine's finite-coefficient vanishing
\cite[Theorem~8.4]{GeisserLevineCharP} implies
\[
 \pi_{2n+s}K(B_g/\pi;\mathbf Q_p)=0,
\]
since $2n+s>D$. The completed relative-to-absolute map is thus
injective in degree $2n+s-1$. In this Laurent extension one uses
this completed injection; no actual even K-group vanishing of
the higher-dimensional special fiber is asserted.
\end{remark}

\begin{remark}
The lemma does not assert surjectivity from actual K-theory to
spectral completion, and it does not identify $r_n$ with a
Bloch--Kato, syntomic, or polylogarithmic coordinate. Independently
known actual classes with independent \'etale Chern values may be
used once the separate factorization through spectral completion
has been established. The hypothesis on $k$ in this lemma is
algebraicity over $\mathbf F_p$; perfectness alone does not imply
the actual even K-group vanishing used in its proof.
\end{remark}


\subsection{Unramified actual coefficients at every prime}
\label{d3:sec:unramifiedhighercoefficients}

\begin{proposition}
\label{d3:prop:unramifiedhighercoefficients}
Let $p$ be any prime, $m\geq2$, and put
\[
 k_m=\F_{p^m},\qquad V_m=W(k_m),\qquad F_m=\Frac V_m.
\]
Take $u=i$ if $p=2$, and take $u$ to be a primitive $p$-th root
if $p>2$. Put $E_m=F_m(u)$, $O_m=V_m[u]$, and let $\sigma$
fix $F_m$ and invert $u$.
For each $n\geq2$ there exist Teichm\"uller roots
$\zeta_1,\ldots,\zeta_{m-1}\in\mu_{p^m-1}(V_m)\setminus\{1\}$
such that the actual rational classes
\[
 \beta_{n,j}=b_n(\zeta_ju)-b_n(\zeta_ju^{-1})
 \in K_{2n-1}(O_m)_{\Q},\qquad 1\leq j\leq m-1,
\]
are $\sigma$-anti-invariant and have $\Q_p$-linearly independent
images in $H^1(E_m,\Q_p(n))$ under $c_{n,1}$.
For each fixed $p,m,n$, a common integer multiple is represented
by strictly anti-invariant integral classes.
\end{proposition}
\begin{proof}
Put $\delta=u-u^{-1}\ne0$, $\epsilon=u^p=u^{-p}$, and
$F_{p,n}(z)=\operatorname{Li}_n(z)-p^{-n}\operatorname{Li}_n(z^p)$.
The equal $p$th powers make the depleted terms cancel exactly:
\[
 \operatorname{Li}_n(\zeta u)-\operatorname{Li}_n(\zeta u^{-1})
 =F_{p,n}(\zeta u)-F_{p,n}(\zeta u^{-1}).
\]
For $1\leq a\leq p-1$ the Laurent polynomial
\[
 Q_a(u)=\frac{u^a-u^{-a}}{u-u^{-1}}
       =\sum_{b=0}^{a-1}u^{a-1-2b}
\]
is integral and reduces to $a$.
Grouping denominators as $a+p\ell$ and using the binomial series
gives, with $S_j(t)=P_j(t)/(1-t)^{j+1}$ and $P_j\in\Z[t]$,
\begin{equation}\label{d3:eq:unramifiedhighercoordinate}
\begin{split}
 a_{n,p}(\zeta)
 &\coloneqq\frac{\operatorname{Li}_n(\zeta u)
                   -\operatorname{Li}_n(\zeta u^{-1})}{\delta}\\
 &=\sum_{j\geq0}(-1)^j\binom{n+j-1}{j}p^j
       \sum_{a=1}^{p-1}\frac{\zeta^aQ_a(u)}{a^{n+j}}
                              S_j(\epsilon\zeta^p)\in O_m.
\end{split}
\end{equation}
Here $S_0(t)=1/(1-t)$.
The series identity first holds for $|z|<1$ and extends uniformly
to $|z|\leq1$, $v_p(1-z^p)\leq2/3$.
The rigid identity principle and
\cite[Theorem~2.7(3)]{BesserDeJeuAlgorithm} identify it with Coleman
continuation. At the displayed points $1-\epsilon\zeta^p$ is a unit,
so every $j$th summand has valuation at least $j$.

If $t=\overline\zeta$, reduction of
\eqref{d3:eq:unramifiedhighercoordinate} is
\[
 h_n(t)=\frac{\sum_{a=1}^{p-1}a^{1-n}t^a}{1-t^p},
 \qquad t\in k_m^\times\setminus\{1\}.
\]
The numerator has $t$-coefficient one and degree $p-1$.
Thus $h_n$ is a nonconstant rational map of degree at most $p$.
If $p>2$, each finite fiber has at most $p$ elements, so
\[
 \#h_n(k_m^\times\setminus\{1\})
 \geq\left\lceil\frac{p^m-2}{p}\right\rceil=p^{m-1}.
\]
Its $\F_p$-span has dimension at least $m-1$.
If $p=2$, then $h_n(t)=t/(1+t)^2=w^2+w$, where
$w=(1+t)^{-1}$ ranges through $k_m\setminus\{0,1\}$.
Its image is exactly the nonzero part of
$\ker\operatorname{Tr}_{k_m/\F_2}$, whose span has dimension $m-1$.
Choose $m-1$ independent image values and Teichm\"uller lifts
$\zeta_j$ of corresponding inputs.
A nonzero $\Q_p$-linear relation among $a_{n,p}(\zeta_j)$ can
be normalized to coefficients in $\Z_p$, with one a unit.
Reduction would contradict the chosen $\F_p$-independence.
The coordinates are therefore $\Q_p$-independent; they need not
belong to $F_m$.

These are actual number-field coefficients. For $N=p^m-1$ put
$L_m=\Q(\mu_N,u)$.
The order of $p$ modulo $N$ is $m$, so the chosen completion
of $\Q(\mu_N)$ is $F_m$, and the completion of $L_m$ is $E_m$.
The shifted cyclotomic polynomial is Eisenstein in the ramified
factor, so its integer ring is $O_m$.
Theorem~1.12 of \cite{BesserDeJeu} gives $b_n(\zeta_ju)$ in
$K_{2n-1}(L_m)_{\Q}$.
Localization at the selected prime gives its rational lift to the
local integer ring, since the positive residue-field $K$-groups
are torsion. Pullback gives a class over $O_m$.
The global automorphism fixing $\mu_N$ and inverting $u$ preserves
the selected prime and induces $\sigma$.
Naturality proves exact anti-invariance of $\beta_{n,j}$.
Clearing denominators of the individual $b_n(\zeta_ju)$, then
subtracting their integral $\sigma$-images, gives the asserted
integral representatives.

Finally the common normalization
\eqref{d3:eq:highercyclotomicregulator} gives
\[
 c_{n,1}(\beta_{n,j})
 =\exp_{E_m,n}\bigl(\kappa_n\delta\,a_{n,p}(\zeta_j)\bigr).
\]
For $n\geq2$ this exponential is a $\Q_p$-linear isomorphism.
The independence of the coordinates and $\kappa_n\delta\ne0$
prove the assertion. No actual-to-completed $K$-theory
surjectivity is used.
\end{proof}

This proposition is a finite-level coefficient statement.
Its use over one fixed ring with infinite perfect residue field
requires the compatible ruling detector after that base change.


\subsection{One strict henselian ring in every odd degree}
\label{d3:sec:strictallprimeallodd}

Put $k=\overline{\F}_p$, $V_\infty=W(k)$, and define
\[
 f_p(z)=
 \begin{cases}
  \Phi_p(1+z),&p>2,\\
  \Phi_4(1+z)=z^2+2z+2,&p=2.
 \end{cases}
\]
The polynomial has Eisenstein degree $p-1$ for odd $p$ and degree
two for $p=2$.

\begin{theorem}\label{d3:thm:strictallprimeallodd}
For every prime $p$, the rings
\[
 R_{\infty,p}=V_\infty[[x,y,z]]/(xy+f_p(z)),\qquad
 S_p=\left(
  \bigl(\Z_{(p)}[x,y,z]/(xy+f_p(z))\bigr)_{(p,x,y,z)}
       \right)^{\mathrm{sh}}
\]
are strictly henselian regular local rings of dimension three,
and $\widehat S_p\cong R_{\infty,p}$ after fixing their residue
identifications with $k$.
For every $n\geq2$, each of these two rings $R$ satisfies
\[
 \dim_{\Q}\ker\left(
 K_{2n-1}(R)_{\Q}\longrightarrow
 K_{2n-1}(\Frac R)_{\Q}\right)=\infty.
\]
More precisely, for every $m\geq2$ the kernel contains $m-1$
rationally independent classes represented by actual integral
$K$-classes whose restrictions to $R[1/x]$ are zero integrally.
The ring for a fixed $p$ is independent of both $m$ and $n$.
\end{theorem}

\begin{proof}
Fix $n\geq2$ and $m\geq2$, and write $q=2n-1$,
$V_m=W(\F_{p^m})$, and $F_m=\Frac V_m$.
Take $\rho=i$ for $p=2$, and a primitive $p$th root for $p>2$.
Put $W_m=V_m[\rho]$, $E_m=F_m(\rho)$ and
$W_\infty=V_\infty[\rho]$, $E_\infty=\Frac W_\infty$.
These are complete coefficient DVRs and their fraction fields;
$\rho-1$ is a uniformizer in each ramified DVR.
Let $\sigma$ fix the unramified coefficient ring and invert $\rho$.

Proposition~\ref{d3:prop:unramifiedhighercoefficients} supplies
$m-1$ actual coefficients $\beta_j\in K_q(W_m)$, after one
common integer multiplication, with $\sigma\beta_j=-\beta_j$
and independent $\Q_p$-valued \'etale Chern classes.
Lemma~\ref{d3:lem:cherncompletionfactor} implies that their images
in $\pi_qK(W_m;\Q_p)$ are $\Q_p$-independent.
Under the AMMN point isomorphism used in
Lemma~\ref{d3:lem:formalhigherTCcup}, denote their coordinates by
\[
 a_j=r_{n,W_m}(\beta_j)\in E_m.
\]
These coordinates are independent and $\sigma a_j=-a_j$.
Naturality of that point calculation under $W_m\to W_\infty$
sends them by the field inclusion $E_m\hookrightarrow E_\infty$,
preserving independence. This argument only transports independence;
it does not identify the AMMN coordinates with a prescribed
polylogarithm normalization.

Set $R_m=V_m[[x,y,z]]/(xy+f_p(z))$.
Weierstrass division identifies its regular divisor $x=0$ with
$W_m[[y]]$, by $1+z\mapsto\rho$.
Pull back $\beta_j$ and push forward along this divisor to obtain
$g_j\in K_q(R_m)$. Localization makes $g_j|_{R_m[1/x]}=0$
integrally.
Both $R_m$ and $R_{\infty,p}$ are regular local of dimension three:
the equation has linear term $p$ in the regular ambient ring,
so $x,y,z$ are regular parameters.
The closed fiber of $R_m\to R_{\infty,p}$ is $k$; hence the map
is flat by \cite[Tag~00R4]{StacksFlatness}.
Its pullback therefore preserves the supported pushforward.

After the finite free extension $V_\infty\to W_\infty$ and
inversion of $p$, the polynomial splits with distinct roots
$a_g=g(\rho)-1$, indexed by
$G=\operatorname{Gal}(E_m/F_m)$.
Write $I_g=(x,z-a_g)$ and $\eta_g=1-[I_g]$ on this generic ring.
The root components of $x=0$ are disjoint Cartier divisors;
their structure-sheaf classes sum to zero because their union is
principal. Flat base change and the projection formula give
\[
 g_j=\sum_{g\in G}g(\beta_j)\eta_g.
\]
It suffices to retain the fixed inverse-root pair.
Put $d=\rho-\rho^{-1}\ne0$,
$H(z)=\prod_{g\ne1,\sigma}(z-a_g)$, and substitute
\[
 z=(\rho-1)+dZ,\qquad x=dX,\qquad
 y=dH((\rho-1)+dZ)Y
\]
into the smooth formal quadric
$B=W_\infty\langle X,Y,Z\rangle/(XY+Z(Z+1))$.
All original formal variables are topologically nilpotent, so the
substitution defines a ring map. It uses $z=t-1$, including at two.
The two selected root ideals pull to $J=(X,Z)$ and $J'=(X,Z+1)$,
up to the invertible generic scalar $d$.
Every omitted root ideal becomes trivial: if
$b=(a_g-a_1)/d\notin\{0,-1\}$, then
\[
 b(b+1)=-XY-(Z-b)(Z+b+1)\in(X,Z-b).
\]
No inversion of $H$ or flatness of this last substitution is used.
The pulled invertible module surjects onto the displayed invertible
ideal, hence is isomorphic to it.
The two remaining structure-sheaf classes sum to zero, so on
$B[1/p]$ the class becomes exactly
$2\beta_j(1-[J])$.
Lemma~\ref{d3:lem:formalhigherTCcup} detects these products by their
independent coordinates $2a_j$ in $K_q(B)$.
Here $K_q(B/\pi)$ is torsion, since the special fiber has the
$K$-theory of $\mathbf P^1_k$ and $k$ is algebraic over $\F_p$.
Localization therefore makes $K_q(B)_{\Q}\to K_q(B[1/p])_{\Q}$
injective, so these products remain independent on the generic
test ring.
Thus the $g_j$ are rationally independent on $R_{\infty,p}$.

For the arithmetic assertion choose $N=p^m-1$ and a prime $v$
above $p$ in $\Q(\mu_N)$ compatible with $k$.
The essentially \'etale local DVR
$V_{m,0}=(\Z_{(p)}[\zeta_N])_v$ has completion $V_m$, since
$p$ has multiplicative order $m$ modulo $N$; put $W_{m,0}=V_{m,0}[\rho]$.
The coefficients in Proposition~\ref{d3:prop:unramifiedhighercoefficients}
come from the actual rational $K_q$ of
$\Frac W_{m,0}=\Q(\mu_N,\rho)$.
The localization map from $K_q(W_{m,0})$ is surjective
\cite[Sections~5--7]{Quillen}, since
$K_{q-1}(\F_{p^m})=0$ by
\cite{QuillenFinite}\cite[IV, Corollary~1.13]{Weibel}.
Clear denominators before lifting and applying $1-\sigma$.
This produces the same independent finite family from actual
classes on $W_{m,0}$.
The chosen pointed \'etale neighborhood gives $V_{m,0}\to S_p$,
and $\rho\mapsto1+z$ gives
$W_{m,0}\to S_p/(x)$.
Pull the coefficients to this regular divisor and push to $S_p$.
They vanish after inverting $x$ integrally.

The universal property gives a compatible local map
$S_p\to R_{\infty,p}$.
Both rings have residue $k$ and regular parameters $x,y,z$;
the induced associated-graded map is the identity on $k[x,y,z]$.
It follows that $\widehat S_p\cong R_{\infty,p}$.
Completion is flat, so these arithmetic supported classes map to
the already detected completed classes and are independent.
Letting $m$ grow proves both assertions.
\end{proof}


\section{One fixed ring in every degree at least three}
\label{d3:sec:singlefixedalldegree}

We use the ordinary continuous de Rham quotient throughout.
First we record why introducing one Laurent variable preserves both
the coefficient--ruling products and their products with that unit.

\begin{lemma}\label{d3:lem:onelaurentcoefficientdetector}
Let $W$ be the integers of a finite extension $E/\Q_p$, with
residue field $k$ and uniformizer $\pi$. Put
\[
 B=W\langle X,Y,Z\rangle/(XY+Z(Z+1)),\qquad
 J=(X,Z),\qquad \eta=1-[J].
\]
Let $0\ne g\in W[T]$ have unit nonzero coefficients, and put
$B_g=(B[T^{\pm1},1/g])^\wedge_\pi$.
For $n\geq2$, suppose actual classes
$\beta_1,\ldots,\beta_r\in K_{2n-1}(W)$ have
$\Q_p$-linearly independent AMMN point coordinates
$a_1,\ldots,a_r\in E$.
Then each of the two families
\[
 \eta\beta_j\in K_{2n-1}(B_g[1/p])_{\Q},\qquad
 \eta\beta_j[T]\in K_{2n}(B_g[1/p])_{\Q}
\]
is rationally independent.
\end{lemma}

\begin{proof}
Choose a common integer $N$, prime to $p$, which kills the
reductions of the $\beta_j$ in $K_{2n-1}(k)$.
Their unique actual relative lifts have point characters
$Nt^{1-n}a_j$, where $|t|=-2$.
Write $\gamma=\operatorname{ch}^-_1(\eta)$ for the nonzero ruling
two-form class of Lemma~\ref{d3:lem:formalhigherTCcup}.
The module construction of that lemma, in the finite mixed-HKR
model of relative dimension three, gives the relative characters
\begin{equation}\label{d3:eq:onelaurentcoefficientcharacters}
 Nt^{2-n}a_j\gamma,\qquad
 Nt^{2-n}a_j\gamma\wedge d\log T.
\end{equation}
Here is the degree check for the second formula, including $n>2$.
The rank-zero negative Chern cycle has terms $t^r e_{2r}$ with
$r\geq1$; a unit cycle has terms $t^a u_{2a+1}$ with $a\geq0$.
A cyclic-shuffle correction increases Hochschild degree by two.
Dimension three therefore retains only $r=1$, $a=0$, and the
ordinary shuffle. Normalized HKR gives exactly the second term
in \eqref{d3:eq:onelaurentcoefficientcharacters}.
This argument does not discard positive cyclic columns
merely because the initial column might be negative.

The degree-zero projection $P_0$ of
Lemma~\ref{d3:lem:continuouslaurentresidue} sends the first form
coefficient in \eqref{d3:eq:onelaurentcoefficientcharacters}
back to $Na_j\gamma$; its residue sends the second coefficient there.
The independence of the $a_j$ thus detects every nonzero rational
linear combination of either relative family.
The special fiber $C_g=B_g/\pi$ is smooth of dimension three.
Lemma~\ref{d3:lem:laurentspecialvanishing} gives
$\pi_{q+1}K(C_g;\Q_p)=0$ for every $q\geq3$.
Consequently the relative-to-absolute map is injective in all the
degrees used here after spectral completion and rationalization.
This proves independence in the actual groups of $B_g$.

The final generic localization requires a separate actual
$K$-theory observation. Let
$U_g=\Spec k[T^{\pm1},1/\bar g]$.
The ruling and projective-bundle formula give
$K(C_g)\simeq K(U_g)\oplus K(U_g)$.
For $q\geq2$, localization on $\A^1_k$
\cite[Sections~5--7]{Quillen} puts $K_q(U_g)$ between
a quotient of $K_q(k)$ and a subgroup of the finite direct sum
$\bigoplus_vK_{q-1}(k(v))$, where the omitted points $v$ have
finite residue fields. These groups are finite by
\cite{QuillenFinite}\cite[IV, Corollary~1.13]{Weibel}.
Thus $K_q(C_g)$ is torsion, and localization makes
\[
 K_q(B_g)_{\Q}\longrightarrow K_q(B_g[1/p])_{\Q}
\]
injective. This proves the stated generic assertions.
\end{proof}

\begin{theorem}\label{d3:thm:singlefixedalldegree}
Fix a prime $p$, let $S=\Z_{(p)}^{\mathrm{sh}}$ with residue
$k=\overline\F_p$, and set
\[
 V_0=(S[T])_{(p)},\qquad
 f_p(z)=
 \begin{cases}
  \Phi_p(1+z),&p>2,\\
  \Phi_4(1+z)=z^2+2z+2,&p=2,
 \end{cases}
\]
\[
 R_p=\bigl(V_0[x,y,z]/(xy+f_p(z))\bigr)_{(x,y,z)}.
\]
This is one Noetherian regular local ring of dimension three,
with residue field $k(T)$, and for every $j\geq3$,
\[
 \dim_{\Q}\ker\left(K_j(R_p)_{\Q}
            \longrightarrow K_j(\Frac R_p)_{\Q}\right)=\infty.
\]
Each of the finite independent families used to establish this
assertion consists of actual integral classes which vanish after
inverting $x$ integrally.
\end{theorem}

\begin{proof}
The strict henselization $S$ is a Noetherian DVR with uniformizer
$p$; see \cite[Tags~06LJ and~0AP3]{StacksHenselization}.
The localization $V_0$ is a DVR with the same uniformizer and
residue field $k(T)$.
The polynomial $f_p$ is Eisenstein of degree at least two.
In the regular ambient local ring of dimension four, the equation
has nonzero linear term $p$, proving the assertions about $R_p$.
Equivalently,
\begin{equation}\label{d3:eq:singlefixedringlocalization}
 R_p=\bigl(S[T,x,y,z]/(xy+f_p(z))\bigr)_{(p,x,y,z)}.
\end{equation}

Fix $n\geq2$ and $m\geq2$, and take $\rho=i$ at two and a
primitive $p$th root otherwise.
Choose a compatible prime $v$ of $\Q(\mu_{p^m-1})$ and put
\[
 V_{m,0}=(\Z_{(p)}[\zeta_{p^m-1}])_v\subset S,
 \qquad W_{m,0}=V_{m,0}[\rho].
\]
The actual number-field classes of
Proposition~\ref{d3:prop:unramifiedhighercoefficients} lift to
$K_{2n-1}(W_{m,0})$ after a common integer multiplication:
the obstruction is $K_{2n-2}(\F_{p^m})=0$.
Choose the integral lifts before applying $1-\sigma$, where
$\sigma$ fixes $V_{m,0}$ and inverts $\rho$.
We obtain $m-1$ strictly anti-invariant classes $\beta_j$.
Their finite-local-field \'etale Chern classes are independent;
Lemma~\ref{d3:lem:cherncompletionfactor} and the AMMN point isomorphism
therefore give independent point coordinates in
$E_m=\Frac W(\F_{p^m})(\rho)$.

The regular divisor $x=0$ in $R_p$ receives $W_{m,0}$ by
$\rho\mapsto1+z$.
Pull back and push forward to obtain $a_j\in K_{2n-1}(R_p)$,
and put $b_j=a_j[T]\in K_{2n}(R_p)$.
Here $T$ is a unit. Both families vanish on $R_p[1/x]$
integrally, by localization and the projection formula.

We check independence after any putative finite relation.
Using \eqref{d3:eq:singlefixedringlocalization} and
filtered-ring continuity of $K$-theory
\cite[IV, Section~6.4]{Weibel}, the coefficients, the classes,
and an actual integral relation all descend to
\[
 \bigl(V_{\lambda,0}[T,x,y,z]/(xy+f_p(z))\bigr)[1/F],
 \qquad \bar F(T)\ne0,
\]
where $V_{\lambda,0}\subset S$ is an essentially \'etale local
DVR containing $V_{m,0}$ and all constants in the finite data.
We include $T$ in the inverted product $F$ when necessary.
Its completion is an unramified coefficient DVR $V_\ell$ with
finite residue field; put $W_\ell=V_\ell[\rho]$ and
$E_\ell=\Frac W_\ell$.
Point-character naturality sends the independent coordinates by
the field inclusion $E_m\hookrightarrow E_\ell$.

Over $W_\ell$, retain the two roots
$\rho-1,\rho^{-1}-1$ of $f_p$.
With $d=\rho-\rho^{-1}$ and
$H(z)=\prod_{a\ne\rho-1,\rho^{-1}-1}(z-a)$, use
\[
 z=(\rho-1)+dZ,\qquad x=dX,\qquad
 y=dH((\rho-1)+dZ)Y.
\]
Choose $g(T)$ by lifting precisely the nonzero coefficients of
$\bar F(T)$ to units in $W_\ell$ and form $B_g$ as in
Lemma~\ref{d3:lem:onelaurentcoefficientdetector}.
All three displayed images are topologically nilpotent.
The image of $F$ is $g+\pi h$, hence a unit of $B_g$.
Thus this finite relation has an actual image in the test ring;
no map from all of $R_p$ to a fixed test algebra is asserted.

Flat coefficient base change and the projection formula, followed
by inversion of $p$, split the supported class as
$\sum_\tau\tau(\beta_j)(1-[I_\tau])$.
The pair substitution makes every omitted root ideal trivial.
The two remaining disjoint components have structure-sheaf
classes summing to zero, and anti-invariance therefore gives
\[
 a_j\longmapsto2\eta\beta_j,\qquad
 b_j\longmapsto2\eta\beta_j[T]
 \quad\text{in }K_*(B_g[1/p]).
\]
Lemma~\ref{d3:lem:onelaurentcoefficientdetector} excludes the asserted
relation in either degree.
Hence both families have rational rank $m-1$ in $R_p$.
As $m$ is arbitrary and the degrees $2n-1,2n$ with $n\geq2$
cover every integer at least three, the theorem follows.
\end{proof}


\subsection{A single complete ring in every degree at least three}
\label{d3:sec:completedfixedalldegree}

The finite-witness proof of Theorem~\ref{d3:thm:singlefixedalldegree}
does not itself pass to arbitrary completed coefficients.  We now
give a complete-ring construction, using
Lemma~\ref{d3:lem:essentiallysmoothAMMN} and a constant-term map
defined on all primitive denominators.

\begin{lemma}
\label{d3:lem:completedcoefficientresidue}
Let $W$ be a complete DVR with uniformizer $\pi$, and set
\[
 V_W=\widehat{W[T]_{(\pi)}}^{\,\pi},\qquad
 B'=W\langle X,Y,Z\rangle/(XY+Z(Z+1)),
\]
\[
 \mathcal B=V_W\langle X,Y,Z\rangle/(XY+Z(Z+1)).
\]
There are continuous $W$-linear maps on the relative form complexes,
a degree-zero chain projection $P_0$ and a degree-minus-one residue,
which fix base forms and satisfy
\[
 P_0(\omega)=\omega,\qquad
 \operatorname{Res}_T(\omega\wedge d\log T)=\omega
 \quad(\omega\in\widehat\Omega^\bullet_{B'/W}).
\]
They preserve actual de Rham boundaries. The complexes here retain
the differential $dT$: their base ring is $W$, not $V_W$.
\end{lemma}

\begin{proof}
For $a\geq1$, put $L_a=(W/\pi^a)((T))$.
Every $g\in W[T]\setminus(\pi)$ is a unit in $L_a$.
Indeed its nonzero reduction is invertible in the field
$k((T))$, and inverses lift across the nilpotent ideal $(\pi)$.
Uniqueness makes these inverses compatible in $a$. Hence
\[
 W[T]_{(\pi)}\longrightarrow\varprojlim_a L_a
\]
is an integral, $\pi$-adically continuous ring homomorphism.
Constant-term extraction is compatible with these reduction maps,
so extends to
\[
 \operatorname{CT}:V_W\longrightarrow W,\qquad
 \operatorname{CT}(1)=1,\qquad
 \operatorname{CT}(T\partial_T f)=0.
\]
The identity with the logarithmic derivative is coefficientwise
at every level and then passes to the inverse limit.

Tensor with the quadric algebra before completion and use
$\varprojlim_a(B'/\pi^a)((T))$ as target.  The completed relative
forms decompose into a component with no $d\log T$ and one
with $d\log T$ at the right. Apply constant term to the first
to define $P_0$, and to the second to define
$\operatorname{Res}_T$.
The derivative identity gives compatibility with the de Rham
differentials (with the conventional shift sign for a residue).
In particular an actual primitive maps to an actual primitive.
All operations are integral and bounded before inverting $p$;
there is no completion of the image of the differential.
\end{proof}

\begin{theorem}
\label{d3:thm:completedfixedalldegree}
For a prime $p$, put $k=\overline{\F}_p$, $W_0=W(k)$, and
\[
 V=\widehat{\,W_0[T]_{(p)}\,}^{\,p},\qquad
 \mathcal R_p=V[[x,y,z]]/(xy+f_p(z)),
\]
where $f_p(z)=\Phi_p(1+z)$ for $p>2$, and
$f_2(z)=\Phi_4(1+z)=z^2+2z+2$.
Then $\mathcal R_p$ is a complete Noetherian regular local
domain of dimension three, with residue field $k(T)$.
For every $j\geq3$,
\[
 \dim_{\Q}\ker\left(K_j(\mathcal R_p)_{\Q}
          \longrightarrow K_j(\Frac\mathcal R_p)_{\Q}\right)
 =\infty.
\]
The finite independent families proving this assertion consist
of actual integral classes with zero image in
$K_j(\mathcal R_p[1/x])$.
\end{theorem}

\begin{proof}
The ring $W_0[T]_{(p)}$ is a Noetherian DVR with uniformizer
$p$ and residue $k(T)$. Its completion $V$ has the same
properties. Since $f_p$ is Eisenstein of degree at least two,
the quotient in the statement has nonzero linear term $p$
in a regular ambient local ring of dimension four.
It is therefore regular of dimension three, with maximal
ideal $(x,y,z)$ and the asserted residue field.

Let $\rho=i$ for $p=2$, and otherwise let $\rho$ be a primitive
$p$th root of unity. Put
\[
 W=W_0[\rho],\qquad \pi=\rho-1,\qquad E=\Frac W.
\]
There are compatible identifications
\begin{equation}\label{d3:eq:completedfixedcoefficients}
 W_0[T]_{(p)}[\rho]=W[T]_{(\pi)},\qquad
 V[\rho]=\widehat{W[T]_{(\pi)}}^{\,\pi}.
\end{equation}
To verify the first, the norm of a polynomial
$g\in W[T]\setminus(\pi)$ reduces to $\bar g^{[W:W_0]}$
and is primitive over $W_0$. The norm and adjugate identity
therefore invert $g$ in the left side. The reverse inclusion
is immediate. Finite freeness over $W_0$ gives the second
identity after completion; the $p$-adic and $\pi$-adic
topologies are cofinal.

Fix $n\geq2$ and an arbitrary $m\geq2$.
Proposition~\ref{d3:prop:unramifiedhighercoefficients} and
Lemma~\ref{d3:lem:cherncompletionfactor} supply $m-1$ actual
classes $\beta_1,\ldots,\beta_{m-1}$ over
$W_m=W(\F_{p^m})[\rho]$, exactly anti-invariant under
$\sigma(\rho)=\rho^{-1}$, with independent AMMN point
coordinates. Their images in $K_{2n-1}(W)$ have coordinates
$b_1,\ldots,b_{m-1}\in E$ still independent over $\Q_p$:
the point character is natural under the coefficient inclusion
$\Frac W_m\hookrightarrow E$.
No surjectivity from actual $K$-groups to all point coordinates
is asserted.

The regular divisor $x=0$ of $\mathcal R_p$ has coordinate ring
\[
 \mathcal R_p/(x)\simeq V[\rho][[y]],
 \qquad \rho\longmapsto1+z.
\]
Pull back the $\beta_\nu$ and push forward to obtain
$a_\nu\in K_{2n-1}(\mathcal R_p)$.
Also put $a_\nu[T]\in K_{2n}(\mathcal R_p)$.
These actual classes vanish after inverting $x$.

Consider the single test algebra
\[
 \mathcal B=V[\rho]\langle X,Y,Z\rangle/(XY+Z(Z+1)),
 \qquad J=(X,Z),\qquad\eta=1-[J].
\]
By \eqref{d3:eq:completedfixedcoefficients} it is the completion
of an essentially smooth $W$-algebra of relative dimension
three. Lemma~\ref{d3:lem:essentiallysmoothAMMN} applies with
this perfect-constant coefficient ring $W$.
Its continuous de Rham model includes the variable $T$.

Retain the inverse pair of roots
$a=\rho-1$, $a'=\rho^{-1}-1$ of $f_p$.
Set $d=\rho-\rho^{-1}$ and
$H(z)=\prod_{\alpha\ne a,a'}(z-\alpha)$.
The substitution
\begin{equation}\label{d3:eq:completedfixedpairmap}
 z=a+dZ,\qquad x=dX,\qquad y=dH(a+dZ)Y
\end{equation}
defines an actual map $\mathcal R_p\to\mathcal B$.
The equation becomes $d^2H(a+dZ)(XY+Z(Z+1))=0$.
All three images lie in $\pi\mathcal B$, so arbitrary formal
power series converge. The coefficient map is the actual
inclusion $V\to V[\rho]$; no finite-stage descent of completed
coefficients is required.

Flat coefficient base change first splits the divisor after
inverting $p$. Pulling its root-ideal line bundles along
\eqref{d3:eq:completedfixedpairmap} leaves only $J$ and
$J'=(X,Z+1)$ on $\mathcal B[1/p]$.
For an omitted root with normalized value
$b=(\alpha-a)/d\notin\{0,-1\}$, the ideal
$(X,Z-b)$ is the unit ideal, since
\[
 b(b+1)=-XY-(Z-b)(Z+b+1).
\]
The two surviving disjoint components have classes summing
to zero, because their union is the principal divisor $X=0$.
The projection formula and exact anti-invariance therefore give
\begin{equation}\label{d3:eq:completedfixedimages}
 a_\nu\longmapsto2\beta_\nu\eta,\qquad
 a_\nu[T]\longmapsto2\beta_\nu\eta[T]
 \quad\text{on }\mathcal B[1/p].
\end{equation}

It remains to check these classes in actual groups.
The reduction of each $\beta_\nu$ lies in
$K_{2n-1}(k)$, a torsion group of order prime to $p$
element by element. Choose one common prime-to-$p$ integer
$N$ killing the finite list of reductions.
Since $K_{2n}(k)=0$, each $N\beta_\nu$ has a unique actual
relative lift at the coefficient point.
Its character is $Nt^{1-n}b_\nu$.
The module comparison and the finite dimension-three HKR
calculation of Lemma~\ref{d3:lem:formalhigherTCcup} give
\[
 Nt^{2-n}b_\nu\gamma,\qquad
 Nt^{2-n}b_\nu\gamma\wedge d\log T,
\]
where $\gamma$ is the nonzero ordinary ruling two-form class
on $B'=W\langle X,Y,Z\rangle/(XY+Z(Z+1))$.
Lemma~\ref{d3:lem:completedcoefficientresidue} sends these
respectively to $Nb_\nu\gamma$.
The independence of the $b_\nu$ detects every nonzero
rational combination of either actual relative family.

Finally,
\[
 C=\mathcal B/\pi
   =k(T)[X,Y,Z]/(XY+Z(Z+1)),\qquad
 K(C)\simeq K(k(T))\oplus K(k(T)).
\]
For $q\geq2$, the actual group $K_q(k(T))$ is torsion.
Every element descends first to a finite field of constants
and then to a finite open of its affine line.
Quillen localization expresses the $K_q$-group of that open
as an extension of a quotient of the finite group $K_q(k_m)$
by a subgroup of a finite sum of $K_{q-1}(k(v))$.
These groups are finite, including at $q=2$.
Continuity proves the claim for $k(T)$.

The algebra $\mathcal B$ is Noetherian and flat over the complete
DVR $V[\rho]$, and its special fiber is smooth.
Consequently $\pi$ is a nonzerodivisor and $\mathcal B$ is regular:
at every maximal ideal the regular quotient by $\pi$ gives the
regularity criterion, and $\pi$ belongs to the Jacobson radical.
Thus $K_{q+1}(C)_{\Q}=K_q(C)_{\Q}=0$ for every $q\geq3$.
The actual relative sequence and the regular-divisor
localization sequence respectively give injections
\[
 K_q(\mathcal B,\pi)_{\Q}\hookrightarrow K_q(\mathcal B)_{\Q},
 \qquad
 K_q(\mathcal B)_{\Q}\hookrightarrow
                     K_q(\mathcal B[1/p])_{\Q}.
\]
The detected families therefore remain independent in the
generic test ring. Equation~\eqref{d3:eq:completedfixedimages}
proves their independence on $\mathcal R_p$.
Their ranks are at least $m-1$ for arbitrary $m$, and the
degrees $2n-1$ and $2n$ cover every integer at least three.
\end{proof}


\section{Milnor, motivic, and higher-Chow consequences}
\label{d3:sec:milnormotivicconsequences}

Put
\[
 A=\bigl(\mathbf Z_{(3)}[x,y,z]/(3+xy+z^2)\bigr)_{(x,y,z)},
 \qquad R\in\{A,A^h,\widehat A,A^{\mathrm{sh}}\},
 \qquad F_R=\operatorname{Frac}(R).
\]
In each ring use the units
\[
 a=\frac{1+z}{2},\qquad \lambda=1+x,\qquad
 \mu=1+\frac{(1+x)y}{4},\qquad
 s_0=a(1-a)=1+\frac{xy}{4}.
\]
Thus $1+x\mu=\lambda s_0$ and
$c_R=2[a][\lambda][\mu]\in K_3(R)$ is the explicit class already
constructed. The superscript $M,\mathrm{naive}$ below means the
tensor algebra on $R^\times$ modulo the two-sided Steinberg ideal;
graded commutativity is not added to its definition.

\begin{corollary}\label{d3:cor:explicitmilnormotivic}
For every such $R$, each of the four generic restriction maps
\[
\begin{gathered}
 K_3^{M,\mathrm{naive}}(R)_{\mathbf Q}
       \longrightarrow K_3^M(F_R)_{\mathbf Q},\qquad
 \widehat K_3^M(R)_{\mathbf Q}
       \longrightarrow K_3^M(F_R)_{\mathbf Q},\\
 H_B^3(R,\mathbf Q(3))\longrightarrow H_B^3(F_R,\mathbf Q(3)),
 \qquad
 CH^3(R,3)_{\mathbf Q}\longrightarrow CH^3(F_R,3)_{\mathbf Q}
\end{gathered}
\]
has a nonzero kernel. Explicit witnesses are, respectively,
$2\{a,\lambda,\mu\}$, its image in improved Milnor theory,
$2[a]_B\cup[\lambda]_B\cup[\mu]_B$, and twice the cubical graph
cycle specified below. The Milnor symbols and the graph cycle vanish
over $F_R$ integrally.
\end{corollary}

\begin{proof}
We first record the stronger Quillen input
\begin{equation}\label{d3:eq:explicitnonzeropopen}
 c_R|_{R[1/3]}\ne0\quad\text{in }K_3(R[1/3])_{\mathbf Q}.
\end{equation}
Proposition~\ref{d3:prop:directtcoriginal} detects the actual image of
$c_R$ on
$B=W\langle X,Y,Z\rangle/(XY+Z^2-1)$, where
$W=\mathbf Z_3[u]/(u^2+3)$, by $(x,y,z)=(uX,uY,uZ)$.
For $A^{\mathrm{sh}}$, first map to its completion and use the same
calculation over
$W_\infty=W(\overline{\mathbf F}_3)[u]/(u^2+3)$:
the nonzero endpoint difference remains nonzero in the larger
coefficient field, and Lemma~\ref{d3:lem:tcreviewresidue} applies there.
The perfect-residue-field scope of
\cite[Proposition~4.10]{AMMN} permits this base change.
In either case $C=B/u$ is an affine-line torsor over
$\mathbf P^1_k$, with $k=\mathbf F_3$ or
$\overline{\mathbf F}_3$. Hence $K_3(C)$ is torsion by the projective
bundle formula, Quillen's finite-field computation, and filtered
continuity. Localization and d\'evissage therefore make
$K_3(B)_{\mathbf Q}\to K_3(B[1/3])_{\mathbf Q}$ injective.
The nonzero detector image on $B[1/3]$ factors through $R[1/3]$,
proving~\eqref{d3:eq:explicitnonzeropopen}.

\emph{Milnor theory.}
There are natural symbol maps
\[
 K_3^{M,\mathrm{naive}}(R)
 \xrightarrow{\eta_R}\widehat K_3^M(R)
 \xrightarrow{\iota_R}K_3(R),
\]
whose composite sends a symbol to the product of its unit classes.
Kerz's universal construction, Proposition~4.2.8(4)--(6), and
Theorem~4.3.2 give these maps, surjectivity of $\eta_R$, and a map
$r_R:K_3(R)\to\widehat K_3^M(R)$ with
$r_R\iota_R=2$; these assertions hold for arbitrary local rings
\cite[Theorem~4.2.5, Proposition~4.2.8, Theorem~4.3.2]{KerzThesis}.
Thus $\iota_R$ is split injective rationally. Both displayed Milnor
classes have infinite order because their Quillen image is $c_R$.
For the first three rings the residue field is $\mathbf F_3$ and no
equality of naive and improved theory is asserted. For
$A^{\mathrm{sh}}$ the residue field is infinite, so $\eta_R$ is an
isomorphism integrally.

In the field $F_R$, Steinberg relations give
\[
 \{\lambda s_0,\mu\}
 =\{-x,\lambda s_0\}=\{-x,s_0\},
 \qquad 2\{a,s_0\}=0.
\]
The first equality uses $\{-x\mu,1+x\mu\}=0$, and the second
cancellation uses $\{-x,1+x\}=0$. Consequently
\[
 2\{a,\lambda,\mu\}
 =2\{a,\lambda s_0,\mu\}
 =-2\{a,s_0,-x\}=0.
\]
These calculations use $-x$ only in the field, where it is a unit.

\emph{Beilinson cohomology.}
All four rings are regular, noetherian, and finite dimensional.
Cisinski--D\'eglise identify
\[
 H_B^q(X,\mathbf Q(p))
 \simeq\operatorname{Gr}_\gamma^pK_{2p-q}(X)_{\mathbf Q}
\]
for every regular $X$, compatibly with the multiplicative Adams
decomposition \cite[Corollaries~14.2.14 and~14.2.19]{CisinskiDegliseMotives}.
A unit has Adams weight one, so $c_R$ has pure weight three and
corresponds to the asserted nonzero cup product. Its generic
restriction is zero by the preceding symbol calculation.

\emph{Ordinary cubical higher Chow groups.}
Let $\square=\mathbf P^1\setminus\{1\}$, with faces $0,\infty$.
Intersect the graph of
$(a,\lambda,\mu):\operatorname{Spec}R\to(\mathbf P^1)^3$ with
$\operatorname{Spec}R\times\square^3$, and call the resulting
closed cycle $\Gamma_R$. It has codimension three and
\[
 \Gamma_R\simeq\operatorname{Spec}R[1/(xy)],
 \qquad \gamma_R=2[\Gamma_R]\in CH^3(R,3).
\]
Indeed $a-1$ is a unit, $\lambda-1=x$, and
$\mu-1=\lambda y/4$; the three graph equations form a regular
sequence. This is not a morphism from all of $\operatorname{Spec}R$
to $\square^3$. Every face intersection is empty because the three
coordinates are units, so $\Gamma_R$ is admissible and has zero
boundary. Over $F_R$, Totaro's field isomorphism sends the Milnor
symbol to this graph; the integral vanishing already proved therefore
gives $\gamma_R|_{F_R}=0$
\cite[Theorem~1 and Sections~1--2]{TotaroMilnor}.

Set $U=\operatorname{Spec}R[1/3]$. This is a regular equal-characteristic
zero scheme, and the compatible comparisons give
\[
 CH^3(U,3)_{\mathbf Q}\simeq H_B^3(U,\mathbf Q(3))
       \simeq K_3(U)_{\mathbf Q}^{(3)}.
\]
Here $U$ is affine, excellent, regular, noetherian, finite dimensional,
and of equal characteristic zero. The comparison with the actual
Bloch cycle complex in precisely this affine scope is recorded in
\cite[Remark~1.3.2]{BondarkoIvanovChowWeights}; see also
\cite[Introduction~B, p.~xxii; Theorem~11.1.24 and
Examples~11.1.25 and~11.2.3]{CisinskiDegliseMotives}.
The cubical and simplicial cycle complexes give the same higher Chow
groups \cite[Theorem~4.7]{LevineHigherChowRevisited}.
The rational monoidal comparison with Beilinson motives is
\cite[Theorem~16.1.4 and Corollary~16.1.7]{CisinskiDegliseMotives}.
The graph-symbol normalization and its compatibility with the product
of unit classes can be checked on a smooth finite-type
characteristic-zero model and then pulled back: each map $A\to R$
is flat, and flat pullback takes the graph cycle to the displayed
graph cycle. Thus the restricted class $\gamma_R|_U$ corresponds to
$c_R|_U$.
Thus~\eqref{d3:eq:explicitnonzeropopen} implies
$\gamma_R|_U\ne0$, and hence $\gamma_R\ne0$.
Only open restriction on $R$ is used in this last step; no comparison
between the full mixed-characteristic cycle complex and Beilinson
cohomology on $R$ is assumed.
\end{proof}

